\documentclass[12pt,a4paper]{article}
%
%	STARWARE - Stile base per la scrittura di documenti interni in LaTex
%

%
%	Pacchetti globali
%	Fare qui eventuali aggiunte!
%
\usepackage[utf8]{inputenc}
\usepackage[italian]{babel}
\usepackage[babel]{csquotes}
\usepackage{url}
\usepackage{graphicx}
\usepackage[colorlinks]{hyperref}
\usepackage{lastpage}
\usepackage{fancyhdr}
\usepackage[top=1cm,bottom=4cm,left=80pt,right=80pt]{geometry} %disegna la linea
\usepackage{listings} %per grandi porzioni di codice
\usepackage{color}
\usepackage[table]{xcolor}
\usepackage{booktabs,tabularx}
\usepackage{makeidx}
\usepackage{fixltx2e}
\usepackage{hyperref}
\usepackage{enumitem}
\usepackage{color}
\usepackage[T1]{fontenc}
\usepackage{float}
\usepackage{svg}
\usepackage{amsmath}
\usepackage[toc]{glossaries}
\usepackage{dirtree}
\usepackage{listings}

\makeglossaries

\bibliographystyle{alpha}

%
%	VARIABILI GLOBALI
%
\newcommand{\nomeGruppo}{StarWare}
\newcommand{\mailGruppo}{starware.swe@gmail.com}
\newcommand{\uni}{Universit\`{a} degli Studi di Padova}
\newcommand{\uniAA}{2015/2016}
\newcommand{\Cardin}{Prof. Riccardo Cardin}
\newcommand{\Vardanega}{Prof. Tullio Vardanega}
\newcommand{\Zucchetti}{Zucchetti S.p.a.}
\newcommand{\prj}{Quizzipedia}
\newcommand{\prjL}{Quizzipedia: software per la gestione di questionari}

\newcommand{\AVI}{Alessio Vitella}
\newcommand{\AVE}{Andrea Venier}
\newcommand{\NDC}{Nicola De Cao}
\newcommand{\IB}{Igor Baylyak}
\newcommand{\WS}{Walter Sandon}
\newcommand{\TP}{Thomas Pigarelli}
\newcommand{\AB}{Anna Bonaldo}

\newcommand{\mgls}[1]{\gls{#1}\textsubscript{G}}
\newcommand{\mglspl}[1]{\glspl{#1}\textsubscript{G}}
\newcommand{\mGls}[1]{\Gls{#1}\textsubscript{G}}
\newcommand{\mGlspl}[1]{\Glspl{#1}\textsubscript{G}}

\newcommand{\AM}{\emph{\mGls{amministratore}}}
\newcommand{\AN}{\emph{\mGls{analista}}}
\newcommand{\PG}{\emph{\mGls{progettista}}}
\newcommand{\PR}{\emph{\mGls{programmatore}}}
\newcommand{\VR}{\emph{\mGls{verificatore}}}
\newcommand{\PM}{\emph{\mGls{project manager}}}

\newcommand{\AMpl}{\emph{\mGlspl{amministratore}}}
\newcommand{\ANpl}{\emph{\mGlspl{analista}}}
\newcommand{\PGpl}{\emph{\mGlspl{progettista}}}
\newcommand{\PRpl}{\emph{\mGlspl{programmatore}}}
\newcommand{\VRpl}{\emph{\mGlspl{verificatore}}}
\newcommand{\PMpl}{\emph{\mGlspl{project manager}}}

\newcommand{\RR}{\emph{\mGls{revisione dei requisiti}}}
\newcommand{\RA}{\emph{\mGls{revisione di accettazione}}}
\newcommand{\RP}{\emph{\mGls{revisione di progettazione}}}
\newcommand{\RQ}{\emph{\mGls{revisione di qualifica}}}

\newcommand{\NdP}{\emph{\mGls{norme di progetto}}}
\newcommand{\SdF}{\emph{\mGls{stidio di fattibilita}}}
\newcommand{\AdR}{\emph{\mGls{analisi dei requisiti}}}
\newcommand{\PdP}{\emph{\mGls{piano di progetto}}}
\newcommand{\PdQ}{\emph{\mGls{piano di qualifica}}}

\newcommand{\latex}[1]{\texttt{#1}}
\newcommand{\fileName}[1]{\texttt{#1}}
\newcommand{\filePath}[1]{\texttt{#1}}
\newcommand{\TODO}[1]{\texttt{\large \color{red} \underline{TODO: #1}}}

\newcommand{\licenza}{GNU GENERAL PUBLIC LICENSE V2}

%per compilare il template usare questi sotto e commentare glia altri:
%\newcommand{\logoLungo}{../imgs/logoLungo.png}
%\newcommand{\logoGrande}{../imgs/logoGrande.png}
\newcommand{\logoLungo}{../../../template/imgs/logoLungo.png}
\newcommand{\logoGrande}{../../../template/imgs/logoGrande.png}

%
%	Setup stili
%

\newcommand{\HRule}{\rule{\linewidth}{0.5mm}}

\definecolor{dkgreen}{rgb}{0,0.6,0}
\definecolor{gray}{rgb}{0.5,0.5,0.5}
\definecolor{mauve}{rgb}{0.58,0,0.82}
\definecolor{light}{RGB}{255,255,190}

%
%	Setup di pagina
%

%colorazione link
\hypersetup
{
	colorlinks=true,
	linkcolor=black,
	urlcolor=blue,
	citecolor=blue
}

%	Setup Header + Footer

\pagestyle{fancy}
\setlength{\headheight}{2cm} %settato grandezza header

\renewcommand{\footrulewidth}{0.5pt} %ridefinisco il valore della riga di intestazione
\renewcommand{\headrulewidth}{0.5pt} %ridefinisco il valore della riga di pie' di pagina
\addtolength{\headwidth}{\marginparsep}
\addtolength{\headwidth}{\marginparwidth}

\fancyhead{} %annulla head di default
\fancyfoot{} %annulla foot di default

%	Logo intestazione
\lhead{\includegraphics[scale=0.12]{\logoLungo}}


%	footer
\cfoot{
	\uni \ - \uniAA \\
	\href{mailto:\mailGruppo}{\mailGruppo}\\
	{\tiny Questo documento è distribuito sotto licenza {\licenza}}
}
\rfoot{
	\thepage\ di \pageref{LastPage}
}


% test subsubsubsection

\usepackage{titlesec}
\usepackage{hyperref}

\titleclass{\subsubsubsection}{straight}[\subsection]

\newcounter{subsubsubsection}[subsubsection]
\renewcommand\thesubsubsubsection{\thesubsubsection.\arabic{subsubsubsection}}
\renewcommand\theparagraph{\thesubsubsubsection.\arabic{paragraph}} % optional; useful if paragraphs are to be numbered

\titleformat{\subsubsubsection}
  {\normalfont\normalsize\bfseries}{\thesubsubsubsection}{1em}{}
\titlespacing*{\subsubsubsection}
{0pt}{3.25ex plus 1ex minus .2ex}{1.5ex plus .2ex}

\makeatletter
\renewcommand\paragraph{\@startsection{paragraph}{5}{\z@}%
  {3.25ex \@plus1ex \@minus.2ex}%
  {-1em}%
  {\normalfont\normalsize\bfseries}}
\renewcommand\subparagraph{\@startsection{subparagraph}{6}{\parindent}%
  {3.25ex \@plus1ex \@minus .2ex}%
  {-1em}%
  {\normalfont\normalsize\bfseries}}
\def\toclevel@subsubsubsection{4}
\def\toclevel@paragraph{5}
\def\toclevel@paragraph{6}
\def\l@subsubsubsection{\@dottedtocline{4}{7em}{4em}}
\def\l@paragraph{\@dottedtocline{5}{10em}{5em}}
\def\l@subparagraph{\@dottedtocline{6}{14em}{6em}}
\makeatother

\setcounter{secnumdepth}{4}
\setcounter{tocdepth}{4}


%pdflatex -synctex=1 -interaction=nonstopmode %.tex|makeglossaries %|pdflatex -synctex=1 -interaction=nonstopmode %.tex|pdflatex -synctex=1 -interaction=nonstopmode %.tex
\makeglossary

\newglossaryentry{responsabile} {
	name=responsabile,
	description={è il responsabile della gestione, pianificazione e realizzazione del progetto},
	plural=Responsabili
}

\newglossaryentry{verificatore} {
	name=verificatore,
	description={è il responsabile dell'attività di verifica},
	plural=Verificatori
}

\newglossaryentry{programmatore} {
	name=programmatore,
	description={è responsabile delle attività di codifica miranti alla realizzazione del prodotto e delle componenti di ausilio necessarie per l'esecuzione delle prove di verifica e validazione},
	plural=programmatori
}

\newglossaryentry{progettista} {
	name=progettista,
	description={è responsabile delle attività di progettazione},
	plural=Progettisti
}

\newglossaryentry{analista} {
	name=analista,
	description={è responsabile delle attività di analisi. },
	plural=Analisti
}

\newglossaryentry{amministratore} {
	name=amministratore,
	description={è responsabile dell'efficienza e dell'operatività dell'ambiente di sviluppo; si occupa della redazione e attuazione di piani e procedure di gestione della qualità; inoltre gestisce l'archivio della documentazione del progetto},
	plural=Amministratori
}

\newglossaryentry{revisione dei requisiti} {
	name=revisione dei requisiti,
	description={è una revisione formale che determina l'accesso del gruppo al progetto didattico e la concordanza con il cliente di una visione condivisa del prodotto atteso}
}

\newglossaryentry{revisione di accettazione} {
	name=revisione di accettazione,
	description={è una revisione formale per l'accertamento del soddisfacimento di tutti i requisiti e il completamento del progetto}
}

\newglossaryentry{revisione di progettazione} {
	name=revisione di progettazione,
	description={è una revisione di progresso che accerta la realizzabilità del prodotto e informa il cliente sulle caratteristiche del prodotto}
}

\newglossaryentry{revisione di qualifica} {
	name=revisione di qualifica,
	description={è una revisione di progresso che approva l'esito finale delle verifiche e attiva la fase di validazione}
}

\newglossaryentry{analisi} {
	name=analisi,
	description={è il periodo di preparazione e produzioni di documenti che precede la Revisione dei requisiti}
}

\newglossaryentry{progettazione} {
	name=progettazione,
	description={è il periodo che intercorre tra la Revisione dei requisiti e la Revisione di progettazione}
}

\newglossaryentry{codifica} {
	name=codifica,
	description={è il periodo che intercorre tra la Revisione di progettazione e la Revisione di qualifica}
}

\newglossaryentry{validazione} {
	name=validazione,
	description={è il periodo che intercorre tra la Revisione di qualifica e la Revisione di accettazione}
}

\newglossaryentry{repository} {
	name=repository,
	description={è dove i file sono memorizzati, spesso su un server}
}

\newglossaryentry{ruolo} {
	name=ruolo,
	description={una delle figure professionali che una persona fisica interpreta nel corso del progetto. I ruoli sono: responsabile, amministratore, analista, progettista, programmatore e verificatore},
    plural=ruoli
}

\newglossaryentry{svg} {
	name=SVG,
	description={è un formato per la visualizzazione di oggetti in grafica vettoriale. Per maggiori informazioni si veda \href{https://it.wikipedia.org/wiki/Scalable_Vector_Graphics}{qui}}
}

\newglossaryentry{png} {
	name=PNG,
	description={abbreviazione di Portable Network Graphics, è un formato di file per memorizzare immagini. Per ulteriori informazioni si veda \href{http://it.wikipedia.org/wiki/Portable_Network_Graphics}{qui}}
}

\newglossaryentry{pdf} {
	name=PDF,
	description={è un formato di file basato su un linguaggio di descrizione di pagina sviluppato da Adobe Systems nel 1993 per rappresentare documenti in modo indipendente dall’hardware e dal software utilizzati per generarli o per visualizzarli. Per ulteriori informazioni si veda \href{http://it.wikipedia.org/wiki/Portable_Document_Format}{qui}}
}

\newglossaryentry{uml} {
	name=UML,
	description={è un linguaggio di modellazione e specifica basato sul paradigma object-oriented. Per ulteriori informazioni si veda \href{http://it.wikipedia.org/wiki/Unified_Modeling_Language}{qui}}
}

\newglossaryentry{walkthrough} {
	name=walkthrough,
    description={consiste nella lettura di un documento o codice cercando errori ed anomalie senza un'idea precisa di quali tipi di errori sarà possibile trovare}
}

\newglossaryentry{lista di controllo} {
	name=lista di controllo,
	description={è un elenco di cose da fare per eseguire una determinata attività}
}

\newglossaryentry{inspection} {
	name=inspection,
	description={è la lettura mirata di un documento o codice cercando errori specifici}
}

\newglossaryentry{milestone} {
	name=milestone,
	description={momento saliente nello sviluppo di un prodotto software per la quale devono essere pronti documenti e/o funzionalità}
}

\newglossaryentry{ticket} {
	name=ticket,
	description={rappresenta un compito nell'organizzazione e distribuzione del lavoro all'interno del progetto},
	plural=tickets
}

\newglossaryentry{commit} {
	name=commit,
	description={è la copia di modifiche fatte su file locali verso la repository remota. Esso rappresenta anche un particolare stato della repository nel tempo}
}

\newglossaryentry{versionamento} {
	name=versionamento,
	description={è la gestione di un versioni multiple di un insieme di informazioni. Per maggiori informazioni si veda \href{http://it. wikipedia.org/wiki/Controllo_versione}{qui}}
}

\newglossaryentry{task} {
	name=task,
	description={è un compito secondo la definizione dello standard IEEE 12207},
	plural=tasks
}

\newglossaryentry{attivita} {
	name=attivita,
	description={è un insieme di task}
}

\newglossaryentry{redattore} {
	name=redattore,
	description={colui che redige un documento},
	plural=redattori
}

\newglossaryentry{proponente} {
	name=proponente,
	description={colui che ha proposto al committente un capitolato d'appalto}
}

\newglossaryentry{committente} {
	name=committente,
	description={colui che assegna un compito. In questo caso è il Professor Tullio Vardanega}
}

\newglossaryentry{quality assurance} {
	name=quality assurance,
	description={è l'insieme delle attività volte a garantire il soddisfacimento degli obiettivi della qualità}
}

\newglossaryentry{telegram} {
	name=telegram,
	description={è un servizio di messaggistica istantanea utilizzato dal gruppo per comunicazioni interne. Per maggiori informazioni si veda \href{https://it.wikipedia.org/wiki/Telegram_(software)}{qui}}
}

\newglossaryentry{browser} {
	name=browser,
	description={è un'applicazione per il recupero, la presentazione e la navigazione di risorse web}
}

\newglossaryentry{google drive} {
	name=Google Drive,
	description={è un servizio di memorizzazione e sincronizzazione online introdotto da Google il 24 aprile 2012. Per maggiori informazioni si veda \href{https://it.wikipedia.org/wiki/Google_Drive}{qui}}
}

\newglossaryentry{skype} {
	name=Skype,
	description={è un software proprietario freeware di messaggistica istantanea e VoIP. Per maggiori informazioni si veda \href{https://it.wikipedia.org/wiki/Skype}{qui}}
}

\newglossaryentry{gantt} {
	name=Gantt,
	description={è un diagramma di supporto alla gestione dei progetti}
}

\newglossaryentry{projectlibre} {
	name=ProjectLibre,
	description={è un software di gestione progettuale}
}

\newglossaryentry{pert} {
	name=PERT,
	description={è uno strumento volto alla programmazione delle attività che compongono il progetto e, più in generale, alla gestione degli aspetti temporali di quest'ultimo}
}

\newglossaryentry{subtask} {
	name=subtask,
	description={è un task compreso all'interno di un altro task. La totalità di tutti i subtasks costituisce un intero task}
	plural=subtasks
}

\newglossaryentry{ticketing} {
	name=ticketing,
	description={procedura con la quale il Responsabile assegna un task}
}

\newglossaryentry{git} {
	name=git,
	description={è un sistema software di controllo di versione distribuito}
}

\newglossaryentry{quizzpedia} {
	name=Quizzpedia,
	description={è il nome del prodotto software richiesto dal capitolato d'appalto scelto}
}

\newglossaryentry{schierabile} {
	name=schierabile,
	description={è la capacità di rilasciare al cliente, con relativa installazione e messa in funzione o esercizio, di una applicazione o di un sistema software tipicamente all'interno di un sistema informatico aziendale}
}

\newglossaryentry{cross-platform} {
	name=cross-platform,
	description={può essere riferito ad un linguaggio di programmazione, ad un'applicazione software o ad un dispositivo hardware che funziona su più di un sistema}
}

\newglossaryentry{qml} {
	name=QML,
	description={è un "Domain Specific Language" richiesto dal capitolato d'appalto per la definizione delle domande all'interno del sistema}
}

\newglossaryentry{tomcat} {
	name=Tomcat,
	description={è un application server nella forma di contenitore servlet open source sviluppato dalla Apache Software Foundation. Per maggiori informazioni si veda \href{https://it.wikipedia.org/wiki/Apache_Tomcat}{qui}}
}

\newglossaryentry{java} {
	name=Java,
	description={è un linguaggio di programmazione orientato agli oggetti, specificatamente progettato per essere il più possibile indipendente dalla piattaforma di esecuzione. Per maggiori informazioni si veda \href{https://it.wikipedia.org/wiki/Java_(linguaggio_di_programmazione)}{qui}}
}

\newglossaryentry{node.js} {
	name=Node.js,
	description={è un framework che permette di realizzare web application usando un linguaggio di programmazione che utilizza la stessa sintassi di JavaScrip. Utilizza un modello event-driven, anzichè il classico modello a processi o thread concorrenti, e ciò significa che si eseguono azioni solo al verificarsi di un evento. Questo modello asincrono rende leggero ed efficiente, ideale per applicazioni real-time in per dispositivi distribuiti}
}

\newglossaryentry{javascript} {
	name=JavaScript,
	description={linguaggio di scripting orientato agli oggetti comunemente usato nella programmazione
Web}
}

\newglossaryentry{postgresql} {
	name=PostgreSQL,
	description={e un sistema di gestione di basi di dati open source usato per applicazioni che richiedono caratteristiche molto complesse complesse}
}

\newglossaryentry{mongodb} {
	name=MongoDB,
	description={è un sistema gestionale di basi di dati NoSQL orientato ai documenti, adatto per ambienti che hanno la necessità d'immagazzinare grosse quantità di dati, e dove il linguaggio utilizzato per la gestione dei dati è \mgls{javascript}}
}

\newglossaryentry{html5} {
	name=HTML5,
	description={linguaggio di markup per la strutturazione delle pagine web}
}

\newglossaryentry{css3} {
	name=CSS3,
	description={è un linguaggio di programmazione web utilizzato per descrivere l'aspetto e la formattazione di un sito web al browser lato client. }
}

\newglossaryentry{xml} {
	name=XML,
	description={è un meta-linguaggio che fornisce un insieme standard di regole sintattiche per modellare la struttura di documenti e dati. Questo insieme di regole, definiscono le modalità secondo cui è possibile crearsi un proprio linguaggio di markup}
}

\newglossaryentry{nosql} {
	name=NoSQL,
	description={come dice anche il termine questi database non sono basati su SQL, non sono basati su uno schema relazionale. I database relazionali sono infatti ottimi quando esistono delle relazioni tra i dati che salviamo, ma sono poco performanti nel caso sia necessario salvare una grande quantità di dati, magari usando la scalabilità orizzontale, quando cioè si utilizzano più server dove salvare questi dati e non solamente incrementando la potenza di un singolo server.
}
}

\newglossaryentry{scala} {
	name=Scala,
	description={è un linguaggio di programmazione di tipo general-purpose multi-paradigma studiato per integrare le caratteristiche e funzionalità dei linguaggi orientati agli oggetti e dei linguaggi funzionali. Per maggiori informazioni si veda \href{https://it.wikipedia.org/wiki/Scala_(linguaggio_di_programmazione)}{qui}}
}

\newglossaryentry{akka} {
	name=Akka,
	description={è un toolkit di strumenti per la costruzione di applicazioni con elevata concorrenza di dati che necessitano di un sistema resilente per l'invio e la ricezione di messaggi}
}

\newglossaryentry{ble} {
	name=BLE,
	description={Bluetooth low energy,  pur mantenendo un range di comunicazione simile a quello classico, fornisce un  consumo energetico dei device notevolmente ridotto}
}

\newglossaryentry{mqtt} {
	name=MQTT,
	description={è un protocollo di messaggistica leggero posizionato in cima a TCP/IP, disegnato per le situazioni in cui è richiesto un basso impatto e dove la banda è limitata. }
}

\newglossaryentry{aws} {
	name=AWS,
	description={è un insieme di servizi di elaborazione rermoti, detti anche servizi web, che costituiscono una piattaforma di cloud computing offerto da Amazon. Per maggiori informazioni si veda \href{https://aws.amazon.com/it/}{qui}}
}

\newglossaryentry{heroku} {
	name=Heroku,
	description={è una  delle prime cloud Platform-as-a-Service (PaaS) che supportava solo Java come linguaggio di programmazione, oggi giorno ha aggiunto molti altri linguaggi come Scala, Phyton, etc. Per maggiori informazioni si veda \href{https://www.heroku.com}{qui}}
}

\newglossaryentry{github} {
	name=Github,
	description={software di controllo di versione che permette di aggiornare un file senza dover sovrascrivere le versioni precedenti}
}

\newglossaryentry{template} {
	name=template,
	description={traducibile in italiano come modello, indica o un programma o un documento idealizzato come un documento semicompilato cartaceo che ha degli spazi bianchi che saranno successivamente riempiti}
	plural=templates
}

\newglossaryentry{dropbox} {
	name=Dropbox,
	description={è un software di cloud storage multipiattaforma, che offre un servizio di file hosting e sincronizzazione automatica di file tramite web}
}

\newglossaryentry{open-source} {
	name=open-source,
	description={un software di cui gli autori ovvero i detentori dei diritti rendono pubblico il codice sorgente, favorendone il libero studio e permettendo a programmatori indipendenti di apportarvi modifiche ed estensioni. Questo è realizzato tramite apposite licenze d'uso}
}

\newglossaryentry{project management} {
	name=project management,
	description={si intende l'insieme di attività aziendali, svolte tipicamene da una figura dedicata e specializzata detta project manager, volte all'analisi, progettazione, pianificazione e realizzazione degli obiettivi di un progetto, gestendolo in tutte le sue caratteristiche e fasi evolutive, nel rispetto di precisi vincoli come i tempi, costi, risorse, scopi, qualità}
}

\newglossaryentry{linux} {
	name=Linux,
	description={è una famiglia di sistemi operativi open-source di tipo Unix-like, rilasciati sotto varie possibili distribuzioni, aventi la caratteristica comune di utilizzare come nucleo il kernel Linux}
}

\newglossaryentry{windows} {
	name=Windows,
	description={è una famiglia di ambienti operativi e sistemi operativi dedicati ai personal computer, alle workstation, ai server e agli smartphone. Il sistema operativo si chiama così per via della sua interfaccia di programmazione di un'applicazione a finestre}
}

\newglossaryentry{mac os} {
	name=Mac OS,
	description={il sistema operativo di Apple dedicato dedicati ai personal computer Macintosh, alle workstation, ai server e agli smartphone  }
}


\newglossaryentry{schedule variance} {
	name=schedule variance,
	description={ogni deviazione alle baseline di un progetto, misurata confrontando costo preventivato di programma di lavoro con il costo preventivato del lavoro svolto. Indica quindi al cliente se il progetto sta procedendo nei tempi stabiliti}
}

\newglossaryentry{cost variance} {
	name=cost variance,
	description={indica al management aziendale se il valore del costo realmente maturato è maggiore, uguale o minore rispetto al costo pianificato}
}


\newglossaryentry{merge} {
	name=merge,
	description={unire due o più quantità}
}

\newglossaryentry{slack} {
	name=slack,
	description={intervallo di tempo entro cui un evento deve avvenire nel rispetto dei vincoli logici e imposti dal reticolo di pianificazione senza compromettere la durata complessiva del progetto}
}

\newglossaryentry{baseline} {
	name=baseline,
	description={di progetto costituisce il punto di riferimento rispetto al quale calcolare gli scostamenti delle principali variabili implicate nella gestione di un progetto}
}

\newglossaryentry{asana} {
	name=Asana,
	description={è un software che permette a dei team di monitorare il loro lavoro e tener traccia dei risultati tramite l'assegnazione di task a componemti specifiche del gruppo di lavoro}
}

\newglossaryentry{deadline} {
	name=deadline,
	description={La deadline di un \mgls{task} è un indicazione dell'urgenza del \mgls{task}; rappresenta un un punto su una linea temporale ideale. Data di scadenza o termine entro il quale deve essere completato un compito assegnato.}
}

\newglossaryentry{revert} {
	name=revert,
	description={per annullare le ultime modifiche effettuate al repository remoto. Per maggiori informazioni si veda \href{https://git-scm.com/docs/}{qui}}
}

\newglossaryentry{backup} {
	name=backup,
	description={si indica la replicazione, su un qualunque supporto di memorizzazione, di materiale informativo archiviato nella memoria di massa, al fine di prevenire la perdita definitiva dei dati in caso di eventi malevoli accidentali o intenzionali}
}

\newglossaryentry{push} {
	name=push,
	description={per inviare modifiche di un documento site in un host locale al repository remoto. Per maggiori informazioni si veda \href{https://git-scm.com/docs/}{qui}}
}

\newglossaryentry{teamwork} {
	name=teamwork,
	description={è un software che permette a dei team di monitorare il loro lavoro e tener traccia dei risultati tramite l'assegnazione di task a componenti specifiche del gruppo di lavoro}
}

\newglossaryentry{evento} {
	name=evento,
	description={\mgls{teamwork}, attraverso il suo calendario, offre la possibilità di impostare eventi in giorni stabiliti dagli utenti. Questi eventi verranno periodicamente segnalati come promemoria a tutte le persone invitate agli stessi.}
}

\newglossaryentry{etichetta} {
	name=etichetta,
	description={è un controllo grafico che mostra informazioni testuali all'interno di un form}
	plural=etichette
}

\newglossaryentry{bug} {
	name=bug,
	description={identifica un errore nella scrittura di un programma software che ne causa un comportamento imprevisto o comunque diverso da quello specificato dal produttore}
	plural=bugs
}

\newglossaryentry{software} {
	name=software,
	description={e’ un termine generico che definisce programmi e procedure utilizzati per far eseguire al computer un determinato compito}
}

\newglossaryentry{desktop} {
	name=desktop,
	description={area dello schermo su cui appaiono le icone e le finestre rappresentanti le memorie di massa collegate al computer ed il loro contenuto}
}

\newglossaryentry{draw.io} {
	name=draw.io,
	description={é un software gratuito per la creazione di diagrammi di flusso, di processo, UML, e diagrammi di rete}
}

\newglossaryentry{tracy} {
	name=tracy,
	description={software opensource per il tracciamento}
}

\newglossaryentry{package} {
	name=package,
	description={è un meccanismo per organizzare classi Java in gruppi logici, principalmente allo scopo di definire namespace distinti per diversi contesti}
}

\newglossaryentry{stakeholder} {
	name=stakeholder,
	description={si indica genericamente un soggetto o un gruppo di soggetti influenti nei confronti di un'iniziativa economica, sia essa un'azienda o un progetto}
}

\newglossaryentry{branch} {
	name=branch,
	description={quando si vuole creare un nuovo ramo al repository remoto. Per maggiori informazioni si veda \href{https://git-scm.com/docs/}{qui}}
	plural=branches
}

\newglossaryentry{pull} {
	name=pull,
	description={un comando gitub per poter ricevere nell'host locale tutte le modifiche fatte nel repository remoto. Per maggiori informazioni si veda \href{https://git-scm.com/docs/}{qui}}
}

\newglossaryentry{underscore} {
	name=underscore,
	description={è un carattere che identifica il trattino basso}
}

\newglossaryentry{spelling} {
	name=spelling,
	description={è l'atto di pronunciare le parole lentamente, separando le singole lettere o le sillabe}
}

\newglossaryentry{cloud} {
	name=cloud,
	description={si indica un sistema di erogazione di risorse informatiche, come l'archiviazione, l'elaborazione o la trasmissione di dati, caratterizzato dalla disponibilità on demand attraverso Internet}
}

\newglossaryentry{smartphone} {
	name=smartphone,
	description={è un telefono cellulare con capacità di calcolo, di memoria e di connessione dati molto più avanzate rispetto ai normali telefoni cellulari}
	plural=smartphones
}

\newglossaryentry{checkbox} {
	name=checkbox,
	description={è un controllo grafico con cui l'utente può effettuare selezioni multiple}
}


\newglossaryentry{makefile} {
	name=makefile,
	description={è usata soprattutto per la compilazione di codice sorgente in codice oggetto, unendo e poi linkando il codice oggetto in programmi eseguibili o in librerie}
}

\newglossaryentry{gulpease} {
	name=Gulpease,
	description={è un indice di leggibilità di un testo tarato sulla lingua italiana. Rispetto ad altri ha il vantaggio di utilizzare la lunghezza delle parole in lettere anziché in sillabe, semplificandone il calcolo automatico}
}

\newglossaryentry{pdca} {
	name=PDCA,
	description={\TODO{}}
}


%Titolo documento
\newcommand{\titoloDocumento}{Analisi dei requisiti}

%Prima data di creazione del documento
\newcommand{\dataCreazione}{30 Novembre 2015}

%Inserite la versione attuale del documento
\newcommand{\versione}{1.0.6}

%Stato in cui si trova il documento: Formale solo all'atto di consegna
\newcommand{\stato}{Formale}

%Uso del documento
\newcommand{\uso}{Esterno}

\rhead{\titoloDocumento}
\lfoot{Versione: \versione}
\title{\titoloDocumento}

\begin{document}
\begin{titlepage}
\begin{center}
\intestazione{\titoloDocumento}

\begin{table}[h]
\begin{center}
\begin{tabular}{r | l}
\multicolumn{2}{c}{\textbf{Informazioni sul documento}}\\
\midrule
\textbf{Nome Documento} & \titoloDocumento \\
\textbf{Versione} & \versione \\
\textbf{Stato} & \emph{\stato} \\
\textbf{Uso} & \emph{\uso} \\
\textbf{Data Creazione} & \dataCreazione \\
\textbf{Data Ultima Modifica} & \today \\
\textbf{Redazione} & \AB\\
\  & \TP\\
\  &  \WS\\
\  &  \AVE\\
\textbf{Verifica} &  \IB\\
\ & \AVI \\
\textbf{Approvazione} &  \NDC\\
\textbf{Lista Distribuzione} & \nomeGruppo \\
\  & \Vardanega \\
\  & \Cardin \\
\  & Il proponente \Zucchetti \\

\end{tabular}
\end{center}
\end{table}

\end{center}
\end{titlepage}
\newpage

\Large{\textbf{Registro delle modifiche}}\\
\normalsize

\begin{center}
	\begin{longtable}[H]{p{0.12\textwidth} p{0.2\textwidth} p{0.18\textwidth} p{0.39\textwidth}}
		\toprule
		\textbf{Versione}	&	\textbf{Autore}	&	\textbf{Data}	&	\textbf{Descrizione}\\*
		\midrule
		\midrule
        1.0.6 & \TP{} \& \AVE{} & 2016-03-05 & Ridefiniti casi d'uso autenticazione e gestione profilo \\*
		\midrule
		1.0.5 & \TP{} & 2016-03-03 & Aggiunti attori Google, Facebook e Microsoft \\*
		\midrule
		1.0.4 & \AVI{} & 2016-02-24 &  Correzioni segnalazioni RR \\*
		\midrule
		1.0.3 & \IB{} & 2016-02-24 &  Modifica glossario \\*
		\midrule
		1.0.2 & \IB{} & 2016-02-24 &  Aggiunta riferimenti normativi e informativi \\*
		\midrule
		1.0.1 & \NDC{} & 2016-02-24 &  Corretti diagrammi use case \\*
		\midrule
		1.2.0 & \NDC{} & 2016-01-21 &  Approvazione \\*
		\midrule
		0.1.1 & \IB{} & 2016-01-21 &  Seconda verifica \\*
		\midrule
		0.1.0 & \AVI{} & 2016-01-20 &  Prima verifica \\*
		\midrule
		0.0.4 & \TP{} & 2016-01-19 & Aggiunta diagrammi casi d'uso\\*
		\midrule
		0.0.3 & \AVE{} & 2016-01-19 &  Introduzione documento e delle sezioni  \\*
		\midrule
		0.0.2 & \AB{} & 2016-01-19 &  Aggiunta tabella dei requisiti \\*
		\midrule
		0.0.1 & \WS{} & 2016-01-19 &  Aggiunta tabella dei casi d'uso \\*
		\midrule
		0.0.0 & \NDC{} & 2016-01-19 &  Creazione documento \\*
		\bottomrule
		\caption{Versionamento del documento}
		\label{tabVers1}
	\end{longtable}
\end{center}

\newpage
\tableofcontents
\newpage
\listoftables
\listoffigures
\newpage

\section{Introduzione}

\subsection{Scopo}
L’analisi del capitolato d’appalto e il successivo incontro con il proponente ha evidenziato un insieme di requisiti che il presente documento ha lo scopo di elencare e descrivere in modo dettagliato. Lo scopo di tale documento è quindi presentare le funzionalità che offrirà il prodotto.

\subsection{Descrizione}
Il documento presenterà una breve descrizione degli attori coinvolti nell'iterazione con il sistema proposto. Successivamente verranno elencati i casi d'uso individuati come essenziali per la creazione di un sistema che corrisponda alle esigenze del proponente.

\subsection{Descrizione Prodotto}
\descrizioneProdotto

\subsection{Glossario}\label{glossario}
\glossarioPrint

\subsection{Riferimenti}\label{riferimenti}
\subsubsection{Normativi}
\begin{itemize}
	\item Capitolato d’Appalto \prjL: \url{http://www.math.unipd.it/~tullio/IS-1/2015/Progetto/C5.pdf};
	\item Vincoli di organigramma e dettagli tecnico-economici: \url{http://www.math.unipd.it/~tullio/IS-1/2015/Progetto/PD01b.html};
	\item \NdP: \textit{NormeDiProgetto\_v2.0.0}. 
\end{itemize}

\newpage
\section{Attori}

\begin{figure}[H]
    \centering
    \begin{tikzpicture}
        \umlactor[x=0, y=0]{Utente}
        \umlactor[x=4, y=0]{Docente}
        \umlactor[x=4, y=4]{Amministratore}
        \umlactor[x=0, y=4]{Proprietario}
        \umlactor[x=0, y=-4]{Ospite}

        \umlinherit{Docente}{Utente}
        \umlinherit{Amministratore}{Docente}
        \umlinherit{Proprietario}{Amministratore}
    \end{tikzpicture}
    \caption{Attori}
    \label{fig:actors} 
\end{figure}

Di seguito viene fornita una breve descrizione degli attori utilizzati nei casi d'uso:

\begin{itemize}
    \item \textbf{Ospite:} attore che non ha ancora effettuato l'accesso ed è quindi limitato
        ad accedere o registrarsi
    \item \textbf{Utente:} attore che è registrato al servizio ed ha effettuato l'accesso al
        sistema
    \item \textbf{Docente:} utente in grado di gestire studenti, classi, argomenti, questionari
        e domande
    \item \textbf{Amministratore:} utente che gestisce i docenti
    \item \textbf{Proprietario:} amministratore iniziale che è in grado di creare nuovi amministratori
        e cancellarne di esistenti
\end{itemize}

Nei diagrammi seguenti, le generalizzazioni tra attori non vengono specificate perché
appesantiscono troppo l'immagine. Come riferimento è possibile utilizzare il diagramma mostrato in questa
sezione.

\newpage
\section{Casi d'uso}

\subsection{Struttura di un caso d'uso}
I casi d'uso sono creati secondo lo standard UML 2.4 e sono identificati dalla seguente notazione:
\begin{center}
	UC[codice]: [Titolo]
\end{center}
dove il \textbf{codice} è un numero progressivo identificativo di ogni requisito, gerarchico nel  caso di sotto-casi d'uso tramite la notazione \textit{CodiceUCPadre.CodiceSottoUC}. Per ogni caso d'uso deve inoltre essere indicato:
\begin{itemize}
	\item \textbf{Titolo:} breve ma non ambiguo
	\item \textbf{Attori:} principali e secondari coinvolti
	\item \textbf{Precondizione}
	\item \textbf{Flusso principale degli eventi:} descrizione, eventualmente per punti, del main scenario del Caso d'uso
	\item \textbf{Postcondizione}
\end{itemize}

% Output di tracy sotto a:
% tracy-tpiga.rhcloud.com/tracy/latex/UseCases?type=latex&raw=1
\hypertarget{UC1}{}
\subsection{Caso d'uso UC1: Autenticazione}
\begin{figure}[H]
	\centering
	\begin{resizedtikzpicture}{\textwidth}
		\umlactor[x=0, y=0]{Ospite}
		\begin{umlsystem}[x=0, fill=lightgray!20]{Quizzipedia}
			\umlusecase[x=5, y=-4, fill=white, width=4cm, name=216]{\textbf{UC1.5:} Recupero password}
			\umlassoc{Ospite}{216}
			\umlusecase[x=5, y=0, fill=white, width=4cm, name=192]{\textbf{UC1.2:} Inserimento password}
			\umlassoc{Ospite}{192}
			\umlusecase[x=5, y=4, fill=white, width=4cm, name=191]{\textbf{UC1.1:} Inserimento username}
			\umlassoc{Ospite}{191}
		\end{umlsystem}
	\end{resizedtikzpicture}
	\caption{\textbf{UC1}: Autenticazione}
	\label{UC1}
\end{figure}
\begin{description}
	\item[Attori:] Ospite;
	\item[Scopo e descrizione:] L'ospite inserisce le credenziali per accedere al sistema
	\item[Precondizione:] L'ospite non è autenticato e possiede un account all'interno del sistema;
	
	\item[Flusso principale degli eventi:] \ 
	\begin{enumerate}
		\item L'ospite inserisce l'username (\hyperlink{UC1.1}{UC1.1});
		\item L'ospite inserisce la password (\hyperlink{UC1.2}{UC1.2});
		\item L'utente richiede il recupero della password (\hyperlink{UC1.5}{UC1.5});
		
	\end{enumerate}
	\item[Postcondizione:] L’ospite è autenticato nel sistema.
\end{description}
\hypertarget{UC1.1}{}
\subsection{Caso d'uso UC1.1: Inserimento username}
\begin{figure}[H]
	\centering
	\begin{resizedtikzpicture}{\textwidth}
		\umlactor[x=0, y=0]{Ospite}
		\begin{umlsystem}[x=0, fill=lightgray!20]{Quizzipedia}
			\umlusecase[x=5, y=0, fill=white, width=4cm, name=191]{\textbf{UC1.1:} Inserimento username}
			\umlassoc{Ospite}{191}
			\umlusecase[x=15, y=0, fill=white, width=4cm, name=193]{\textbf{UC1.3:} Errore username non presente}
			\umlextend[name=ext-193-191]{193}{191}
			\umlnote[x=23, y=0, fill=white, width=4cm, geometry=|-|, arm=-2cm]{ext-193-191-1}{L'ospite ha inserito un username che non esiste nel sistema}
		\end{umlsystem}
	\end{resizedtikzpicture}
	\caption{\textbf{UC1.1}: Inserimento username}
	\label{UC1.1}
\end{figure}
\begin{description}
	\item[Attori:] Ospite;
	\item[Scopo e descrizione:] L'ospite inserisce l'username per l'accesso
	\item[Precondizione:] L'ospite non è autenticato e possiede un account all'interno del sistema;
	
	\item[Flusso principale degli eventi:] \ 
	\begin{enumerate}
		\item L'ospite inserisce l'username per l'accesso;
		
	\end{enumerate}
	\item[Estensioni:]
	\begin{enumerate}
		\item Se l'username inserito non è presente nel sistema viene visualizzato un messaggio di errore (\hyperlink{UC1.3}{UC1.3});
		
	\end{enumerate}
	\item[Postcondizione:] L'ospite ha inserito l'username.
\end{description}
\hypertarget{UC1.2}{}
\subsection{Caso d'uso UC1.2: Inserimento password}
\begin{figure}[H]
	\centering
	\begin{resizedtikzpicture}{\textwidth}
		\umlactor[x=0, y=0]{Ospite}
		\begin{umlsystem}[x=0, fill=lightgray!20]{Quizzipedia}
			\umlusecase[x=5, y=0, fill=white, width=4cm, name=192]{\textbf{UC1.2:} Inserimento password}
			\umlassoc{Ospite}{192}
			\umlusecase[x=15, y=0, fill=white, width=4cm, name=194]{\textbf{UC1.4:} Errore password errata}
			\umlextend[name=ext-194-192]{194}{192}
			\umlnote[x=23, y=0, fill=white, width=4cm, geometry=|-|, arm=-2cm]{ext-194-192-1}{L'ospite ha inserito la sua password per l'accesso}
		\end{umlsystem}
	\end{resizedtikzpicture}
	\caption{\textbf{UC1.2}: Inserimento password}
	\label{UC1.2}
\end{figure}
\begin{description}
	\item[Attori:] Ospite;
	\item[Scopo e descrizione:] L'ospite inserisce la sua password per l'accesso
	\item[Precondizione:] L'ospite non è autenticato e possiede un account all'interno del sistema
	;
	
	\item[Flusso principale degli eventi:] \ 
	\begin{enumerate}
		\item L'ospite inserisce la sua password per l'accesso
		;
		
	\end{enumerate}
	\item[Estensioni:]
	\begin{enumerate}
		\item Se la password inserita non è corretta viene visualizzato un messaggio di errore (\hyperlink{UC1.4}{UC1.4});
		
	\end{enumerate}
	\item[Postcondizione:] L'ospite ha inserito la sua password per l'accesso.
\end{description}
\hypertarget{UC1.3}{}
\subsection{Caso d'uso UC1.3: Errore username non presente}\begin{description}
	\item[Attori:] Ospite;
	\item[Scopo e descrizione:] L'username inserito non esiste nel sistema
	\item[Precondizione:] L'ospite ha inserito un username che non esiste nel sistema;
	
	\item[Flusso principale degli eventi:] \ 
	\begin{enumerate}
		\item L'username inserito non esiste nel sistema
		;
		
	\end{enumerate}
	\item[Postcondizione:] L'utente non viene autenticato e viene visualizzato un messaggio di errore.
\end{description}
\hypertarget{UC1.4}{}
\subsection{Caso d'uso UC1.4: Errore password errata}\begin{description}
	\item[Attori:] Ospite;
	\item[Scopo e descrizione:] La password per l'accesso inserita dall'ospite è incorretta
	\item[Precondizione:] L'ospite ha inserito la sua password per l'accesso;
	
	\item[Flusso principale degli eventi:] \ 
	\begin{enumerate}
		\item La password per l'accesso inserita dall'ospite non corrisponde a quella salvata nel sistema
		;
		
	\end{enumerate}
	\item[Postcondizione:] L'utente non viene autenticato e viene visualizzato un messaggio d'errore.
\end{description}
\hypertarget{UC1.5}{}
\subsection{Caso d'uso UC1.5: Recupero password}
\begin{figure}[H]
	\centering
	\begin{resizedtikzpicture}{\textwidth}
		\umlactor[x=0, y=0]{Ospite}
		\begin{umlsystem}[x=0, fill=lightgray!20]{Quizzipedia}
			\umlusecase[x=5, y=0, fill=white, width=4cm, name=216]{\textbf{UC1.5:} Recupero password}
			\umlassoc{Ospite}{216}
			\umlusecase[x=15, y=0, fill=white, width=4cm, name=217]{\textbf{UC1.6:} Errore mail per recupero password non presente}
			\umlextend[name=ext-217-216]{217}{216}
			\umlnote[x=23, y=0, fill=white, width=4cm, geometry=|-|, arm=-2cm]{ext-217-216-1}{L'ospite non ricorda la password
			}
		\end{umlsystem}
	\end{resizedtikzpicture}
	\caption{\textbf{UC1.5}: Recupero password}
	\label{UC1.5}
\end{figure}
\begin{description}
	\item[Attori:] Ospite;
	\item[Scopo e descrizione:] L'ospite non ricorda la password di accesso al sistema e inizia la procedura di recupero password
	\item[Precondizione:] L'ospite non ricorda la password
	;
	
	\item[Flusso principale degli eventi:] \ 
	\begin{enumerate}
		\item L'ospite inserisce la mail registrata nel sistema sulla quale riceve la nuova password;
		\item L'ospite conferma la procedura;
		
	\end{enumerate}
	\item[Estensioni:]
	\begin{enumerate}
		\item Se la mail inserita dall'ospite non è presente nel sistema, viene visualizzato un messaggio d'errore (\hyperlink{UC1.6}{UC1.6});
		
	\end{enumerate}
	\item[Postcondizione:] L'ospite ha ricevuto una mail con la nuova password
	.
\end{description}
\hypertarget{UC1.6}{}
\subsection{Caso d'uso UC1.6: Errore mail per recupero password non presente}\begin{description}
	\item[Attori:] Ospite;
	\item[Scopo e descrizione:] L'ospite inserisce una mail per il recupero password che non è presente nel sistema
	
	\item[Precondizione:] L'ospite non ricorda la password
	;
	
	\item[Flusso principale degli eventi:] \ 
	\begin{enumerate}
		\item L'ospite ha sbagliato ad inserire la sua mail nella procedura di recupero password e il sistema mostra un messaggio d'errore;
		
	\end{enumerate}
	\item[Postcondizione:] Non viene inviata la mail per il recupero della password e viene visualizzato un messaggio di errore
	.
\end{description}
\hypertarget{UC2}{}
\subsection{Caso d'uso UC2: Registrazione}
\begin{figure}[H]
	\centering
	\begin{resizedtikzpicture}{\textwidth}
		\umlactor[x=0, y=0]{Ospite}
		\begin{umlsystem}[x=0, fill=lightgray!20]{Quizzipedia}
			\umlusecase[x=5, y=-4, fill=white, width=4cm, name=186]{\textbf{UC2.3:} Inserimento password}
			\umlassoc{Ospite}{186}
			\umlusecase[x=5, y=0, fill=white, width=4cm, name=185]{\textbf{UC2.2:} Inserimento username}
			\umlassoc{Ospite}{185}
			\umlusecase[x=5, y=4, fill=white, width=4cm, name=184]{\textbf{UC2.1:} Inserimento nome completo}
			\umlassoc{Ospite}{184}
		\end{umlsystem}
	\end{resizedtikzpicture}
	\caption{\textbf{UC2}: Registrazione}
	\label{UC2}
\end{figure}
\begin{description}
	\item[Attori:] Ospite;
	\item[Scopo e descrizione:] L'utente inserisce i dati del profilo e le credenziali per l'accesso
	\item[Precondizione:] L'ospite non è autenticato nel sistema e non ha un account presso il sistema;
	
	\item[Flusso principale degli eventi:] \ 
	\begin{enumerate}
		\item L'utente inserisce il nome completo (\hyperlink{UC2.1}{UC2.1});
		\item L'utente inserisce l'username (\hyperlink{UC2.2}{UC2.2});
		\item L'utente inserisce la password (\hyperlink{UC2.3}{UC2.3});
		
	\end{enumerate}
	\item[Postcondizione:] L’ospite possiede un account di ruolo studente presso il sistema.
\end{description}
\hypertarget{UC2.1}{}
\subsection{Caso d'uso UC2.1: Inserimento nome completo}
\begin{figure}[H]
	\centering
	\begin{resizedtikzpicture}{\textwidth}
		\umlactor[x=0, y=0]{Ospite}
		\begin{umlsystem}[x=0, fill=lightgray!20]{Quizzipedia}
			\umlusecase[x=5, y=0, fill=white, width=4cm, name=184]{\textbf{UC2.1:} Inserimento nome completo}
			\umlassoc{Ospite}{184}
			\umlusecase[x=15, y=0, fill=white, width=4cm, name=187]{\textbf{UC2.4:} Errore nome completo troppo corto}
			\umlextend[name=ext-187-184]{187}{184}
			\umlnote[x=23, y=0, fill=white, width=4cm, geometry=|-|, arm=-2cm]{ext-187-184-1}{Il nome completo inserito dall'ospite è troppo corto}
		\end{umlsystem}
	\end{resizedtikzpicture}
	\caption{\textbf{UC2.1}: Inserimento nome completo}
	\label{UC2.1}
\end{figure}
\begin{description}
	\item[Attori:] Ospite;
	\item[Scopo e descrizione:] L'ospite inserisce il nome completo
	\item[Precondizione:] L'ospite non è autenticato nel sistema e non ha un account presso il sistema;
	
	\item[Flusso principale degli eventi:] \ 
	\begin{enumerate}
		\item L'ospite inserisce il nome completo di almeno 2 caratteri;
		
	\end{enumerate}
	\item[Estensioni:]
	\begin{enumerate}
		\item Se il nome completo è troppo corto viene visualizzato un errore (\hyperlink{UC2.4}{UC2.4});
		
	\end{enumerate}
	\item[Postcondizione:] L'ospite ha specificato il nome completo.
\end{description}
\hypertarget{UC2.2}{}
\subsection{Caso d'uso UC2.2: Inserimento username}
\begin{figure}[H]
	\centering
	\begin{resizedtikzpicture}{\textwidth}
		\umlactor[x=0, y=0]{Ospite}
		\begin{umlsystem}[x=0, fill=lightgray!20]{Quizzipedia}
			\umlusecase[x=5, y=0, fill=white, width=4cm, name=185]{\textbf{UC2.2:} Inserimento username}
			\umlassoc{Ospite}{185}
			\umlusecase[x=15, y=-2.5, fill=white, width=4cm, name=190]{\textbf{UC2.7:} Errore username già utilizzato}
			\umlextend[name=ext-190-185]{190}{185}
			\umlnote[x=23, y=-2.5, fill=white, width=4cm, geometry=|-|, arm=-2.25cm]{ext-190-185-1}{L'ospite ha inserito un username già utilizzato da qualche altro utente}
			\umlusecase[x=15, y=2.75, fill=white, width=4cm, name=188]{\textbf{UC2.5:} Errore username troppo corto}
			\umlextend[name=ext-188-185]{188}{185}
			\umlnote[x=23, y=2.75, fill=white, width=4cm, geometry=|-|, arm=-2cm]{ext-188-185-1}{L'username inserito dall'ospite è troppo corto
				
			}
		\end{umlsystem}
	\end{resizedtikzpicture}
	\caption{\textbf{UC2.2}: Inserimento username}
	\label{UC2.2}
\end{figure}
\begin{description}
	\item[Attori:] Ospite;
	\item[Scopo e descrizione:] L'ospite inserisce l'username con la quale effettuare l'autenticazione
	\item[Precondizione:] L'ospite non è autenticato nel sistema e non ha un account presso il sistema;
	
	\item[Flusso principale degli eventi:] \ 
	\begin{enumerate}
		\item L'ospite inserisce l'username, di almeno 6 caratteri, con la quale effettuare l'autenticazione  (\hyperlink{UC2.4}{UC2.4});
		
	\end{enumerate}
	\item[Estensioni:]
	\begin{enumerate}
		\item Se l'username è troppo corto viene visualizzato un errore (\hyperlink{UC2.5}{UC2.5});
		\item Se l'username è già stato utilizzato viene visualizzato un errore (\hyperlink{UC2.7}{UC2.7});
		
	\end{enumerate}
	\item[Postcondizione:] L'ospite ha specificato l'username.
\end{description}
\hypertarget{UC2.3}{}
\subsection{Caso d'uso UC2.3: Inserimento password}
\begin{figure}[H]
	\centering
	\begin{resizedtikzpicture}{\textwidth}
		\umlactor[x=0, y=0]{Ospite}
		\begin{umlsystem}[x=0, fill=lightgray!20]{Quizzipedia}
			\umlusecase[x=5, y=0, fill=white, width=4cm, name=186]{\textbf{UC2.3:} Inserimento password}
			\umlassoc{Ospite}{186}
			\umlusecase[x=15, y=0, fill=white, width=4cm, name=189]{\textbf{UC2.6:} Errore password troppo corta}
			\umlextend[name=ext-189-186]{189}{186}
			\umlnote[x=23, y=0, fill=white, width=4cm, geometry=|-|, arm=-2cm]{ext-189-186-1}{La password inserita dall'ospite è troppo corta
			}
		\end{umlsystem}
	\end{resizedtikzpicture}
	\caption{\textbf{UC2.3}: Inserimento password}
	\label{UC2.3}
\end{figure}
\begin{description}
	\item[Attori:] Ospite;
	\item[Scopo e descrizione:] L'ospite inserisce la password con la quale effettuare l'autenticazione
	\item[Precondizione:] L'ospite non è autenticato nel sistema e non ha un account presso il sistema;
	
	\item[Flusso principale degli eventi:] \ 
	\begin{enumerate}
		\item L'ospite inserisce la password di almeno 8 caratteri con la quale effettuare l'autenticazione;
		
	\end{enumerate}
	\item[Estensioni:]
	\begin{enumerate}
		\item Se la password è troppo corta viene visualizzato un errore	 (\hyperlink{UC2.6}{UC2.6});
		
	\end{enumerate}
	\item[Postcondizione:] L'ospite ha specificato la password.
\end{description}
\hypertarget{UC2.4}{}
\subsection{Caso d'uso UC2.4: Errore nome completo troppo corto}\begin{description}
	\item[Attori:] Ospite;
	\item[Scopo e descrizione:] Viene visualizzato un messaggio di errore nel caso il nome completo sia inferiore ai 2 caratteri
	\item[Precondizione:] Il nome completo inserito dall'ospite è troppo corto;
	
	\item[Flusso principale degli eventi:] \ 
	\begin{enumerate}
		\item Viene visualizzato un messaggio di errore nel caso il nome completo sia inferiore ai 2 caratteri;
		
	\end{enumerate}
	\item[Postcondizione:] Non viene registrato l'utente e viene visualizzato un errore.
\end{description}
\hypertarget{UC2.5}{}
\subsection{Caso d'uso UC2.5: Errore username troppo corto}\begin{description}
	\item[Attori:] Ospite;
	\item[Scopo e descrizione:] Viene visualizzato un messaggio d'errore nel caso in cui l'ospite inserisca uno username inferiore ai 6 caratteri
	\item[Precondizione:] L'username inserito dall'ospite è troppo corto
	
	;
	
	\item[Flusso principale degli eventi:] \ 
	\begin{enumerate}
		\item Viene visualizzato un messaggio d'errore nel caso in cui l'ospite inserisca uno username inferiore ai 6 caratteri;
		
	\end{enumerate}
	\item[Postcondizione:] Non viene registrato l'utente viene visualizzato un messaggio d'errore.
\end{description}
\hypertarget{UC2.6}{}
\subsection{Caso d'uso UC2.6: Errore password troppo corta}\begin{description}
	\item[Attori:] Ospite;
	\item[Scopo e descrizione:] Viene visualizzato un messaggio d'errore nel caso in cui l'ospite inserisca una password inferiore agli 8 caratteri
	\item[Precondizione:] La password inserita dall'ospite è troppo corta
	;
	
	\item[Flusso principale degli eventi:] \ 
	\begin{enumerate}
		\item Viene visualizzato un messaggio d'errore nel caso in cui l'ospite inserisca una password inferiore agli 8 caratteri;
		
	\end{enumerate}
	\item[Postcondizione:] Non viene registrato l'utente e viene visualizzato un messaggio d'errore
	.
\end{description}
\hypertarget{UC2.7}{}
\subsection{Caso d'uso UC2.7: Errore username già utilizzato}\begin{description}
	\item[Attori:] Ospite;
	\item[Scopo e descrizione:] Viene visualizzato un messaggio d'errore nel caso in cui l'ospite inserisca un username già utilizzato da qualche altro utente
	\item[Precondizione:] L'ospite ha inserito un username già utilizzato da qualche altro utente;
	
	\item[Flusso principale degli eventi:] \ 
	\begin{enumerate}
		\item Viene visualizzato un messaggio d'errore nel caso in cui l'ospite inserisca un username già utilizzato da qualche altro utente;
		
	\end{enumerate}
	\item[Postcondizione:] Non viene registrato l'utente e viene visualizzato un messaggio d'errore.
\end{description}
\hypertarget{UC3}{}
\subsection{Caso d'uso UC3: Esegui questionario}
\begin{figure}[H]
	\centering
	\begin{resizedtikzpicture}{\textwidth}
		\umlactor[x=0, y=0]{Ospite}
		\begin{umlsystem}[x=0, fill=lightgray!20]{Quizzipedia}
			\umlusecase[x=5, y=-6.25, fill=white, width=4cm, name=132]{\textbf{UC3.3:} Conferma questionario}
			\umlassoc{Ospite}{132}
			\umlusecase[x=5, y=-2, fill=white, width=4cm, name=205]{\textbf{UC3.11:} Visualizza valutazione questionario}
			\umlassoc{Ospite}{205}
			\umlusecase[x=5, y=2.25, fill=white, width=4cm, name=160]{\textbf{UC3.10:} Feedback questionario}
			\umlassoc{Ospite}{160}
			\umlusecase[x=5, y=6.25, fill=white, width=4cm, name=124]{\textbf{UC3.1:} Rispondi domanda}
			\umlassoc{Ospite}{124}
		\end{umlsystem}
	\end{resizedtikzpicture}
	\caption{\textbf{UC3}: Esegui questionario}
	\label{UC3}
\end{figure}
\begin{description}
	\item[Attori:] Ospite;
	\item[Scopo e descrizione:] L'ospite effettua un questionario
	\item[Precondizione:] L'ospite ha scelto e selezionato il questionario da eseguire;
	
	\item[Flusso principale degli eventi:] \ 
	\begin{enumerate}
		\item L'ospite risponde ad una domanda (\hyperlink{UC3.1}{UC3.1});
		\item L'ospite può andare alla domanda successiva;
		\item L'ospite può andare alla domanda precedente;
		\item L'ospite conferma il questionario (\hyperlink{UC3.3}{UC3.3});
		\item L'ospite visualizza la valutazione per il questionario (\hyperlink{UC3.11}{UC3.11});
		\item Lo studente può lasciare un feedback (\hyperlink{UC3.10}{UC3.10});
		
	\end{enumerate}
	\item[Postcondizione:] Il questionario è stato compilato e vengono visualizzati i risultati.
\end{description}
\hypertarget{UC3.1}{}
\subsection{Caso d'uso UC3.1: Rispondi domanda}
\begin{figure}[H]
	\centering
	\begin{resizedtikzpicture}{\textwidth}
		\umlactor[x=0, y=0]{Ospite}
		\begin{umlsystem}[x=0, fill=lightgray!20]{Quizzipedia}
			\umlusecase[x=5, y=0, fill=white, width=4cm, name=124]{\textbf{UC3.1:} Rispondi domanda}
			\umlassoc{Ospite}{124}
			\umlusecase[x=15, y=-11, fill=white, width=4cm, name=159]{\textbf{UC3.9:} Rispondi domanda a risposta aperta}
			\umlinherit{159}{124}
			\umlusecase[x=15, y=-6.5, fill=white, width=4cm, name=158]{\textbf{UC3.8:} Rispondi domanda con associazione di parole}
			\umlinherit{158}{124}
			\umlusecase[x=15, y=-2, fill=white, width=4cm, name=157]{\textbf{UC3.7:} Rispondi domanda di tipo testo con parole omesse}
			\umlinherit{157}{124}
			\umlusecase[x=15, y=2.5, fill=white, width=4cm, name=156]{\textbf{UC3.6:} Rispondi domanda a risposta multipla}
			\umlinherit{156}{124}
			\umlusecase[x=15, y=7, fill=white, width=4cm, name=155]{\textbf{UC3.5:} Rispondi domanda a scelta multipla}
			\umlinherit{155}{124}
			\umlusecase[x=15, y=11.25, fill=white, width=4cm, name=154]{\textbf{UC3.4:} Rispondi domanda vero/falso}
			\umlinherit{154}{124}
		\end{umlsystem}
	\end{resizedtikzpicture}
	\caption{\textbf{UC3.1}: Rispondi domanda}
	\label{UC3.1}
\end{figure}
\begin{description}
	\item[Attori:] Ospite;
	\item[Scopo e descrizione:] L'ospite riporta la risposta alla domanda corrente
	\item[Precondizione:] L'ospite ha iniziato l'esecuzione di un questionario;
	
	\item[Flusso principale degli eventi:] \ 
	\begin{enumerate}
		\item L'ospite seleziona la risposta che ritiene esatta con le modalità previste dal tipo di domanda;
		\item Lo studente può lasciare un feedback per dire che la domanda gli piace oppure per segnalare un errore (\hyperlink{UC3.10}{UC3.10});
		
	\end{enumerate}
	\item[Postcondizione:] L'ospite ha risposto alla domanda.
\end{description}
\hypertarget{UC3.2}{}
\subsection{Caso d'uso UC3.2: Errore domanda non risposta}\begin{description}
	\item[Attori:] Ospite;
	\item[Scopo e descrizione:] Il sistema non permette all'ospite di terminare un questionario in cui non tutte le domande hanno una risposta.
	\item[Precondizione:] L'ospite non ha riposto ad alcune domande nel momento in cui conferma il questionario;
	
	\item[Flusso principale degli eventi:] \ 
	\begin{enumerate}
		\item Viene visualizzato un messaggio di errore;
		
	\end{enumerate}
	\item[Postcondizione:] Il questionario non viene terminato.
\end{description}
\hypertarget{UC3.3}{}
\subsection{Caso d'uso UC3.3: Conferma questionario}
\begin{figure}[H]
	\centering
	\begin{resizedtikzpicture}{\textwidth}
		\umlactor[x=0, y=0]{Ospite}
		\begin{umlsystem}[x=0, fill=lightgray!20]{Quizzipedia}
			\umlusecase[x=5, y=0, fill=white, width=4cm, name=132]{\textbf{UC3.3:} Conferma questionario}
			\umlassoc{Ospite}{132}
			\umlusecase[x=15, y=0, fill=white, width=4cm, name=125]{\textbf{UC3.2:} Errore domanda non risposta}
			\umlextend[name=ext-125-132]{125}{132}
			\umlnote[x=23, y=0, fill=white, width=4cm, geometry=|-|, arm=-2.5cm]{ext-125-132-1}{L'ospite non ha riposto ad alcune domande nel momento in cui conferma il questionario}
		\end{umlsystem}
	\end{resizedtikzpicture}
	\caption{\textbf{UC3.3}: Conferma questionario}
	\label{UC3.3}
\end{figure}
\begin{description}
	\item[Attori:] Ospite;
	\item[Scopo e descrizione:] L'ospite conferma il questionario eseguito
	\item[Precondizione:] L'ospite è autenticato nel sistema;
	
	\item[Flusso principale degli eventi:] \ 
	\begin{enumerate}
		\item L'ospite conferma il questionario;
		
	\end{enumerate}
	\item[Estensioni:]
	\begin{enumerate}
		\item Se quando l'ospite conferma il questionario ci sono domande non risposte viene visualizzato un errore (\hyperlink{UC3.2}{UC3.2});
		
	\end{enumerate}
	\item[Postcondizione:] Il questionario viene confermato.
\end{description}
\hypertarget{UC3.4}{}
\subsection{Caso d'uso UC3.4: Rispondi domanda vero/falso}\begin{description}
	\item[Attori:] Ospite;
	\item[Scopo e descrizione:] L'ospite risponde ad una domanda di tipo vero/falso
	\item[Precondizione:] È stata iniziata l'esecuzione di un questionario;
	
	\item[Flusso principale degli eventi:] \ 
	\begin{enumerate}
		\item L'ospite seleziona vero o falso;
		
	\end{enumerate}
	\item[Postcondizione:] L'ospite ha risposto alla domanda di tipo vero/falso.
\end{description}
\hypertarget{UC3.5}{}
\subsection{Caso d'uso UC3.5: Rispondi domanda a scelta multipla}\begin{description}
	\item[Attori:] Ospite;
	\item[Scopo e descrizione:] L'ospite risponde ad una domanda a scelta multipla
	\item[Precondizione:] È stata iniziata l'esecuzione di un questionario;
	
	\item[Flusso principale degli eventi:] \ 
	\begin{enumerate}
		\item L'ospite seleziona la risposta che ritiene corretta;
		
	\end{enumerate}
	\item[Postcondizione:] L'ospite ha risposto alla domanda a scelta multipla.
\end{description}
\hypertarget{UC3.6}{}
\subsection{Caso d'uso UC3.6: Rispondi domanda a risposta multipla}\begin{description}
	\item[Attori:] Ospite;
	\item[Scopo e descrizione:] L'ospite risponde ad una domanda a risposta multipla
	\item[Precondizione:] È stata iniziata l'esecuzione di un questionario;
	
	\item[Flusso principale degli eventi:] \ 
	\begin{enumerate}
		\item L'ospite e seleziona la o le risposte che ritiene corrette;
		
	\end{enumerate}
	\item[Postcondizione:] L'ospite ha risposto alla domanda a risposta multipla.
\end{description}
\hypertarget{UC3.7}{}
\subsection{Caso d'uso UC3.7: Rispondi domanda di tipo testo con parole omesse}\begin{description}
	\item[Attori:] Ospite;
	\item[Scopo e descrizione:] L'ospite risponde ad una domanda di tipo testo con parole omesse
	\item[Precondizione:] È stata iniziata l'esecuzione di un questionario;
	
	\item[Flusso principale degli eventi:] \ 
	\begin{enumerate}
		\item L'ospite riempie il testo con parole omesse scegliendole da una lista;
		
	\end{enumerate}
	\item[Postcondizione:] L'ospite ha risposto alla domanda di tipo testo con parole omesse.
\end{description}
\hypertarget{UC3.8}{}
\subsection{Caso d'uso UC3.8: Rispondi domanda con associazione di parole}\begin{description}
	\item[Attori:] Ospite;
	\item[Scopo e descrizione:] L'ospite risponde ad una domanda con associazione di parole
	\item[Precondizione:] È stata iniziata l'esecuzione di un questionario;
	
	\item[Flusso principale degli eventi:] \ 
	\begin{enumerate}
		\item L'ospite e seleziona le associazioni di parole che ritiene corrette;
		
	\end{enumerate}
	\item[Postcondizione:] L'ospite ha risposto alla domanda con associazione di parole.
\end{description}
\hypertarget{UC3.9}{}
\subsection{Caso d'uso UC3.9: Rispondi domanda a risposta aperta}\begin{description}
	\item[Attori:] Ospite;
	\item[Scopo e descrizione:] L'ospite risponde ad una domanda a risposta aperta
	\item[Precondizione:] È stata iniziata l'esecuzione di un questionario;
	
	\item[Flusso principale degli eventi:] \ 
	\begin{enumerate}
		\item L'ospite inserisce la risposta come testo;
		
	\end{enumerate}
	\item[Postcondizione:] L'ospite ha risposto alla domanda a risposta aperta.
\end{description}
\hypertarget{UC3.10}{}
\subsection{Caso d'uso UC3.10: Feedback questionario}\begin{description}
	\item[Attori:] Ospite;
	\item[Scopo e descrizione:] L'ospite lascia un feedback positivo al questionario
	\item[Precondizione:] L'ospite ha completato il questionario e vede i suoi risultati;
	
	\item[Flusso principale degli eventi:] \ 
	\begin{enumerate}
		\item L'ospite può lasciare un feedback al questionario;
		
	\end{enumerate}
	\item[Postcondizione:] Il feedback dell'ospite viene memorizzato per il questionario svolto .
\end{description}
\hypertarget{UC3.11}{}
\subsection{Caso d'uso UC3.11: Visualizza valutazione questionario}\begin{description}
	\item[Attori:] Ospite;
	\item[Scopo e descrizione:] L'ospite visualizza la valutazione del questionario eseguito
	\item[Precondizione:] L'ospite ha eseguito e confermato il questionario;
	
	\item[Flusso principale degli eventi:] \ 
	\begin{enumerate}
		\item L'ospite visualizza la valutazione del questionario eseguito;
		
	\end{enumerate}
	\item[Postcondizione:] L'ospite visualizza la valutazione del questionario eseguito.
\end{description}
\hypertarget{UC4}{}
\subsection{Caso d'uso UC4: Ricerca questionario}
\begin{figure}[H]
	\centering
	\begin{resizedtikzpicture}{\textwidth}
		\umlactor[x=0, y=0]{Ospite}
		\begin{umlsystem}[x=0, fill=lightgray!20]{Quizzipedia}
			\umlusecase[x=5, y=0, fill=white, width=4cm, name=17]{\textbf{UC4:} Ricerca questionario}
			\umlassoc{Ospite}{17}
			\umlusecase[x=15, y=-8, fill=white, width=4cm, name=22]{\textbf{UC9:} Ricerca questionario per difficoltà}
			\umlinherit{22}{17}
			\umlusecase[x=15, y=-3.75, fill=white, width=4cm, name=21]{\textbf{UC8:} Ricerca questionario per docente}
			\umlinherit{21}{17}
			\umlusecase[x=15, y=0.25, fill=white, width=4cm, name=20]{\textbf{UC7:} Ricerca questionario per argomento}
			\umlinherit{20}{17}
			\umlusecase[x=15, y=4.25, fill=white, width=4cm, name=19]{\textbf{UC6:} Ricerca questionario per classe}
			\umlinherit{19}{17}
			\umlusecase[x=15, y=8.25, fill=white, width=4cm, name=18]{\textbf{UC5:} Ricerca questionario per titolo}
			\umlinherit{18}{17}
		\end{umlsystem}
	\end{resizedtikzpicture}
	\caption{\textbf{UC4}: Ricerca questionario}
	\label{UC4}
\end{figure}
\begin{description}
	\item[Attori:] Ospite;
	\item[Scopo e descrizione:] L'ospite ricerca un questionario
	\item[Precondizione:] L'ospite è nella sezione di ricerca dei questionari;
	
	\item[Flusso principale degli eventi:] \ 
	\begin{enumerate}
		\item L'ospite inserisce i dati per la ricerca;
		
	\end{enumerate}
	\item[Postcondizione:] Il sistema mostra la lista dei questionari che soddisfano la ricerca.
\end{description}
\hypertarget{UC5}{}
\subsection{Caso d'uso UC5: Ricerca questionario per titolo}\begin{description}
	\item[Attori:] Ospite;
	\item[Scopo e descrizione:] L'ospite ricerca un questionario per titolo
	\item[Precondizione:] L'ospite è nella sezione di ricerca dei questionari;
	
	\item[Flusso principale degli eventi:] \ 
	\begin{enumerate}
		\item L'ospite inserisce il titolo del questionario che vuole cercare;
		
	\end{enumerate}
	\item[Postcondizione:] Il sistema mostra la lista dei questionari il cui titolo corrisponde al titolo ricercato.
\end{description}
\hypertarget{UC6}{}
\subsection{Caso d'uso UC6: Ricerca questionario per classe}\begin{description}
	\item[Attori:] Ospite;
	\item[Scopo e descrizione:] Lo studente ricerca il questionario per classe
	\item[Precondizione:] Lo studente è autenticato presso il sistema;
	
	\item[Flusso principale degli eventi:] \ 
	\begin{enumerate}
		\item Lo studente seleziona la classe di cui vuole visualizzare i questionari;
		
	\end{enumerate}
	\item[Postcondizione:] Il sistema mostra la lista dei questionari la cui classe corrisponde alla classe ricercata.
\end{description}
\hypertarget{UC7}{}
\subsection{Caso d'uso UC7: Ricerca questionario per argomento}\begin{description}
	\item[Attori:] Ospite;
	\item[Scopo e descrizione:] L'ospite ricerca questionario per argomento
	\item[Precondizione:] L'ospite è nella sezione di ricerca dei questionari;
	
	\item[Flusso principale degli eventi:] \ 
	\begin{enumerate}
		\item L'ospite seleziona gli argomento di cui vuole visualizzare i questionari;
		
	\end{enumerate}
	\item[Postcondizione:] Il sistema mostra la lista dei questionari i cui argomenti corrispondono agli argomenti ricercati.
\end{description}
\hypertarget{UC8}{}
\subsection{Caso d'uso UC8: Ricerca questionario per docente}\begin{description}
	\item[Attori:] Ospite;
	\item[Scopo e descrizione:] L'ospite ricerca questionario per docente
	\item[Precondizione:] L'ospite è nella sezione di ricerca dei questionari;
	
	\item[Flusso principale degli eventi:] \ 
	\begin{enumerate}
		\item L'ospite inserisce il nome del docente di cui vuole cercare i questionari;
		
	\end{enumerate}
	\item[Postcondizione:] Il sistema mostra la lista dei questionari il cui docente corrisponde al docente ricercato.
\end{description}
\hypertarget{UC9}{}
\subsection{Caso d'uso UC9: Ricerca questionario per difficoltà}\begin{description}
	\item[Attori:] Ospite;
	\item[Scopo e descrizione:] Lo studente ricerca un questionario per difficoltà
	\item[Precondizione:] Lo studente è autenticato presso il sistema;
	
	\item[Flusso principale degli eventi:] \ 
	\begin{enumerate}
		\item Lo studente inserisce un limite inferiore e superiore di difficoltà di cui vuole cercare questionari;
		
	\end{enumerate}
	\item[Postcondizione:] Il sistema mostra la lista dei questionari la cui difficoltà corrisponde alla difficoltà ricercata.
\end{description}
\hypertarget{UC10}{}
\subsection{Caso d'uso UC10: Gestione profilo}
\begin{figure}[H]
	\centering
	\begin{resizedtikzpicture}{\textwidth}
		\umlactor[x=0, y=0]{Studente}
		\begin{umlsystem}[x=0, fill=lightgray!20]{Quizzipedia}
			\umlusecase[x=5, y=-4, fill=white, width=4cm, name=195]{\textbf{UC10.3:} Modifica password}
			\umlassoc{Studente}{195}
			\umlusecase[x=5, y=0, fill=white, width=4cm, name=175]{\textbf{UC10.2:} Modifica username}
			\umlassoc{Studente}{175}
			\umlusecase[x=5, y=4, fill=white, width=4cm, name=41]{\textbf{UC10.1:} Modifica nome completo}
			\umlassoc{Studente}{41}
		\end{umlsystem}
	\end{resizedtikzpicture}
	\caption{\textbf{UC10}: Gestione profilo}
	\label{UC10}
\end{figure}
\begin{description}
	\item[Attori:] Studente;
	\item[Scopo e descrizione:] Lo studente gestisce il proprio profilo potendo cambiare il suo nome, aggiungendo o rimuovendo le connessioni agli account dei provider 
	\item[Precondizione:] Lo studente è autenticato nel sistema;
	
	\item[Flusso principale degli eventi:] \ 
	\begin{enumerate}
		\item Lo studente visualizza le proprie informazioni personali;
		\item Lo studente può modificare il proprio nome completo (\hyperlink{UC10.1}{UC10.1});
		\item Lo studente può modificare il proprio username	 (\hyperlink{UC10.2}{UC10.2});
		\item Lo studente può modificare la propria password (\hyperlink{UC10.3}{UC10.3});
		
	\end{enumerate}
	\item[Postcondizione:] Il sistema ha apportato le modifiche al profilo dello studente.
\end{description}
\hypertarget{UC10.1}{}
\subsection{Caso d'uso UC10.1: Modifica nome completo}
\begin{figure}[H]
	\centering
	\begin{resizedtikzpicture}{\textwidth}
		\umlactor[x=0, y=0]{Studente}
		\begin{umlsystem}[x=0, fill=lightgray!20]{Quizzipedia}
			\umlusecase[x=5, y=0, fill=white, width=4cm, name=41]{\textbf{UC10.1:} Modifica nome completo}
			\umlassoc{Studente}{41}
			\umlusecase[x=15, y=0, fill=white, width=4cm, name=197]{\textbf{UC10.5:} Errore nome completo troppo corto}
			\umlextend[name=ext-197-41]{197}{41}
			\umlnote[x=23, y=0, fill=white, width=4cm, geometry=|-|, arm=-2cm]{ext-197-41-1}{Il nome completo inserito dallo studente è troppo corto
			}
		\end{umlsystem}
	\end{resizedtikzpicture}
	\caption{\textbf{UC10.1}: Modifica nome completo}
	\label{UC10.1}
\end{figure}
\begin{description}
	\item[Attori:] Studente;
	\item[Scopo e descrizione:] Lo studente modifica il proprio nome completo
	\item[Precondizione:] Lo studente è autenticato nel sistema;
	
	\item[Flusso principale degli eventi:] \ 
	\begin{enumerate}
		\item Lo studente modifica il proprio nome completo;
		
	\end{enumerate}
	\item[Estensioni:]
	\begin{enumerate}
		\item Se il nome completo è troppo corto viene visualizzato un errore (\hyperlink{UC10.5}{UC10.5});
		
	\end{enumerate}
	\item[Postcondizione:] Lo studente ha modificato il proprio nome completo.
\end{description}
\hypertarget{UC10.2}{}
\subsection{Caso d'uso UC10.2: Modifica username}
\begin{figure}[H]
	\centering
	\begin{resizedtikzpicture}{\textwidth}
		\umlactor[x=0, y=0]{Studente}
		\begin{umlsystem}[x=0, fill=lightgray!20]{Quizzipedia}
			\umlusecase[x=5, y=0, fill=white, width=4cm, name=175]{\textbf{UC10.2:} Modifica username}
			\umlassoc{Studente}{175}
			\umlusecase[x=15, y=-2.5, fill=white, width=4cm, name=198]{\textbf{UC10.6:} Errore username già utilizzato}
			\umlextend[name=ext-198-175]{198}{175}
			\umlnote[x=23, y=-2.5, fill=white, width=4cm, geometry=|-|, arm=-2.25cm]{ext-198-175-1}{Lo studente ha inserito un username già utilizzato da qualche altro utente}
			\umlusecase[x=15, y=2.75, fill=white, width=4cm, name=196]{\textbf{UC10.4:} Errore username troppo corto}
			\umlextend[name=ext-196-175]{196}{175}
			\umlnote[x=23, y=2.75, fill=white, width=4cm, geometry=|-|, arm=-2cm]{ext-196-175-1}{L'username inserito dallo studente è troppo corto}
		\end{umlsystem}
	\end{resizedtikzpicture}
	\caption{\textbf{UC10.2}: Modifica username}
	\label{UC10.2}
\end{figure}
\begin{description}
	\item[Attori:] Studente;
	\item[Scopo e descrizione:] Lo studente modifica il proprio username che deve contenere almeno 6 caratteri
	\item[Precondizione:] Lo studente è autenticato nel sistema;
	
	\item[Flusso principale degli eventi:] \ 
	\begin{enumerate}
		\item Lo studente modifica il proprio username;
		
	\end{enumerate}
	\item[Estensioni:]
	\begin{enumerate}
		\item Se l'username è troppo corto viene visualizzato un errore (\hyperlink{UC10.4}{UC10.4});
		\item Se l'username è già stato utilizzato viene visualizzato un errore (\hyperlink{UC10.6}{UC10.6});
		
	\end{enumerate}
	\item[Postcondizione:] Lo studente ha modificato il proprio username.
\end{description}
\hypertarget{UC10.3}{}
\subsection{Caso d'uso UC10.3: Modifica password}
\begin{figure}[H]
	\centering
	\begin{resizedtikzpicture}{\textwidth}
		\umlactor[x=0, y=0]{Studente}
		\begin{umlsystem}[x=0, fill=lightgray!20]{Quizzipedia}
			\umlusecase[x=5, y=-2, fill=white, width=4cm, name=212]{\textbf{UC10.3.3:} Inserimento nuova password}
			\umlassoc{Studente}{212}
			\umlusecase[x=5, y=2, fill=white, width=4cm, name=207]{\textbf{UC10.3.1:} Inserimento vecchia password}
			\umlassoc{Studente}{207}
		\end{umlsystem}
	\end{resizedtikzpicture}
	\caption{\textbf{UC10.3}: Modifica password}
	\label{UC10.3}
\end{figure}

\begin{figure}[H]
	\centering
	\begin{resizedtikzpicture}{\textwidth}
		\umlactor[x=0, y=0]{Studente}
		\begin{umlsystem}[x=0, fill=lightgray!20]{Quizzipedia}
			\umlusecase[x=5, y=0, fill=white, width=4cm, name=195]{\textbf{UC10.3:} Modifica password}
			\umlassoc{Studente}{195}
			\umlusecase[x=15, y=0, fill=white, width=4cm, name=199]{\textbf{UC10.7:} Errore password troppo corta}
			\umlextend[name=ext-199-195]{199}{195}
			\umlnote[x=23, y=0, fill=white, width=4cm, geometry=|-|, arm=-2cm]{ext-199-195-1}{La password inserita dallo studente è troppo corta
			}
		\end{umlsystem}
	\end{resizedtikzpicture}
	\caption{\textbf{UC10.3}: Modifica password}
	\label{UC10.3}
\end{figure}
\begin{description}
	\item[Attori:] Studente;
	\item[Scopo e descrizione:] Lo studente modifica la propria password che deve contenere almeno 8 caratteri
	\item[Precondizione:] Lo studente è autenticato nel sistema
	;
	
	\item[Flusso principale degli eventi:] \ 
	\begin{enumerate}
		\item Lo studente inserisce la vecchia password (\hyperlink{UC10.3.1}{UC10.3.1});
		\item Lo studente inserisce la nuova password per effettuare l'autenticazione nel sistema (\hyperlink{UC10.3.3}{UC10.3.3});
		\item Lo studente modifica la propria password che deve contenere almeno 8 caratteri;
		
	\end{enumerate}
	\item[Estensioni:]
	\begin{enumerate}
		\item Se la password è troppo corta viene visualizzato un errore	 (\hyperlink{UC10.7}{UC10.7});
		
	\end{enumerate}
	\item[Postcondizione:] Lo studente ha modificato il proprio nome completo.
\end{description}
\hypertarget{UC10.3.1}{}
\subsection{Caso d'uso UC10.3.1: Inserimento vecchia password}
\begin{figure}[H]
	\centering
	\begin{resizedtikzpicture}{\textwidth}
		\umlactor[x=0, y=0]{Studente}
		\begin{umlsystem}[x=0, fill=lightgray!20]{Quizzipedia}
			\umlusecase[x=5, y=0, fill=white, width=4cm, name=207]{\textbf{UC10.3.1:} Inserimento vecchia password}
			\umlassoc{Studente}{207}
			\umlusecase[x=15, y=0, fill=white, width=4cm, name=208]{\textbf{UC10.3.2:} Errore password non corrispondente}
			\umlextend[name=ext-208-207]{208}{207}
			\umlnote[x=23, y=0, fill=white, width=4cm, geometry=|-|, arm=-2cm]{ext-208-207-1}{Lo studente ha inserito la sua vecchia password ma è errata}
		\end{umlsystem}
	\end{resizedtikzpicture}
	\caption{\textbf{UC10.3.1}: Inserimento vecchia password}
	\label{UC10.3.1}
\end{figure}
\begin{description}
	\item[Attori:] Studente;
	\item[Scopo e descrizione:] Lo studente inserisce la sua vecchia password
	\item[Precondizione:] Lo studente è autenticato nel sistema;
	
	\item[Flusso principale degli eventi:] \ 
	\begin{enumerate}
		\item Lo studente inserisce la sua vecchia password corretta;
		
	\end{enumerate}
	\item[Estensioni:]
	\begin{enumerate}
		\item Se alla richiesta di inserimento della vecchia password essa non corrisponde con quella presente nel sistema viene visualizzato un messaggio d'errore (\hyperlink{UC10.3.2}{UC10.3.2});
		
	\end{enumerate}
	\item[Postcondizione:] Lo studente ha inserito la sua vecchia password corretta.
\end{description}
\hypertarget{UC10.3.2}{}
\subsection{Caso d'uso UC10.3.2: Errore password non corrispondente}\begin{description}
	\item[Attori:] Studente;
	\item[Scopo e descrizione:] La vecchia password inserita dallo studente non è corretta
	\item[Precondizione:] Lo studente ha inserito la sua vecchia password ma è errata;
	
	\item[Flusso principale degli eventi:] \ 
	\begin{enumerate}
		\item La vecchia password inserita dallo studente non è corretta
		;
		
	\end{enumerate}
	\item[Postcondizione:] La password non viene modificata e viene visualizzato un messaggio d'errore.
\end{description}
\hypertarget{UC10.3.3}{}
\subsection{Caso d'uso UC10.3.3: Inserimento nuova password}\begin{description}
	\item[Attori:] Studente;
	\item[Scopo e descrizione:] Lo studente ha inserito la nuova password
	\item[Precondizione:] Lo studente è autenticato nel sistema;
	
	\item[Flusso principale degli eventi:] \ 
	\begin{enumerate}
		\item Lo studente inserisce la password di almeno 8 caratteri con la quale effettuare l'autenticazione	;
		
	\end{enumerate}
	\item[Postcondizione:] Lo studente ha inserito la sua nuova password.
\end{description}
\hypertarget{UC10.4}{}
\subsection{Caso d'uso UC10.4: Errore username troppo corto}\begin{description}
	\item[Attori:] Studente;
	\item[Scopo e descrizione:] Viene visualizzato un messaggio d'errore nel caso in cui lo studente inserisca uno username inferiore ai 6 caratteri
	\item[Precondizione:] L'username inserito dallo studente è troppo corto;
	
	\item[Flusso principale degli eventi:] \ 
	\begin{enumerate}
		\item Viene visualizzato un messaggio d'errore nel caso in cui lo studente inserisca un username inferiore ai 6 caratteri;
		
	\end{enumerate}
	\item[Postcondizione:] Non viene modificato l'username dell'utente e viene visualizzato un messaggio d'errore.
\end{description}
\hypertarget{UC10.5}{}
\subsection{Caso d'uso UC10.5: Errore nome completo troppo corto}\begin{description}
	\item[Attori:] Studente;
	\item[Scopo e descrizione:] Viene visualizzato un messaggio di errore nel caso il nome completo sia inferiore ai 2 caratteri
	
	\item[Precondizione:] Il nome completo inserito dallo studente è troppo corto
	;
	
	\item[Flusso principale degli eventi:] \ 
	\begin{enumerate}
		\item Viene visualizzato un messaggio di errore nel caso il nome completo sia inferiore ai 2 caratteri	;
		
	\end{enumerate}
	\item[Postcondizione:] Non viene modificato il nome e viene visualizzato un errore
	.
\end{description}
\hypertarget{UC10.6}{}
\subsection{Caso d'uso UC10.6: Errore username già utilizzato}\begin{description}
	\item[Attori:] Studente;
	\item[Scopo e descrizione:] Viene visualizzato un messaggio d'errore nel caso in cui lo studente inserisca un username già utilizzato da qualche altro utente
	\item[Precondizione:] Lo studente ha inserito un username già utilizzato da qualche altro utente;
	
	\item[Flusso principale degli eventi:] \ 
	\begin{enumerate}
		\item Viene visualizzato un messaggio d'errore nel caso in cui lo studente inserisca un username già utilizzato da qualche altro utente;
		
	\end{enumerate}
	\item[Postcondizione:] Non viene modificato l'username e viene visualizzato un messaggio d'errore.
\end{description}
\hypertarget{UC10.7}{}
\subsection{Caso d'uso UC10.7: Errore password troppo corta}\begin{description}
	\item[Attori:] Studente;
	\item[Scopo e descrizione:] Viene visualizzato un messaggio d'errore nel caso in cui lo studente inserisca una password inferiore agli 8 caratteri
	
	\item[Precondizione:] La password inserita dallo studente è troppo corta
	;
	
	\item[Flusso principale degli eventi:] \ 
	\begin{enumerate}
		\item Viene visualizzato un messaggio d'errore nel caso in cui lo studente inserisca una password inferiore agli 8 caratteri	;
		
	\end{enumerate}
	\item[Postcondizione:] Non viene modificata la password e viene visualizzato un messaggio d'errore
	.
\end{description}
\hypertarget{UC11}{}
\subsection{Caso d'uso UC11: Logout}\begin{description}
	\item[Attori:] Studente;
	\item[Scopo e descrizione:] Lo studente effettua il logout
	\item[Precondizione:] Lo studente è autenticato nel sistema;
	
	\item[Flusso principale degli eventi:] \ 
	\begin{enumerate}
		\item Lo studente seleziona la funzionalità di Logout;
		
	\end{enumerate}
	\item[Postcondizione:] Lo studente non è più autenticato nel sistema.
\end{description}
\hypertarget{UC12}{}
\subsection{Caso d'uso UC12: Iscrizione ad una classe}
\begin{figure}[H]
	\centering
	\begin{resizedtikzpicture}{\textwidth}
		\umlactor[x=0, y=0]{Studente}
		\begin{umlsystem}[x=0, fill=lightgray!20]{Quizzipedia}
			\umlusecase[x=5, y=0, fill=white, width=4cm, name=97]{\textbf{UC12.1:} Inserisci password classe}
			\umlassoc{Studente}{97}
		\end{umlsystem}
	\end{resizedtikzpicture}
	\caption{\textbf{UC12}: Iscrizione ad una classe}
	\label{UC12}
\end{figure}
\begin{description}
	\item[Attori:] Studente;
	\item[Scopo e descrizione:] Lo studente si iscrive ad una classe
	\item[Precondizione:] Lo studente è autenticato presso il sistema;
	
	\item[Flusso principale degli eventi:] \ 
	\begin{enumerate}
		\item Lo studente cerca una classe (\hyperlink{UC14}{UC14});
		\item Lo studente inserisce la password (che gli deve essere già stata consegnata personalmente dal docente della classe) per iscriversi alla classe (\hyperlink{UC12.1}{UC12.1});
		\item Lo studente conferma l'iscrizione alla classe;
		
	\end{enumerate}
	\item[Postcondizione:] Lo studente si è iscritto alla classe desiderata.
\end{description}
\hypertarget{UC12.1}{}
\subsection{Caso d'uso UC12.1: Inserisci password classe}
\begin{figure}[H]
	\centering
	\begin{resizedtikzpicture}{\textwidth}
		\umlactor[x=0, y=0]{Studente}
		\begin{umlsystem}[x=0, fill=lightgray!20]{Quizzipedia}
			\umlusecase[x=5, y=0, fill=white, width=4cm, name=97]{\textbf{UC12.1:} Inserisci password classe}
			\umlassoc{Studente}{97}
			\umlusecase[x=15, y=0, fill=white, width=4cm, name=98]{\textbf{UC12.2:} Errore password classe}
			\umlextend[name=ext-98-97]{98}{97}
			\umlnote[x=23, y=0, fill=white, width=4cm, geometry=|-|, arm=-2.25cm]{ext-98-97-1}{Lo studente ha inserito una password errata per iscriversi alla classe}
		\end{umlsystem}
	\end{resizedtikzpicture}
	\caption{\textbf{UC12.1}: Inserisci password classe}
	\label{UC12.1}
\end{figure}
\begin{description}
	\item[Attori:] Studente;
	\item[Scopo e descrizione:] Lo studente inserisce la password per l'iscrizione alla classe
	\item[Precondizione:] Lo studente è autenticato presso il sistema;
	
	\item[Flusso principale degli eventi:] \ 
	\begin{enumerate}
		\item Lo studente inserisce la password corretta (che gli deve essere già stata consegnata personalmente dal docente della classe) per iscriversi alla classe;
		
	\end{enumerate}
	\item[Estensioni:]
	\begin{enumerate}
		\item Se la password inserita non è corretta viene visualizzato un messaggio di errore (\hyperlink{UC12.2}{UC12.2});
		
	\end{enumerate}
	\item[Postcondizione:] Lo studente ha inserito la password per iscriversi alla classe.
\end{description}
\hypertarget{UC12.2}{}
\subsection{Caso d'uso UC12.2: Errore password classe}\begin{description}
	\item[Attori:] Studente;
	\item[Scopo e descrizione:] Il sistema avverte lo studente che ha inserito una password errata per l'iscrizione alla classe
	\item[Precondizione:] Lo studente ha inserito una password errata per iscriversi alla classe;
	
	\item[Flusso principale degli eventi:] \ 
	\begin{enumerate}
		\item Viene visualizzato un messaggio di errore;
		
	\end{enumerate}
	\item[Postcondizione:] Lo studente non viene iscritto alla classe.
\end{description}
\hypertarget{UC13}{}
\subsection{Caso d'uso UC13: Visualizza storico studente}
\begin{figure}[H]
	\centering
	\begin{resizedtikzpicture}{\textwidth}
		\umlactor[x=0, y=0]{Studente}
		\begin{umlsystem}[x=0, fill=lightgray!20]{Quizzipedia}
			\umlusecase[x=5, y=-4.5, fill=white, width=4cm, name=139]{\textbf{UC13.3:} Visualizza sommario statistiche studente}
			\umlassoc{Studente}{139}
			\umlusecase[x=5, y=0, fill=white, width=4cm, name=138]{\textbf{UC13.2:} Visualizza statistiche questionari studente}
			\umlassoc{Studente}{138}
			\umlusecase[x=5, y=4.5, fill=white, width=4cm, name=136]{\textbf{UC13.1:} Visualizza statistiche domande studente}
			\umlassoc{Studente}{136}
		\end{umlsystem}
	\end{resizedtikzpicture}
	\caption{\textbf{UC13}: Visualizza storico studente}
	\label{UC13}
\end{figure}
\begin{description}
	\item[Attori:] Studente;
	\item[Scopo e descrizione:] Lo studente visualizza il proprio storico
	\item[Precondizione:] Lo studente è autenticato presso il sistema;
	
	\item[Flusso principale degli eventi:] \ 
	\begin{enumerate}
		\item Lo studente può visualizzare le sue risposte e le statistiche generali relative alle domande eseguite (\hyperlink{UC13.1}{UC13.1});
		\item Lo studente può visualizzare le sue risposte e le statistiche generali relative ai questionari eseguiti (\hyperlink{UC13.2}{UC13.2});
		\item Lo studente può visualizzare un sommario delle proprie statistiche (\hyperlink{UC13.3}{UC13.3});
		
	\end{enumerate}
	\item[Postcondizione:] Lo studente ha visualizzato le sue risposte e le statistiche generali relative ai test eseguiti.
\end{description}
\hypertarget{UC13.1}{}
\subsection{Caso d'uso UC13.1: Visualizza statistiche domande studente}\begin{description}
	\item[Attori:] Studente;
	\item[Scopo e descrizione:] Lo studente visualizza la percentuale delle proprie risposte corrette e le statistiche generali della domanda
	\item[Precondizione:] Lo studente è autenticato presso il sistema;
	
	\item[Flusso principale degli eventi:] \ 
	\begin{enumerate}
		\item Lo studente cerca una domanda tra quelle che ha eseguito in passato (\hyperlink{UC25}{UC25});
		\item Lo studente seleziona la domanda per visualizzarne le risposte date e le statistiche generali;
		\item Lo studente visualizza le risposte già date in passato alla domanda selezionata;
		\item Lo studente visualizza la difficoltà della domanda calcolata come funzione delle risposte date da tutti gli utenti;
		
	\end{enumerate}
	\item[Postcondizione:] Lo studente ha visualizzato i risultati e le statistiche relative alle domande desiderate.
\end{description}
\hypertarget{UC13.2}{}
\subsection{Caso d'uso UC13.2: Visualizza statistiche questionari studente}\begin{description}
	\item[Attori:] Studente;
	\item[Scopo e descrizione:] Lo studente visualizza i punteggi che ha ottenuto nel questionario e le statistiche generali dello stesso
	\item[Precondizione:] Lo studente è autenticato presso il sistema;
	
	\item[Flusso principale degli eventi:] \ 
	\begin{enumerate}
		\item Lo studente cerca un questionario tra quelli che ha eseguito in passato (\hyperlink{UC4}{UC4});
		\item Lo studente seleziona il questionario per visualizzarne le risposte date e le statistiche generali;
		\item Lo studente può vedere le statistiche di ogni singola domanda del questionario selezionato (\hyperlink{UC13.1}{UC13.1});
		\item Lo studente vede le proprie risposte al questionario e il punteggio totale;
		\item Lo studente vede la difficoltà del questionario calcolata come funzione delle risposte date da tutti gli utenti;
		
	\end{enumerate}
	\item[Postcondizione:] Lo studente ha visualizzato i risultati e le statistiche relative ai questionari desiderati.
\end{description}
\hypertarget{UC13.3}{}
\subsection{Caso d'uso UC13.3: Visualizza sommario statistiche studente}\begin{description}
	\item[Attori:] Studente;
	\item[Scopo e descrizione:] Lo studente visualizza il totale delle risposte corrette sul totale delle risposte date e la media dei punteggi su tutti i questionari
	\item[Precondizione:] Lo studente è autenticato presso il sistema;
	
	\item[Flusso principale degli eventi:] \ 
	\begin{enumerate}
		\item Lo studente visualizza il totale delle domande eseguite e il numero di risposte corrette;
		\item Lo studente visualizza il totale dei questionari eseguiti e la media dei risultati ottenuti;
		\item Lo studente visualizza la media delle difficoltà delle domande a cui ha risposto correttamente;
		\item Lo studente visualizza la media delle difficoltà delle domande a cui non ha risposto correttamente;
		\item Per ogni argomento lo studente visualizza il totale delle domande eseguite, il numero di risposte corrette, la media delle difficoltà delle risposte corrette e la media delle difficoltà delle risposte non corrette;
		
	\end{enumerate}
	\item[Postcondizione:] Lo studente ha visualizzato i risultati e le statistiche relative alle domande desiderate.
\end{description}
\hypertarget{UC14}{}
\subsection{Caso d'uso UC14: Ricerca classe}
\begin{figure}[H]
	\centering
	\begin{resizedtikzpicture}{\textwidth}
		\umlactor[x=0, y=0]{Studente}
		\begin{umlsystem}[x=0, fill=lightgray!20]{Quizzipedia}
			\umlusecase[x=5, y=0, fill=white, width=4cm, name=28]{\textbf{UC14:} Ricerca classe}
			\umlassoc{Studente}{28}
			\umlusecase[x=15, y=-2, fill=white, width=4cm, name=30]{\textbf{UC16:} Ricerca classe per argomento}
			\umlinherit{30}{28}
			\umlusecase[x=15, y=2, fill=white, width=4cm, name=29]{\textbf{UC15:} Ricerca classe per docente}
			\umlinherit{29}{28}
		\end{umlsystem}
	\end{resizedtikzpicture}
	\caption{\textbf{UC14}: Ricerca classe}
	\label{UC14}
\end{figure}
\begin{description}
	\item[Attori:] Studente;
	\item[Scopo e descrizione:] Lo studente ricerca una classe
	\item[Precondizione:] Lo studente è autenticato presso il sistema;
	
	\item[Flusso principale degli eventi:] \ 
	\begin{enumerate}
		\item Lo studente inserisce i dati per la ricerca;
		
	\end{enumerate}
	\item[Postcondizione:] Il sistema mostra la lista delle classi che soddisfano la ricerca.
\end{description}
\hypertarget{UC15}{}
\subsection{Caso d'uso UC15: Ricerca classe per docente}\begin{description}
	\item[Attori:] Studente;
	\item[Scopo e descrizione:] Lo studente ricerca una classe per docente
	\item[Precondizione:] Lo studente è autenticato presso il sistema;
	
	\item[Flusso principale degli eventi:] \ 
	\begin{enumerate}
		\item Lo studente inserisce il nome del docente di cui le classi;
		
	\end{enumerate}
	\item[Postcondizione:] Il sistema mostra la lista delle classi che soddisfano la ricerca in base al docente selezionato.
\end{description}
\hypertarget{UC16}{}
\subsection{Caso d'uso UC16: Ricerca classe per argomento}\begin{description}
	\item[Attori:] Studente;
	\item[Scopo e descrizione:] Lo studente ricerca una classe per argomento
	\item[Precondizione:] Lo studente è autenticato presso il sistema;
	
	\item[Flusso principale degli eventi:] \ 
	\begin{enumerate}
		\item Lo studente seleziona gli argomenti di cui vuole visualizzare le classi;
		
	\end{enumerate}
	\item[Postcondizione:] Il sistema mostra la lista delle classi che soddisfano la ricerca in base agli argomenti selezionati.
\end{description}
\hypertarget{UC17}{}
\subsection{Caso d'uso UC17: Gestione domande}
\begin{figure}[H]
	\centering
	\begin{resizedtikzpicture}{\textwidth}
		\umlactor[x=0, y=0]{Docente}
		\begin{umlsystem}[x=0, fill=lightgray!20]{Quizzipedia}
			\umlusecase[x=5, y=-4, fill=white, width=4cm, name=50]{\textbf{UC17.3:} Elimina domanda}
			\umlassoc{Docente}{50}
			\umlusecase[x=5, y=0, fill=white, width=4cm, name=46]{\textbf{UC17.2:} Modifica domanda}
			\umlassoc{Docente}{46}
			\umlusecase[x=5, y=4, fill=white, width=4cm, name=43]{\textbf{UC17.1:} Inserimento domanda}
			\umlassoc{Docente}{43}
		\end{umlsystem}
	\end{resizedtikzpicture}
	\caption{\textbf{UC17}: Gestione domande}
	\label{UC17}
\end{figure}
\begin{description}
	\item[Attori:] Docente;
	\item[Scopo e descrizione:] Il docente gestisce le proprie domande
	\item[Precondizione:] Il docente è autenticato nel sistema;
	
	\item[Flusso principale degli eventi:] \ 
	\begin{enumerate}
		\item Il docente può creare una nuova domanda (\hyperlink{UC17.1}{UC17.1});
		\item Il docente può modificare una domanda (\hyperlink{UC17.2}{UC17.2});
		\item Il docente può eliminare una domanda (\hyperlink{UC17.3}{UC17.3});
		
	\end{enumerate}
	\item[Postcondizione:] Il sistema ha ottenuto le informazioni sulle operazioni che il docente desidera eseguire sulle domande.
\end{description}
\hypertarget{UC17.1}{}
\subsection{Caso d'uso UC17.1: Inserimento domanda}
\begin{figure}[H]
	\centering
	\begin{resizedtikzpicture}{\textwidth}
		\umlactor[x=0, y=0]{Docente}
		\begin{umlsystem}[x=0, fill=lightgray!20]{Quizzipedia}
			\umlusecase[x=5, y=-4.5, fill=white, width=4cm, name=162]{\textbf{UC17.1.5:} Scrittura nuova domanda da interfaccia grafica}
			\umlassoc{Docente}{162}
			\umlusecase[x=5, y=0, fill=white, width=4cm, name=114]{\textbf{UC17.1.2:} Scrittura domanda in QML della nuova domanda}
			\umlassoc{Docente}{114}
			\umlusecase[x=5, y=4.5, fill=white, width=4cm, name=112]{\textbf{UC17.1.1:} Selezione argomenti nuova domanda}
			\umlassoc{Docente}{112}
		\end{umlsystem}
	\end{resizedtikzpicture}
	\caption{\textbf{UC17.1}: Inserimento domanda}
	\label{UC17.1}
\end{figure}

\begin{figure}[H]
	\centering
	\begin{resizedtikzpicture}{\textwidth}
		\umlactor[x=0, y=0]{Docente}
		\begin{umlsystem}[x=0, fill=lightgray!20]{Quizzipedia}
			\umlusecase[x=5, y=0, fill=white, width=4cm, name=43]{\textbf{UC17.1:} Inserimento domanda}
			\umlassoc{Docente}{43}
			\umlusecase[x=15, y=-11.25, fill=white, width=4cm, name=153]{\textbf{UC17.9:} Inserimento domanda a risposta aperta}
			\umlinherit{153}{43}
			\umlusecase[x=15, y=-6.75, fill=white, width=4cm, name=152]{\textbf{UC17.8:} Inserimento domanda con l'associazione di parole}
			\umlinherit{152}{43}
			\umlusecase[x=15, y=-2.25, fill=white, width=4cm, name=151]{\textbf{UC17.7:} Inserimento domanda di tipo testo con parole omesse}
			\umlinherit{151}{43}
			\umlusecase[x=15, y=2.25, fill=white, width=4cm, name=150]{\textbf{UC17.6:} Inserimento domanda a risposta multipla}
			\umlinherit{150}{43}
			\umlusecase[x=15, y=6.75, fill=white, width=4cm, name=149]{\textbf{UC17.5:} Inserimento domanda a scelta multipla}
			\umlinherit{149}{43}
			\umlusecase[x=15, y=11.25, fill=white, width=4cm, name=148]{\textbf{UC17.4:} Inserimento domanda di tipo vero/falso}
			\umlinherit{148}{43}
		\end{umlsystem}
	\end{resizedtikzpicture}
	\caption{\textbf{UC17.1}: Inserimento domanda}
	\label{UC17.1}
\end{figure}
\begin{description}
	\item[Attori:] Docente;
	\item[Scopo e descrizione:] Il docente compone una domanda in linguaggio QML, che verrà salvata nel sistema e potrà essere utilizzata nei questionari
	\item[Precondizione:] Il docente è autenticato nel sistema;
	
	\item[Flusso principale degli eventi:] \ 
	\begin{enumerate}
		\item Il docente seleziona gli argomenti della domanda (\hyperlink{UC17.1.1}{UC17.1.1});
		\item Il docente compone la domanda in QML  (\hyperlink{UC17.1.2}{UC17.1.2});
		\item Il docente compone la domanda attraverso l'interfaccia grafica (\hyperlink{UC17.1.5}{UC17.1.5});
		\item Il docente conferma la creazione della domanda;
		
	\end{enumerate}
	\item[Postcondizione:] È stata creata una nuova domanda.
\end{description}
\hypertarget{UC17.1.1}{}
\subsection{Caso d'uso UC17.1.1: Selezione argomenti nuova domanda}
\begin{figure}[H]
	\centering
	\begin{resizedtikzpicture}{\textwidth}
		\umlactor[x=0, y=0]{Docente}
		\begin{umlsystem}[x=0, fill=lightgray!20]{Quizzipedia}
			\umlusecase[x=5, y=0, fill=white, width=4cm, name=112]{\textbf{UC17.1.1:} Selezione argomenti nuova domanda}
			\umlassoc{Docente}{112}
			\umlusecase[x=15, y=0, fill=white, width=4cm, name=143]{\textbf{UC17.1.4:} Errore argomento mancante}
			\umlextend[name=ext-143-112]{143}{112}
			\umlnote[x=23, y=0, fill=white, width=4cm, geometry=|-|, arm=-2cm]{ext-143-112-1}{Non è stato selezionato almeno un argomento}
		\end{umlsystem}
	\end{resizedtikzpicture}
	\caption{\textbf{UC17.1.1}: Selezione argomenti nuova domanda}
	\label{UC17.1.1}
\end{figure}
\begin{description}
	\item[Attori:] Docente;
	\item[Scopo e descrizione:] Il docente seleziona gli argomenti corrispondenti alla domanda selezionata
	\item[Precondizione:] Il docente sta creando una nuova domanda;
	
	\item[Flusso principale degli eventi:] \ 
	\begin{enumerate}
		\item Il docente seleziona gli argomenti della domanda;
		
	\end{enumerate}
	\item[Estensioni:]
	\begin{enumerate}
		\item Se non è stato selezionato nessun argomento viene visualizzato un messaggio di errore (\hyperlink{UC17.1.4}{UC17.1.4});
		
	\end{enumerate}
	\item[Postcondizione:] Il docente ha definito gli argomenti per classificare la domanda
	.
\end{description}
\hypertarget{UC17.1.2}{}
\subsection{Caso d'uso UC17.1.2: Scrittura domanda in QML della nuova domanda}
\begin{figure}[H]
	\centering
	\begin{resizedtikzpicture}{\textwidth}
		\umlactor[x=0, y=0]{Docente}
		\begin{umlsystem}[x=0, fill=lightgray!20]{Quizzipedia}
			\umlusecase[x=5, y=0, fill=white, width=4cm, name=114]{\textbf{UC17.1.2:} Scrittura domanda in QML della nuova domanda}
			\umlassoc{Docente}{114}
			\umlusecase[x=15, y=-2.25, fill=white, width=4cm, name=142]{\textbf{UC17.1.3:} Errore QML non valido}
			\umlextend[name=ext-142-114]{142}{114}
			\umlnote[x=23, y=-2.25, fill=white, width=4cm, geometry=|-|, arm=-1.75cm]{ext-142-114-1}{Il QML non è valido}
			\umlusecase[x=15, y=2.25, fill=white, width=4cm, name=162]{\textbf{UC17.1.5:} Scrittura nuova domanda da interfaccia grafica}
			\umlinherit{162}{114}
		\end{umlsystem}
	\end{resizedtikzpicture}
	\caption{\textbf{UC17.1.2}: Scrittura domanda in QML della nuova domanda}
	\label{UC17.1.2}
\end{figure}
\begin{description}
	\item[Attori:] Docente;
	\item[Scopo e descrizione:] Il docente compone una domanda in linguaggio QML
	\item[Precondizione:] Il docente sta creando una nuova domanda;
	
	\item[Flusso principale degli eventi:] \ 
	\begin{enumerate}
		\item Il docente compone la domanda in QML;
		
	\end{enumerate}
	\item[Estensioni:]
	\begin{enumerate}
		\item Se il codice QML inserito non è valido viene visualizzato un messaggio di errore (\hyperlink{UC17.1.3}{UC17.1.3});
		
	\end{enumerate}
	\item[Postcondizione:] La domanda è stata scritta in linguaggio QML.
\end{description}
\hypertarget{UC17.1.3}{}
\subsection{Caso d'uso UC17.1.3: Errore QML non valido}\begin{description}
	\item[Attori:] Docente;
	\item[Scopo e descrizione:] Il sistema avvisa il docente di uno o più errori nel codice QML inserito
	\item[Precondizione:] Il QML non è valido;
	
	\item[Flusso principale degli eventi:] \ 
	\begin{enumerate}
		\item Viene visualizzato un messaggio d'errore;
		
	\end{enumerate}
	\item[Postcondizione:] Non viene inserita la domanda.
\end{description}
\hypertarget{UC17.1.4}{}
\subsection{Caso d'uso UC17.1.4: Errore argomento mancante}\begin{description}
	\item[Attori:] Docente;
	\item[Scopo e descrizione:] Il sistema avvisa il docente che deve essere inserito almeno un argomento
	\item[Precondizione:] Non è stato selezionato almeno un argomento;
	
	\item[Flusso principale degli eventi:] \ 
	\begin{enumerate}
		\item Viene visualizzato un messaggio d'errore;
		
	\end{enumerate}
	\item[Postcondizione:] Non viene inserita la domanda.
\end{description}
\hypertarget{UC17.1.5}{}
\subsection{Caso d'uso UC17.1.5: Scrittura nuova domanda da interfaccia grafica}\begin{description}
	\item[Attori:] Docente;
	\item[Scopo e descrizione:] Il docente compone la domanda attraverso un interfaccia grafica
	\item[Precondizione:] Il docente sta creando una nuova domanda;
	
	\item[Flusso principale degli eventi:] \ 
	\begin{enumerate}
		\item Il docente compone la domanda attraverso l'interfaccia grafica;
		
	\end{enumerate}
	\item[Postcondizione:] La domanda è stata composta attraverso l'interfaccia grafica.
\end{description}
\hypertarget{UC17.2}{}
\subsection{Caso d'uso UC17.2: Modifica domanda}
\begin{figure}[H]
	\centering
	\begin{resizedtikzpicture}{\textwidth}
		\umlactor[x=0, y=0]{Docente}
		\begin{umlsystem}[x=0, fill=lightgray!20]{Quizzipedia}
			\umlusecase[x=5, y=-4.75, fill=white, width=4cm, name=165]{\textbf{UC17.2.3:} Modifica domanda da interfaccia grafica}
			\umlassoc{Docente}{165}
			\umlusecase[x=5, y=0, fill=white, width=4cm, name=164]{\textbf{UC17.2.2:} Scrittura domanda in QML della domanda da modificare}
			\umlassoc{Docente}{164}
			\umlusecase[x=5, y=4.75, fill=white, width=4cm, name=163]{\textbf{UC17.2.1:} Selezione argomenti modifica domanda}
			\umlassoc{Docente}{163}
		\end{umlsystem}
	\end{resizedtikzpicture}
	\caption{\textbf{UC17.2}: Modifica domanda}
	\label{UC17.2}
\end{figure}
\begin{description}
	\item[Attori:] Docente;
	\item[Scopo e descrizione:] Il docente effettua delle modifiche ad una domanda 
	\item[Precondizione:] Il docente è autenticato nel sistema;
	
	\item[Flusso principale degli eventi:] \ 
	\begin{enumerate}
		\item Ricerca della domanda da modificare (\hyperlink{UC25}{UC25});
		\item Selezione della domanda da modificare;
		\item Il docente seleziona gli argomenti della domanda	 (\hyperlink{UC17.2.1}{UC17.2.1});
		\item Il docente modifica la domanda in QML	 (\hyperlink{UC17.2.2}{UC17.2.2});
		\item Il docente compone la domanda attraverso l'interfaccia grafica (\hyperlink{UC17.2.3}{UC17.2.3});
		\item Conferma delle modifiche;
		
	\end{enumerate}
	\item[Postcondizione:] La domanda è stata modificata.
\end{description}
\hypertarget{UC17.2.1}{}
\subsection{Caso d'uso UC17.2.1: Selezione argomenti modifica domanda}
\begin{figure}[H]
	\centering
	\begin{resizedtikzpicture}{\textwidth}
		\umlactor[x=0, y=0]{Docente}
		\begin{umlsystem}[x=0, fill=lightgray!20]{Quizzipedia}
			\umlusecase[x=5, y=0, fill=white, width=4cm, name=163]{\textbf{UC17.2.1:} Selezione argomenti modifica domanda}
			\umlassoc{Docente}{163}
			\umlusecase[x=15, y=0, fill=white, width=4cm, name=215]{\textbf{UC17.2.5:} Errore argomento mancante}
			\umlextend[name=ext-215-163]{215}{163}
			\umlnote[x=23, y=0, fill=white, width=4cm, geometry=|-|, arm=-2cm]{ext-215-163-1}{Non è stato selezionato alcun argomento
			}
		\end{umlsystem}
	\end{resizedtikzpicture}
	\caption{\textbf{UC17.2.1}: Selezione argomenti modifica domanda}
	\label{UC17.2.1}
\end{figure}
\begin{description}
	\item[Attori:] Docente;
	\item[Scopo e descrizione:] Il docente seleziona gli argomenti corrispondenti alla domanda selezionata
	\item[Precondizione:] Il docente sta modificando una domanda
	;
	
	\item[Flusso principale degli eventi:] \ 
	\begin{enumerate}
		\item Il docente seleziona gli argomenti della domanda;
		
	\end{enumerate}
	\item[Estensioni:]
	\begin{enumerate}
		\item Se la domanda non ha alcun argomento associato viene visualizzato un messaggio di errore	 (\hyperlink{UC17.2.5}{UC17.2.5});
		
	\end{enumerate}
	\item[Postcondizione:] Il docente ha definito gli argomenti per classificare la domanda.
\end{description}
\hypertarget{UC17.2.2}{}
\subsection{Caso d'uso UC17.2.2: Scrittura domanda in QML della domanda da modificare}
\begin{figure}[H]
	\centering
	\begin{resizedtikzpicture}{\textwidth}
		\umlactor[x=0, y=0]{Docente}
		\begin{umlsystem}[x=0, fill=lightgray!20]{Quizzipedia}
			\umlusecase[x=5, y=0, fill=white, width=4cm, name=164]{\textbf{UC17.2.2:} Scrittura domanda in QML della domanda da modificare}
			\umlassoc{Docente}{164}
			\umlusecase[x=15, y=-2.25, fill=white, width=4cm, name=213]{\textbf{UC17.2.4:} Errore modifica QML non valido}
			\umlextend[name=ext-213-164]{213}{164}
			\umlnote[x=23, y=-2.25, fill=white, width=4cm, geometry=|-|, arm=-1.75cm]{ext-213-164-1}{Il QML non è valido
			}
			\umlusecase[x=15, y=2.25, fill=white, width=4cm, name=165]{\textbf{UC17.2.3:} Modifica domanda da interfaccia grafica}
			\umlinherit{165}{164}
		\end{umlsystem}
	\end{resizedtikzpicture}
	\caption{\textbf{UC17.2.2}: Scrittura domanda in QML della domanda da modificare}
	\label{UC17.2.2}
\end{figure}
\begin{description}
	\item[Attori:] Docente;
	\item[Scopo e descrizione:] Il docente compone una domanda in linguaggio QML
	\item[Precondizione:] Il docente sta modificando una domanda;
	
	\item[Flusso principale degli eventi:] \ 
	\begin{enumerate}
		\item Il docente compone la domanda in QML;
		
	\end{enumerate}
	\item[Estensioni:]
	\begin{enumerate}
		\item Se il codice QML inserito non è valido viene visualizzato un messaggio di errore (\hyperlink{UC17.2.4}{UC17.2.4});
		
	\end{enumerate}
	\item[Postcondizione:] La domanda è stata scritta in linguaggio QML.
\end{description}
\hypertarget{UC17.2.3}{}
\subsection{Caso d'uso UC17.2.3: Modifica domanda da interfaccia grafica}\begin{description}
	\item[Attori:] Docente;
	\item[Scopo e descrizione:] Il docente modifica la domanda attraverso un interfaccia grafica
	
	\item[Precondizione:] Il docente sta modificando una domanda
	;
	
	\item[Flusso principale degli eventi:] \ 
	\begin{enumerate}
		\item Il docente modifica la domanda attraverso l'interfaccia grafica;
		
	\end{enumerate}
	\item[Postcondizione:] La domanda è stata modificata attraverso l'interfaccia grafica.
\end{description}
\hypertarget{UC17.2.4}{}
\subsection{Caso d'uso UC17.2.4: Errore modifica QML non valido}\begin{description}
	\item[Attori:] Docente;
	\item[Scopo e descrizione:] Il sistema avvisa il docente di uno o più errori nel codice QML inserito
	
	\item[Precondizione:] Il QML non è valido
	;
	
	\item[Flusso principale degli eventi:] \ 
	\begin{enumerate}
		\item Viene visualizzato un messaggio d'errore;
		
	\end{enumerate}
	\item[Postcondizione:] Non viene modificata la domanda
	.
\end{description}
\hypertarget{UC17.2.5}{}
\subsection{Caso d'uso UC17.2.5: Errore argomento mancante}\begin{description}
	\item[Attori:] Docente;
	\item[Scopo e descrizione:] Il sistema avvisa il docente che deve essere inserito almeno un argomento
	
	\item[Precondizione:] Non è stato selezionato alcun argomento
	;
	
	\item[Flusso principale degli eventi:] \ 
	\begin{enumerate}
		\item Viene visualizzato un messaggio d'errore	;
		
	\end{enumerate}
	\item[Postcondizione:] Non viene modificata la domanda
	.
\end{description}
\hypertarget{UC17.3}{}
\subsection{Caso d'uso UC17.3: Elimina domanda}
\begin{figure}[H]
	\centering
	\begin{resizedtikzpicture}{\textwidth}
		\umlactor[x=0, y=0]{Docente}
		\begin{umlsystem}[x=0, fill=lightgray!20]{Quizzipedia}
			\umlusecase[x=5, y=0, fill=white, width=4cm, name=50]{\textbf{UC17.3:} Elimina domanda}
			\umlassoc{Docente}{50}
			\umlusecase[x=15, y=0, fill=white, width=4cm, name=206]{\textbf{UC17.11:} Errore eliminazione domanda utilizzata}
			\umlextend[name=ext-206-50]{206}{50}
			\umlnote[x=23, y=0, fill=white, width=4cm, geometry=|-|, arm=-2.25cm]{ext-206-50-1}{Il docente tenta di eliminare una domanda utilizzata da almeno un questionario}
		\end{umlsystem}
	\end{resizedtikzpicture}
	\caption{\textbf{UC17.3}: Elimina domanda}
	\label{UC17.3}
\end{figure}
\begin{description}
	\item[Attori:] Docente;
	\item[Scopo e descrizione:] Il docente rimuove dal sistema una domanda da lui creata
	\item[Precondizione:] Il docente è autenticato nel sistema;
	
	\item[Flusso principale degli eventi:] \ 
	\begin{enumerate}
		\item Ricerca della domanda da eliminare (\hyperlink{UC25}{UC25});
		\item Selezione della domanda da eliminare;
		\item Conferma per l'eliminazione della domanda	;
		
	\end{enumerate}
	\item[Estensioni:]
	\begin{enumerate}
		\item Se il docente tenta di eliminare una domanda utilizzata da almeno un questionario viene mostrato un messaggio di errore (\hyperlink{UC17.11}{UC17.11});
		
	\end{enumerate}
	\item[Postcondizione:] La domanda è stata eliminata.
\end{description}
\hypertarget{UC17.4}{}
\subsection{Caso d'uso UC17.4: Inserimento domanda di tipo vero/falso}\begin{description}
	\item[Attori:] Docente;
	\item[Scopo e descrizione:] Il docente inserisce una domanda di tipo vero/falso
	\item[Precondizione:] Il docente è autenticato presso il sistema;
	
	\item[Flusso principale degli eventi:] \ 
	\begin{enumerate}
		\item Il docente seleziona gli argomenti della domanda (\hyperlink{UC17.1.1}{UC17.1.1});
		\item Il docente compone la domanda vero/falso in QML specificando la risposta esatta (\hyperlink{UC17.1.2}{UC17.1.2});
		\item Il docente conferma la creazione della domanda;
		
	\end{enumerate}
	\item[Postcondizione:] È stata creata una nuova domanda di tipo vero/falso.
\end{description}
\hypertarget{UC17.5}{}
\subsection{Caso d'uso UC17.5: Inserimento domanda a scelta multipla}\begin{description}
	\item[Attori:] Docente;
	\item[Scopo e descrizione:] Il docente inserisce una domanda a scelta multipla
	\item[Precondizione:] Il docente è autenticato presso il sistema;
	
	\item[Flusso principale degli eventi:] \ 
	\begin{enumerate}
		\item Il docente seleziona gli argomenti della domanda (\hyperlink{UC17.1.1}{UC17.1.1});
		\item Il docente compone la domanda a scelta multipla in QML specificando testo della domanda, risposte possibili e risposta corretta (\hyperlink{UC17.1.2}{UC17.1.2});
		\item Il docente conferma la creazione della domanda;
		
	\end{enumerate}
	\item[Postcondizione:] È stata creata una nuova domanda a scelta multipla.
\end{description}
\hypertarget{UC17.6}{}
\subsection{Caso d'uso UC17.6: Inserimento domanda a risposta multipla}\begin{description}
	\item[Attori:] Docente;
	\item[Scopo e descrizione:] Il docente inserisce una domanda a risposta multipla
	\item[Precondizione:] Il docente è autenticato presso il sistema;
	
	\item[Flusso principale degli eventi:] \ 
	\begin{enumerate}
		\item Il docente seleziona gli argomenti della domanda (\hyperlink{UC17.1.1}{UC17.1.1});
		\item Il docente compone la domanda a risposta multipla in QML specificando il testo della domanda, risposte possibili e risposta/e esatte (\hyperlink{UC17.1.2}{UC17.1.2});
		\item Il docente conferma la creazione della domanda;
		
	\end{enumerate}
	\item[Postcondizione:] È stata creata una nuova domanda a risposta multipla.
\end{description}
\hypertarget{UC17.7}{}
\subsection{Caso d'uso UC17.7: Inserimento domanda di tipo testo con parole omesse}\begin{description}
	\item[Attori:] Docente;
	\item[Scopo e descrizione:] Il docente inserisce una domanda di tipo testo con parole omesse
	\item[Precondizione:] Il docente è autenticato presso il sistema;
	
	\item[Flusso principale degli eventi:] \ 
	\begin{enumerate}
		\item Il docente seleziona gli argomenti della domanda (\hyperlink{UC17.1.1}{UC17.1.1});
		\item Il docente compone il testo della domanda in QML specificando quali parole possano essere scelte e dove vadano messe  (\hyperlink{UC17.1.2}{UC17.1.2});
		\item Il docente conferma la creazione della domanda;
		
	\end{enumerate}
	\item[Postcondizione:] È stata creata una domanda di tipo testo con parole omesse.
\end{description}
\hypertarget{UC17.8}{}
\subsection{Caso d'uso UC17.8: Inserimento domanda con l'associazione di parole}\begin{description}
	\item[Attori:] Docente;
	\item[Scopo e descrizione:] Il docente inserisce una domanda con associazione di parole
	\item[Precondizione:] Il docente è autenticato presso il sistema;
	
	\item[Flusso principale degli eventi:] \ 
	\begin{enumerate}
		\item Il docente seleziona gli argomenti della domanda (\hyperlink{UC17.1.1}{UC17.1.1});
		\item Il docente compone la domanda in QML specificando le parole che possono essere combinate e le giuste combinazioni (\hyperlink{UC17.1.2}{UC17.1.2});
		\item Il docente conferma la creazione della domanda;
		
	\end{enumerate}
	\item[Postcondizione:] Stata inserita una domanda con associazione di parole.
\end{description}
\hypertarget{UC17.9}{}
\subsection{Caso d'uso UC17.9: Inserimento domanda a risposta aperta}\begin{description}
	\item[Attori:] Docente;
	\item[Scopo e descrizione:] Il docente inserisce una domanda a risposta aperta
	\item[Precondizione:] Il docente è autenticato presso il sistema;
	
	\item[Flusso principale degli eventi:] \ 
	\begin{enumerate}
		\item Il docente seleziona gli argomenti della domanda (\hyperlink{UC17.1.1}{UC17.1.1});
		\item Il docente compone la domanda in QML specificando il testo della domanda (ma non la risposta) (\hyperlink{UC17.1.2}{UC17.1.2});
		\item Il docente conferma la creazione della domanda;
		
	\end{enumerate}
	\item[Postcondizione:] È stata inserita una domanda a risposta aperta.
\end{description}
\hypertarget{UC17.10}{}
\subsection{Caso d'uso UC17.10: Visualizza domanda}\begin{description}
	\item[Attori:] Docente;
	\item[Scopo e descrizione:] Il docente visualizza una domanda
	\item[Precondizione:] Il docente è autenticato nel sistema
	;
	
	\item[Flusso principale degli eventi:] \ 
	\begin{enumerate}
		\item Ricerca della domanda da visualizzare (\hyperlink{UC25}{UC25});
		\item Selezione della domanda da visualizzare;
		
	\end{enumerate}
	\item[Postcondizione:] Il docente visualizza una domanda.
\end{description}
\hypertarget{UC17.11}{}
\subsection{Caso d'uso UC17.11: Errore eliminazione domanda utilizzata}\begin{description}
	\item[Attori:] Docente;
	\item[Scopo e descrizione:] Viene mostrato un errore nel caso in cui il docente tenti di eliminare una domanda utilizzata da almeno un questionario
	\item[Precondizione:] Il docente tenta di eliminare una domanda utilizzata da almeno un questionario;
	
	\item[Flusso principale degli eventi:] \ 
	\begin{enumerate}
		\item Viene mostrato un errore nel caso in cui il docente tenti di eliminare una domanda utilizzata da almeno un questionario;
		
	\end{enumerate}
	\item[Postcondizione:] Non viene eliminata la domanda e viene mostrato un messaggio di errore.
\end{description}
\hypertarget{UC18}{}
\subsection{Caso d'uso UC18: Gestione questionari}
\begin{figure}[H]
	\centering
	\begin{resizedtikzpicture}{\textwidth}
		\umlactor[x=0, y=0]{Docente}
		\begin{umlsystem}[x=0, fill=lightgray!20]{Quizzipedia}
			\umlusecase[x=5, y=-4, fill=white, width=4cm, name=48]{\textbf{UC18.3:} Elimina questionario}
			\umlassoc{Docente}{48}
			\umlusecase[x=5, y=0, fill=white, width=4cm, name=47]{\textbf{UC18.2:} Modifica questionario}
			\umlassoc{Docente}{47}
			\umlusecase[x=5, y=4, fill=white, width=4cm, name=42]{\textbf{UC18.1:} Inserisci questionario}
			\umlassoc{Docente}{42}
		\end{umlsystem}
	\end{resizedtikzpicture}
	\caption{\textbf{UC18}: Gestione questionari}
	\label{UC18}
\end{figure}
\begin{description}
	\item[Attori:] Docente;
	\item[Scopo e descrizione:] Il docente gestisce i propri questionari
	\item[Precondizione:] Il docente è autenticato nel sistema;
	
	\item[Flusso principale degli eventi:] \ 
	\begin{enumerate}
		\item Il docente può creare un nuovo questionario (\hyperlink{UC18.1}{UC18.1});
		\item Il docente può modificare un questionario (\hyperlink{UC18.2}{UC18.2});
		\item Il docente può eliminare un questionario (\hyperlink{UC18.3}{UC18.3});
		
	\end{enumerate}
	\item[Postcondizione:] Il sistema ha ottenuto le informazioni sulle operazioni che il docente desidera eseguire su un questionario.
\end{description}
\hypertarget{UC18.1}{}
\subsection{Caso d'uso UC18.1: Inserisci questionario}
\begin{figure}[H]
	\centering
	\begin{resizedtikzpicture}{\textwidth}
		\umlactor[x=0, y=0]{Docente}
		\begin{umlsystem}[x=0, fill=lightgray!20]{Quizzipedia}
			\umlusecase[x=5, y=-4.5, fill=white, width=4cm, name=118]{\textbf{UC18.1.3:} Seleziona argomenti del nuovo questionario}
			\umlassoc{Docente}{118}
			\umlusecase[x=5, y=0, fill=white, width=4cm, name=93]{\textbf{UC18.1.2:} Elimina domanda da un nuovo questionario}
			\umlassoc{Docente}{93}
			\umlusecase[x=5, y=4.5, fill=white, width=4cm, name=91]{\textbf{UC18.1.1:} Aggiungi domanda in un nuovo questionario }
			\umlassoc{Docente}{91}
		\end{umlsystem}
	\end{resizedtikzpicture}
	\caption{\textbf{UC18.1}: Inserisci questionario}
	\label{UC18.1}
\end{figure}

\begin{figure}[H]
	\centering
	\begin{resizedtikzpicture}{\textwidth}
		\umlactor[x=0, y=0]{Docente}
		\begin{umlsystem}[x=0, fill=lightgray!20]{Quizzipedia}
			\umlusecase[x=5, y=0, fill=white, width=4cm, name=42]{\textbf{UC18.1:} Inserisci questionario}
			\umlassoc{Docente}{42}
			\umlusecase[x=15, y=0, fill=white, width=4cm, name=144]{\textbf{UC18.4:} Errore questionario vuoto}
			\umlextend[name=ext-144-42]{144}{42}
			\umlnote[x=23, y=0, fill=white, width=4cm, geometry=|-|, arm=-2cm]{ext-144-42-1}{Il questionario selezionato non ha domande}
		\end{umlsystem}
	\end{resizedtikzpicture}
	\caption{\textbf{UC18.1}: Inserisci questionario}
	\label{UC18.1}
\end{figure}
\begin{description}
	\item[Attori:] Docente;
	\item[Scopo e descrizione:] Il docente crea un questionario
	\item[Precondizione:] Il docente è autenticato nel sistema;
	
	\item[Flusso principale degli eventi:] \ 
	\begin{enumerate}
		\item Il docente può aggiungere domande al questionario (\hyperlink{UC18.1.1}{UC18.1.1});
		\item Il docente può togliere domande precedentemente aggiunte al questionario (\hyperlink{UC18.1.2}{UC18.1.2});
		\item Il docente seleziona gli argomenti del questionario (\hyperlink{UC18.1.3}{UC18.1.3});
		\item Il docente aggiunge il titolo del questionario;
		\item Il docente conferma la creazione del questionario;
		
	\end{enumerate}
	\item[Estensioni:]
	\begin{enumerate}
		\item Se il questionario non ha domande viene visualizzato un messaggio d'errore (\hyperlink{UC18.4}{UC18.4});
		
	\end{enumerate}
	\item[Postcondizione:] È stato creato un nuovo questionario.
\end{description}
\hypertarget{UC18.1.1}{}
\subsection{Caso d'uso UC18.1.1: Aggiungi domanda in un nuovo questionario }\begin{description}
	\item[Attori:] Docente;
	\item[Scopo e descrizione:] Il docente ricerca e seleziona una domanda da inserire in un nuovo questionario
	\item[Precondizione:] Il docente sta creando un nuovo questionario;
	
	\item[Flusso principale degli eventi:] \ 
	\begin{enumerate}
		\item Viene ricercata un domanda (\hyperlink{UC25}{UC25});
		\item Selezione della domanda;
		\item Conferma inserimento domanda;
		
	\end{enumerate}
	\item[Postcondizione:] È stata aggiunta una domanda al nuovo questionario.
\end{description}
\hypertarget{UC18.1.2}{}
\subsection{Caso d'uso UC18.1.2: Elimina domanda da un nuovo questionario}\begin{description}
	\item[Attori:] Docente;
	\item[Scopo e descrizione:] Il docente elimina una domanda da un nuovo questionario
	\item[Precondizione:] Il docente sta creando un nuovo questionario e ha selezionato una domanda da eliminare;
	
	\item[Flusso principale degli eventi:] \ 
	\begin{enumerate}
		\item Viene eliminata la domanda;
		
	\end{enumerate}
	\item[Postcondizione:] È stato eliminata la domanda dal nuovo questionario.
\end{description}
\hypertarget{UC18.1.3}{}
\subsection{Caso d'uso UC18.1.3: Seleziona argomenti del nuovo questionario}\begin{description}
	\item[Attori:] Docente;
	\item[Scopo e descrizione:] Il docente seleziona gli argomenti relativi al nuovo questionario 
	\item[Precondizione:] Il docente sta creando un nuovo questionario;
	
	\item[Flusso principale degli eventi:] \ 
	\begin{enumerate}
		\item Il docente aggiunge un argomento al questionario;
		\item Il docente toglie un argomento al questionario;
		
	\end{enumerate}
	\item[Postcondizione:] Il docente ha selezionato gli argomenti del nuovo questionario.
\end{description}
\hypertarget{UC18.2}{}
\subsection{Caso d'uso UC18.2: Modifica questionario}
\begin{figure}[H]
	\centering
	\begin{resizedtikzpicture}{\textwidth}
		\umlactor[x=0, y=0]{Docente}
		\begin{umlsystem}[x=0, fill=lightgray!20]{Quizzipedia}
			\umlusecase[x=5, y=-4.5, fill=white, width=4cm, name=168]{\textbf{UC18.2.3:} Selezione argomenti modifica questionario}
			\umlassoc{Docente}{168}
			\umlusecase[x=5, y=0, fill=white, width=4cm, name=167]{\textbf{UC18.2.2:} Elimina domanda da un questionario da  modificare}
			\umlassoc{Docente}{167}
			\umlusecase[x=5, y=4.5, fill=white, width=4cm, name=166]{\textbf{UC18.2.1:} Aggiungi domanda in un questionario}
			\umlassoc{Docente}{166}
		\end{umlsystem}
	\end{resizedtikzpicture}
	\caption{\textbf{UC18.2}: Modifica questionario}
	\label{UC18.2}
\end{figure}

\begin{figure}[H]
	\centering
	\begin{resizedtikzpicture}{\textwidth}
		\umlactor[x=0, y=0]{Docente}
		\begin{umlsystem}[x=0, fill=lightgray!20]{Quizzipedia}
			\umlusecase[x=5, y=0, fill=white, width=4cm, name=47]{\textbf{UC18.2:} Modifica questionario}
			\umlassoc{Docente}{47}
			\umlusecase[x=15, y=0, fill=white, width=4cm, name=144]{\textbf{UC18.4:} Errore questionario vuoto}
			\umlextend[name=ext-144-47]{144}{47}
			\umlnote[x=23, y=0, fill=white, width=4cm, geometry=|-|, arm=-2cm]{ext-144-47-1}{Il questionario selezionato non ha domande}
		\end{umlsystem}
	\end{resizedtikzpicture}
	\caption{\textbf{UC18.2}: Modifica questionario}
	\label{UC18.2}
\end{figure}
\begin{description}
	\item[Attori:] Docente;
	\item[Scopo e descrizione:] Il docente modifica il questionario che ha selezionato potendo aggiungere e rimuovere domande e modificare gli argomenti del questionario
	\item[Precondizione:] Il docente è autenticato nel sistema;
	
	\item[Flusso principale degli eventi:] \ 
	\begin{enumerate}
		\item Ricerca del questionario da modificare (\hyperlink{UC4}{UC4});
		\item Selezione del questionario da modificare;
		\item Il docente può aggiungere domande al questionario (\hyperlink{UC18.2.1}{UC18.2.1});
		\item Il docente può togliere domande precedentemente aggiunte al questionario (\hyperlink{UC18.2.2}{UC18.2.2});
		\item Il docente può modificare gli argomenti del questionario (\hyperlink{UC18.2.3}{UC18.2.3});
		\item Il docente conferma la modifica del questionario;
		
	\end{enumerate}
	\item[Estensioni:]
	\begin{enumerate}
		\item Se il questionario non ha domande viene visualizzato un messaggio d'errore (\hyperlink{UC18.4}{UC18.4});
		
	\end{enumerate}
	\item[Postcondizione:] È stato modificato il questionario.
\end{description}
\hypertarget{UC18.2.1}{}
\subsection{Caso d'uso UC18.2.1: Aggiungi domanda in un questionario}\begin{description}
	\item[Attori:] Docente;
	\item[Scopo e descrizione:] Il docente inserisce una domanda da un questionario da modificare
	\item[Precondizione:] Il docente ha selezionato un questionario da modificare e una domanda da inserire;
	
	\item[Flusso principale degli eventi:] \ 
	\begin{enumerate}
		\item Viene ricercata una domanda (\hyperlink{UC25}{UC25});
		\item Selezione della domanda	;
		\item Conferma inserimento domanda;
		
	\end{enumerate}
	\item[Postcondizione:] È stata aggiunta una domanda al questionario.
\end{description}
\hypertarget{UC18.2.2}{}
\subsection{Caso d'uso UC18.2.2: Elimina domanda da un questionario da  modificare}\begin{description}
	\item[Attori:] Docente;
	\item[Scopo e descrizione:] Il docente elimina una domanda da un questionario da modificare
	\item[Precondizione:] Il docente ha selezionato un questionario da modificare e una domanda da eliminare;
	
	\item[Flusso principale degli eventi:] \ 
	\begin{enumerate}
		\item Viene eliminata la domanda;
		
	\end{enumerate}
	\item[Postcondizione:] È stato eliminata la domanda questionario da modificare.
\end{description}
\hypertarget{UC18.2.3}{}
\subsection{Caso d'uso UC18.2.3: Selezione argomenti modifica questionario}\begin{description}
	\item[Attori:] Docente;
	\item[Scopo e descrizione:] Il docente seleziona gli argomenti corrispondenti al questionario selezionato
	\item[Precondizione:] Il docente sta modificando un questionario;
	
	\item[Flusso principale degli eventi:] \ 
	\begin{enumerate}
		\item Il docente aggiunge un argomento al questionario;
		\item Il docente toglie un argomento al questionario;
		
	\end{enumerate}
	\item[Postcondizione:] Il docente ha definito gli argomenti per classificare il questionario.
\end{description}
\hypertarget{UC18.3}{}
\subsection{Caso d'uso UC18.3: Elimina questionario}\begin{description}
	\item[Attori:] Docente;
	\item[Scopo e descrizione:] Il docente rimuove un questionario dal sistema
	\item[Precondizione:] Il docente è autenticato nel sistema;
	
	\item[Flusso principale degli eventi:] \ 
	\begin{enumerate}
		\item Ricerca il questionario da eliminare (\hyperlink{UC4}{UC4});
		\item Selezione del questionario da eliminare;
		\item Il docente conferma l'eliminazione;
		
	\end{enumerate}
	\item[Postcondizione:] È stato eliminato il questionario.
\end{description}
\hypertarget{UC18.4}{}
\subsection{Caso d'uso UC18.4: Errore questionario vuoto}\begin{description}
	\item[Attori:] Docente;
	\item[Scopo e descrizione:] Il sistema avvisa il docente che nel questionario deve essere presente almeno una domanda
	\item[Precondizione:] Il questionario selezionato non ha domande;
	
	\item[Flusso principale degli eventi:] \ 
	\begin{enumerate}
		\item Viene visualizzato un messaggio di errore;
		
	\end{enumerate}
	\item[Postcondizione:] Il questionario non viene inserito nel sistema.
\end{description}
\hypertarget{UC18.5}{}
\subsection{Caso d'uso UC18.5: Visualizza questionario}\begin{description}
	\item[Attori:] Docente;
	\item[Scopo e descrizione:] Il docente visualizza un questionario
	\item[Precondizione:] Il docente è autenticato nel sistema
	;
	
	\item[Flusso principale degli eventi:] \ 
	\begin{enumerate}
		\item Ricerca del questionario da visualizzare	 (\hyperlink{UC4}{UC4});
		\item Selezione del questionario da visualizzare	;
		
	\end{enumerate}
	\item[Postcondizione:] Il docente ha visualizzato un questionario.
\end{description}
\hypertarget{UC19}{}
\subsection{Caso d'uso UC19: Gestione classi}
\begin{figure}[H]
	\centering
	\begin{resizedtikzpicture}{\textwidth}
		\umlactor[x=0, y=0]{Docente}
		\begin{umlsystem}[x=0, fill=lightgray!20]{Quizzipedia}
			\umlusecase[x=5, y=-4, fill=white, width=4cm, name=49]{\textbf{UC19.3:} Elimina classe}
			\umlassoc{Docente}{49}
			\umlusecase[x=5, y=0, fill=white, width=4cm, name=45]{\textbf{UC19.2:} Modifica classe}
			\umlassoc{Docente}{45}
			\umlusecase[x=5, y=4, fill=white, width=4cm, name=44]{\textbf{UC19.1:} Inserisci classe}
			\umlassoc{Docente}{44}
		\end{umlsystem}
	\end{resizedtikzpicture}
	\caption{\textbf{UC19}: Gestione classi}
	\label{UC19}
\end{figure}
\begin{description}
	\item[Attori:] Docente;
	\item[Scopo e descrizione:] Il docente gestisce le proprie classi
	\item[Precondizione:] Il docente è autenticato nel sistema;
	
	\item[Flusso principale degli eventi:] \ 
	\begin{enumerate}
		\item Il docente può creare una nuova classe (\hyperlink{UC19.1}{UC19.1});
		\item Il docente può modificare una classe (\hyperlink{UC19.2}{UC19.2});
		\item Il docente può eliminare una classe (\hyperlink{UC19.3}{UC19.3});
		
	\end{enumerate}
	\item[Postcondizione:] Il sistema ha ottenuto le informazioni sulle operazioni che il docente desidera eseguire sulla classe.
\end{description}
\hypertarget{UC19.1}{}
\subsection{Caso d'uso UC19.1: Inserisci classe}
\begin{figure}[H]
	\centering
	\begin{resizedtikzpicture}{\textwidth}
		\umlactor[x=0, y=0]{Docente}
		\begin{umlsystem}[x=0, fill=lightgray!20]{Quizzipedia}
			\umlusecase[x=5, y=-4, fill=white, width=4cm, name=95]{\textbf{UC19.1.3:} Inserisci password classe}
			\umlassoc{Docente}{95}
			\umlusecase[x=5, y=0, fill=white, width=4cm, name=94]{\textbf{UC19.1.2:} Inserisci argomenti classe}
			\umlassoc{Docente}{94}
			\umlusecase[x=5, y=4, fill=white, width=4cm, name=90]{\textbf{UC19.1.1:} Inserisci nome classe}
			\umlassoc{Docente}{90}
		\end{umlsystem}
	\end{resizedtikzpicture}
	\caption{\textbf{UC19.1}: Inserisci classe}
	\label{UC19.1}
\end{figure}
\begin{description}
	\item[Attori:] Docente;
	\item[Scopo e descrizione:] Il docente crea una nuova classe alla quale gli studenti potranno iscriversi 
	\item[Precondizione:] Il docente è autenticato nel sistema;
	
	\item[Flusso principale degli eventi:] \ 
	\begin{enumerate}
		\item Viene inserito il nome della classe (\hyperlink{UC19.1.1}{UC19.1.1});
		\item Vengono inseriti gli argomenti della classe (\hyperlink{UC19.1.2}{UC19.1.2});
		\item Viene inserita la password della classe (\hyperlink{UC19.1.3}{UC19.1.3});
		\item Viene confermato l'inserimento della classe;
		
	\end{enumerate}
	\item[Postcondizione:] È stato creata una nuova classe.
\end{description}
\hypertarget{UC19.1.1}{}
\subsection{Caso d'uso UC19.1.1: Inserisci nome classe}
\begin{figure}[H]
	\centering
	\begin{resizedtikzpicture}{\textwidth}
		\umlactor[x=0, y=0]{Docente}
		\begin{umlsystem}[x=0, fill=lightgray!20]{Quizzipedia}
			\umlusecase[x=5, y=0, fill=white, width=4cm, name=90]{\textbf{UC19.1.1:} Inserisci nome classe}
			\umlassoc{Docente}{90}
			\umlusecase[x=15, y=0, fill=white, width=4cm, name=92]{\textbf{UC19.4:} Errore nome classe già presente}
			\umlextend[name=ext-92-90]{92}{90}
			\umlnote[x=23, y=0, fill=white, width=4cm, geometry=|-|, arm=-2cm]{ext-92-90-1}{Il nome della classe è già presente nel sistema}
		\end{umlsystem}
	\end{resizedtikzpicture}
	\caption{\textbf{UC19.1.1}: Inserisci nome classe}
	\label{UC19.1.1}
\end{figure}
\begin{description}
	\item[Attori:] Docente;
	\item[Scopo e descrizione:] Il docente inserisce il nome scelto per la classe
	\item[Precondizione:] Il docente ha selezionato la classe;
	
	\item[Flusso principale degli eventi:] \ 
	\begin{enumerate}
		\item Il docente inserisce un nome per la classe non ancora presente nel sistema;
		
	\end{enumerate}
	\item[Estensioni:]
	\begin{enumerate}
		\item Se il nome della classe è già presente nel sistema viene visualizzato un messaggio di errore (\hyperlink{UC19.4}{UC19.4});
		
	\end{enumerate}
	\item[Postcondizione:] Il docente ha inserito il nome della classe.
\end{description}
\hypertarget{UC19.1.2}{}
\subsection{Caso d'uso UC19.1.2: Inserisci argomenti classe}\begin{description}
	\item[Attori:] Docente;
	\item[Scopo e descrizione:] Il docente specifica di quali argomenti tratta la classe selezionata
	\item[Precondizione:] Il docente ha selezionato una classe;
	
	\item[Flusso principale degli eventi:] \ 
	\begin{enumerate}
		\item Il docente specifica di quali argomenti tratta la classe selezionata;
		
	\end{enumerate}
	\item[Postcondizione:] Il docente ha inserito gli argomenti della classe.
\end{description}
\hypertarget{UC19.1.3}{}
\subsection{Caso d'uso UC19.1.3: Inserisci password classe}\begin{description}
	\item[Attori:] Docente;
	\item[Scopo e descrizione:] Il docente specifica la password che verrà fornita poi agli studenti in classe o per altre vie per poter accedere alla classe
	\item[Precondizione:] Il docente ha selezionato una classe;
	
	\item[Flusso principale degli eventi:] \ 
	\begin{enumerate}
		\item Il docente inserisce una password che servirà per accedere alla classe che verrà creata;
		
	\end{enumerate}
	\item[Postcondizione:] Il docente ha inserito la password per accedere alla classe.
\end{description}
\hypertarget{UC19.2}{}
\subsection{Caso d'uso UC19.2: Modifica classe}\begin{description}
	\item[Attori:] Docente;
	\item[Scopo e descrizione:] Il docente modifica degli attributi o dei dati legati alla classe scelta
	\item[Precondizione:] Il docente è autenticato nel sistema, è presente nel sistema la classe da modificare;
	
	\item[Flusso principale degli eventi:] \ 
	\begin{enumerate}
		\item Viene selezionata una delle proprie classi da modificare;
		\item Può venire modificato il nome (\hyperlink{UC19.1.1}{UC19.1.1});
		\item Possono venire modificati gli argomenti (\hyperlink{UC19.1.2}{UC19.1.2});
		\item Può venire modificata la password (\hyperlink{UC19.1.3}{UC19.1.3});
		\item Viene confermata la modifica;
		
	\end{enumerate}
	\item[Postcondizione:] È stata modificata la classe.
\end{description}
\hypertarget{UC19.3}{}
\subsection{Caso d'uso UC19.3: Elimina classe}\begin{description}
	\item[Attori:] Docente;
	\item[Scopo e descrizione:] Il docente rimuove dal sistema una classe da lui creata 
	\item[Precondizione:] Il docente è autenticato nel sistema, è presente nel sistema la classe da eliminare;
	
	\item[Flusso principale degli eventi:] \ 
	\begin{enumerate}
		\item Viene sezionata una delle proprie classi;
		\item Viene sezionata la funzionalità elimina;
		\item Viene confermata l'eliminazione;
		
	\end{enumerate}
	\item[Postcondizione:] È stata eliminata la classe.
\end{description}
\hypertarget{UC19.4}{}
\subsection{Caso d'uso UC19.4: Errore nome classe già presente}\begin{description}
	\item[Attori:] Docente;
	\item[Scopo e descrizione:] Viene visualizzato un errore nel caso il nome della classe che si sta tentando di inserire esiste già
	\item[Precondizione:] Il nome della classe è già presente nel sistema;
	
	\item[Flusso principale degli eventi:] \ 
	\begin{enumerate}
		\item Viene visualizzato un messaggio di errore;
		
	\end{enumerate}
	\item[Postcondizione:] La classe non viene inserita nel sistema.
\end{description}
\hypertarget{UC20}{}
\subsection{Caso d'uso UC20: Gestione argomenti}
\begin{figure}[H]
	\centering
	\begin{resizedtikzpicture}{\textwidth}
		\umlactor[x=0, y=0]{Docente}
		\begin{umlsystem}[x=0, fill=lightgray!20]{Quizzipedia}
			\umlusecase[x=5, y=-4, fill=white, width=4cm, name=146]{\textbf{UC20.4:} Eliminazione argomento}
			\umlassoc{Docente}{146}
			\umlusecase[x=5, y=0, fill=white, width=4cm, name=141]{\textbf{UC20.3:} Modifica argomento}
			\umlassoc{Docente}{141}
			\umlusecase[x=5, y=4, fill=white, width=4cm, name=137]{\textbf{UC20.1:} Crea argomento}
			\umlassoc{Docente}{137}
		\end{umlsystem}
	\end{resizedtikzpicture}
	\caption{\textbf{UC20}: Gestione argomenti}
	\label{UC20}
\end{figure}
\begin{description}
	\item[Attori:] Docente;
	\item[Scopo e descrizione:] Il docente gestisce gli argomenti
	\item[Precondizione:] Il docente è autenticato nel sistema;
	
	\item[Flusso principale degli eventi:] \ 
	\begin{enumerate}
		\item Il docente può esplorare gli argomenti;
		\item Il docente può creare un argomento (\hyperlink{UC20.1}{UC20.1});
		\item Il docente può modificare un argomento (\hyperlink{UC20.3}{UC20.3});
		\item Il docente può eliminare un argomento (\hyperlink{UC20.4}{UC20.4});
		
	\end{enumerate}
	\item[Postcondizione:] Il sistema ha ottenuto le informazioni sulle operazioni che il docente desidera eseguire su un argomento.
\end{description}
\hypertarget{UC20.1}{}
\subsection{Caso d'uso UC20.1: Crea argomento}
\begin{figure}[H]
	\centering
	\begin{resizedtikzpicture}{\textwidth}
		\umlactor[x=0, y=0]{Docente}
		\begin{umlsystem}[x=0, fill=lightgray!20]{Quizzipedia}
			\umlusecase[x=5, y=0, fill=white, width=4cm, name=137]{\textbf{UC20.1:} Crea argomento}
			\umlassoc{Docente}{137}
			\umlusecase[x=15, y=0, fill=white, width=4cm, name=140]{\textbf{UC20.2:} Errore argomento già presente nel sistema}
			\umlextend[name=ext-140-137]{140}{137}
			\umlnote[x=23, y=0, fill=white, width=4cm, geometry=|-|, arm=-2cm]{ext-140-137-1}{L'argomento è già presente nel sistema}
		\end{umlsystem}
	\end{resizedtikzpicture}
	\caption{\textbf{UC20.1}: Crea argomento}
	\label{UC20.1}
\end{figure}
\begin{description}
	\item[Attori:] Docente;
	\item[Scopo e descrizione:] Il docente aggiunge un nuovo argomento nel sistema
	\item[Precondizione:] Il docente è autenticato nel sistema;
	
	\item[Flusso principale degli eventi:] \ 
	\begin{enumerate}
		\item Il docente inserisce il nome dell'argomento;
		\item Il docente conferma la creazione dell'argomento;
		
	\end{enumerate}
	\item[Estensioni:]
	\begin{enumerate}
		\item Se l'argomento è già presente nel sistema viene visualizzato un messaggio di errore (\hyperlink{UC20.2}{UC20.2});
		
	\end{enumerate}
	\item[Postcondizione:] L'argomento inserito è ora presente tra gli argomenti del sistema.
\end{description}
\hypertarget{UC20.2}{}
\subsection{Caso d'uso UC20.2: Errore argomento già presente nel sistema}\begin{description}
	\item[Attori:] Docente;
	\item[Scopo e descrizione:] Il sistema avvisa il docente che non è possibile inserire un argomento già presente
	\item[Precondizione:] L'argomento è già presente nel sistema;
	
	\item[Flusso principale degli eventi:] \ 
	\begin{enumerate}
		\item Viene visualizzato un messaggio d'errore;
		
	\end{enumerate}
	\item[Postcondizione:] L'argomento non viene inserito nel sistema.
\end{description}
\hypertarget{UC20.3}{}
\subsection{Caso d'uso UC20.3: Modifica argomento}
\begin{figure}[H]
	\centering
	\begin{resizedtikzpicture}{\textwidth}
		\umlactor[x=0, y=0]{Docente}
		\begin{umlsystem}[x=0, fill=lightgray!20]{Quizzipedia}
			\umlusecase[x=5, y=0, fill=white, width=4cm, name=141]{\textbf{UC20.3:} Modifica argomento}
			\umlassoc{Docente}{141}
			\umlusecase[x=15, y=0, fill=white, width=4cm, name=140]{\textbf{UC20.2:} Errore argomento già presente nel sistema}
			\umlextend[name=ext-140-141]{140}{141}
			\umlnote[x=23, y=0, fill=white, width=4cm, geometry=|-|, arm=-2cm]{ext-140-141-1}{L'argomento è già presente nel sistema}
		\end{umlsystem}
	\end{resizedtikzpicture}
	\caption{\textbf{UC20.3}: Modifica argomento}
	\label{UC20.3}
\end{figure}
\begin{description}
	\item[Attori:] Docente;
	\item[Scopo e descrizione:] Il docente modifica il nome di un argomento
	\item[Precondizione:] Il docente è autenticato nel sistema;
	
	\item[Flusso principale degli eventi:] \ 
	\begin{enumerate}
		\item Il docente seleziona un argomento da modificare;
		\item Il docente può modificare il nome dell'argomento;
		\item Il docente conferma la modifica dell'argomento;
		
	\end{enumerate}
	\item[Estensioni:]
	\begin{enumerate}
		\item Se l'argomento già presente nel sistema viene visualizzato un messaggio di errore (\hyperlink{UC20.2}{UC20.2});
		
	\end{enumerate}
	\item[Postcondizione:] È stato modificato un argomento.
\end{description}
\hypertarget{UC20.4}{}
\subsection{Caso d'uso UC20.4: Eliminazione argomento}
\begin{figure}[H]
	\centering
	\begin{resizedtikzpicture}{\textwidth}
		\umlactor[x=0, y=0]{Docente}
		\begin{umlsystem}[x=0, fill=lightgray!20]{Quizzipedia}
			\umlusecase[x=5, y=0, fill=white, width=4cm, name=146]{\textbf{UC20.4:} Eliminazione argomento}
			\umlassoc{Docente}{146}
			\umlusecase[x=15, y=0, fill=white, width=4cm, name=147]{\textbf{UC20.5:} Errore l'argomento ha domande o questionari}
			\umlextend[name=ext-147-146]{147}{146}
			\umlnote[x=23, y=0, fill=white, width=4cm, geometry=|-|, arm=-2cm]{ext-147-146-1}{L'argomento ha domande o questionari al suo interno}
		\end{umlsystem}
	\end{resizedtikzpicture}
	\caption{\textbf{UC20.4}: Eliminazione argomento}
	\label{UC20.4}
\end{figure}
\begin{description}
	\item[Attori:] Docente;
	\item[Scopo e descrizione:] Il docente elimina un argomento già presente nel sistema
	\item[Precondizione:] Il docente è autenticato nel sistema;
	
	\item[Flusso principale degli eventi:] \ 
	\begin{enumerate}
		\item Il docente seleziona l'argomento da eliminare;
		\item Il docente conferma l'eliminazione;
		
	\end{enumerate}
	\item[Estensioni:]
	\begin{enumerate}
		\item Se l'argomento ha domande o questionari viene visualizzato un messaggio d'errore (\hyperlink{UC20.5}{UC20.5});
		
	\end{enumerate}
	\item[Postcondizione:] È stato eliminato un argomento.
\end{description}
\hypertarget{UC20.5}{}
\subsection{Caso d'uso UC20.5: Errore l'argomento ha domande o questionari}\begin{description}
	\item[Attori:] Docente;
	\item[Scopo e descrizione:] Il sistema avvisa il docente che non può essere eliminato un argomento di cui esistano ancora domande o questionari
	\item[Precondizione:] L'argomento ha domande o questionari al suo interno;
	
	\item[Flusso principale degli eventi:] \ 
	\begin{enumerate}
		\item Viene visualizzato un messaggio d'errore;
		
	\end{enumerate}
	\item[Postcondizione:] L'argomento non viene eliminato.
\end{description}
\hypertarget{UC21}{}
\subsection{Caso d'uso UC21: Visualizza statistiche}
\begin{figure}[H]
	\centering
	\begin{resizedtikzpicture}{\textwidth}
		\umlactor[x=0, y=0]{Docente}
		\begin{umlsystem}[x=0, fill=lightgray!20]{Quizzipedia}
			\umlusecase[x=5, y=0, fill=white, width=4cm, name=10]{\textbf{UC21:} Visualizza statistiche}
			\umlassoc{Docente}{10}
			\umlusecase[x=15, y=-4.25, fill=white, width=4cm, name=13]{\textbf{UC24:} Visualizza statistiche classe}
			\umlinherit{13}{10}
			\umlusecase[x=15, y=0, fill=white, width=4cm, name=12]{\textbf{UC23:} Visualizza statistiche questionario}
			\umlinherit{12}{10}
			\umlusecase[x=15, y=4.25, fill=white, width=4cm, name=11]{\textbf{UC22:} Visualizza statistiche domanda}
			\umlinherit{11}{10}
		\end{umlsystem}
	\end{resizedtikzpicture}
	\caption{\textbf{UC21}: Visualizza statistiche}
	\label{UC21}
\end{figure}
\begin{description}
	\item[Attori:] Docente;
	\item[Scopo e descrizione:] Il docente visualizza le statistiche
	\item[Precondizione:] Il docente è autenticato presso il sistema;
	
	\item[Flusso principale degli eventi:] \ 
	\begin{enumerate}
		\item Il docente può visualizzare le statistiche;
		
	\end{enumerate}
	\item[Postcondizione:] Il docente ha visualizzato le statistiche a cui era interessato.
\end{description}
\hypertarget{UC22}{}
\subsection{Caso d'uso UC22: Visualizza statistiche domanda}\begin{description}
	\item[Attori:] Docente;
	\item[Scopo e descrizione:] Il docente visualizza il numero di risposte totali, risposte corrette, risposte errate e la percentuale di risposte corrette sul totale
	\item[Precondizione:] Il docente è autenticato presso il sistema;
	
	\item[Flusso principale degli eventi:] \ 
	\begin{enumerate}
		\item Il docente ricerca una domanda (\hyperlink{UC25}{UC25});
		\item Il docente seleziona la domanda interessata;
		
	\end{enumerate}
	\item[Postcondizione:] Il docente ha visualizzato le statistiche relative alla domanda a cui era interessato.
\end{description}
\hypertarget{UC23}{}
\subsection{Caso d'uso UC23: Visualizza statistiche questionario}\begin{description}
	\item[Attori:] Docente;
	\item[Scopo e descrizione:] Il docente visualizza per ogni punteggio la percentuale degli studenti che ottenuto tale punteggio e la media del punteggio
	\item[Precondizione:] Il docente è autenticato presso il sistema;
	
	\item[Flusso principale degli eventi:] \ 
	\begin{enumerate}
		\item Il docente ricerca un questionario (\hyperlink{UC4}{UC4});
		\item Il docente seleziona il questionario interessato;
		
	\end{enumerate}
	\item[Postcondizione:] Il docente ha visualizzato le statistiche relative al questionario a cui era interessato.
\end{description}
\hypertarget{UC24}{}
\subsection{Caso d'uso UC24: Visualizza statistiche classe}
\begin{figure}[H]
	\centering
	\begin{resizedtikzpicture}{\textwidth}
		\umlactor[x=0, y=0]{Docente}
		\begin{umlsystem}[x=0, fill=lightgray!20]{Quizzipedia}
			\umlusecase[x=5, y=-6.75, fill=white, width=4cm, name=110]{\textbf{UC24.4:} Visualizza statistiche studente della classe}
			\umlassoc{Docente}{110}
			\umlusecase[x=5, y=-2.25, fill=white, width=4cm, name=109]{\textbf{UC24.3:} Visualizza sommario statistiche classe}
			\umlassoc{Docente}{109}
			\umlusecase[x=5, y=2.25, fill=white, width=4cm, name=108]{\textbf{UC24.2:} Visualizza risultati questionari della classe}
			\umlassoc{Docente}{108}
			\umlusecase[x=5, y=6.75, fill=white, width=4cm, name=107]{\textbf{UC24.1:} Visualizza risultati domande della classe}
			\umlassoc{Docente}{107}
		\end{umlsystem}
	\end{resizedtikzpicture}
	\caption{\textbf{UC24}: Visualizza statistiche classe}
	\label{UC24}
\end{figure}
\begin{description}
	\item[Attori:] Docente;
	\item[Scopo e descrizione:] Il docente visualizza le statistiche di una sua classe
	\item[Precondizione:] Il docente è autenticato presso il sistema;
	
	\item[Flusso principale degli eventi:] \ 
	\begin{enumerate}
		\item Il docente seleziona la classe di cui vuole visualizzare le statistiche;
		\item Il docente può visualizzare i risultati delle sue domande della classe (\hyperlink{UC24.1}{UC24.1});
		\item Il docente può visualizzare i risultati dei suoi questionari della classe (\hyperlink{UC24.2}{UC24.2});
		\item Il docente può visualizzare un sommario delle statistiche della classe (\hyperlink{UC24.3}{UC24.3});
		\item Il docente può visualizzare i risultati dei test di uno studente della classe (\hyperlink{UC24.4}{UC24.4});
		
	\end{enumerate}
	\item[Postcondizione:] Il docente ha visualizzato le statistiche relative ad una delle sue classi.
\end{description}
\hypertarget{UC24.1}{}
\subsection{Caso d'uso UC24.1: Visualizza risultati domande della classe}\begin{description}
	\item[Attori:] Docente;
	\item[Scopo e descrizione:] Il docente visualizza i risultati e le statistiche relative alle domande della classe selezionata
	\item[Precondizione:] Il docente ha selezionato una classe;
	
	\item[Flusso principale degli eventi:] \ 
	\begin{enumerate}
		\item Il docente cerca una domanda tra quelle della classe (\hyperlink{UC25}{UC25});
		\item Il docente seleziona la domanda per visualizzarne i risultati e le statistiche;
		
	\end{enumerate}
	\item[Postcondizione:] Il docente ha visualizzato i risultati e le statistiche relative alle domande desiderate.
\end{description}
\hypertarget{UC24.2}{}
\subsection{Caso d'uso UC24.2: Visualizza risultati questionari della classe}\begin{description}
	\item[Attori:] Docente;
	\item[Scopo e descrizione:] Il docente visualizza i risultati e le statistiche relative ai questionari della classe
	\item[Precondizione:] Il docente ha selezionato una classe;
	
	\item[Flusso principale degli eventi:] \ 
	\begin{enumerate}
		\item Il docente cerca un questionario tra quelli della classe (\hyperlink{UC4}{UC4});
		\item Il docente seleziona il questionario per visualizzarne i risultati e le statistiche;
		
	\end{enumerate}
	\item[Postcondizione:] Il docente ha visualizzato i risultati e le statistiche relative ai questionari desiderate.
\end{description}
\hypertarget{UC24.3}{}
\subsection{Caso d'uso UC24.3: Visualizza sommario statistiche classe}\begin{description}
	\item[Attori:] Docente;
	\item[Scopo e descrizione:] Il docente visualizza le statistiche generali relative alla classe selezionata
	\item[Precondizione:] Il docente ha selezionato una classe;
	
	\item[Flusso principale degli eventi:] \ 
	\begin{enumerate}
		\item Il docente seleziona la funzionalità per visualizzare il sommario delle statistiche della classe;
		
	\end{enumerate}
	\item[Postcondizione:] Il docente ha visualizzato le statistiche generali relative alla classe selezionata.
\end{description}
\hypertarget{UC24.4}{}
\subsection{Caso d'uso UC24.4: Visualizza statistiche studente della classe}
\begin{figure}[H]
	\centering
	\begin{resizedtikzpicture}{\textwidth}
		\umlactor[x=0, y=0]{Docente}
		\begin{umlsystem}[x=0, fill=lightgray!20]{Quizzipedia}
			\umlusecase[x=5, y=0, fill=white, width=4cm, name=111]{\textbf{UC24.4.1:} Visualizza risultati questionari dello studente}
			\umlassoc{Docente}{111}
		\end{umlsystem}
	\end{resizedtikzpicture}
	\caption{\textbf{UC24.4}: Visualizza statistiche studente della classe}
	\label{UC24.4}
\end{figure}
\begin{description}
	\item[Attori:] Docente;
	\item[Scopo e descrizione:] Il docente visualizza i risultati e le statistiche relative alla classe selezionata
	\item[Precondizione:] Il docente ha selezionato una classe;
	
	\item[Flusso principale degli eventi:] \ 
	\begin{enumerate}
		\item Il docente seleziona lo studente di cui vuole vedere i risultati e le statistiche;
		\item Il docente visualizza i risultati e le statistiche di un questionario della classe eseguito dallo studente selezionato (\hyperlink{UC24.4.1}{UC24.4.1});
		
	\end{enumerate}
	\item[Postcondizione:] Il docente ha visualizzato i risultati e le statistiche relative allo studente desiderate.
\end{description}
\hypertarget{UC24.4.1}{}
\subsection{Caso d'uso UC24.4.1: Visualizza risultati questionari dello studente}\begin{description}
	\item[Attori:] Docente;
	\item[Scopo e descrizione:] Il docente visualizza i risultati e le statistiche relative ai questionari dello studente selezionato
	\item[Precondizione:] Il docente ha selezionato uno studente della classe;
	
	\item[Flusso principale degli eventi:] \ 
	\begin{enumerate}
		\item Il docente cerca un questionario tra quelli dello studente nella classe (\hyperlink{UC4}{UC4});
		\item Il docente seleziona il questionario per visualizzarne i risultati e le statistiche;
		
	\end{enumerate}
	\item[Postcondizione:] Il docente ha visualizzato i risultati e le statistiche relative ai questionari dello studente selezionato.
\end{description}
\hypertarget{UC25}{}
\subsection{Caso d'uso UC25: Ricerca domanda}
\begin{figure}[H]
	\centering
	\begin{resizedtikzpicture}{\textwidth}
		\umlactor[x=0, y=0]{Docente}
		\begin{umlsystem}[x=0, fill=lightgray!20]{Quizzipedia}
			\umlusecase[x=5, y=0, fill=white, width=4cm, name=23]{\textbf{UC25:} Ricerca domanda}
			\umlassoc{Docente}{23}
			\umlusecase[x=15, y=-6, fill=white, width=4cm, name=27]{\textbf{UC29:} Ricerca domanda per docente}
			\umlinherit{27}{23}
			\umlusecase[x=15, y=-2, fill=white, width=4cm, name=26]{\textbf{UC28:} Ricerca domanda per difficoltà}
			\umlinherit{26}{23}
			\umlusecase[x=15, y=2, fill=white, width=4cm, name=25]{\textbf{UC27:} Ricerca domanda per argomento}
			\umlinherit{25}{23}
			\umlusecase[x=15, y=6, fill=white, width=4cm, name=24]{\textbf{UC26:} Ricerca domanda per keywords}
			\umlinherit{24}{23}
		\end{umlsystem}
	\end{resizedtikzpicture}
	\caption{\textbf{UC25}: Ricerca domanda}
	\label{UC25}
\end{figure}
\begin{description}
	\item[Attori:] Docente;
	\item[Scopo e descrizione:] Il docente ricerca una domanda
	\item[Precondizione:] Il docente è autenticato presso il sistema;
	
	\item[Flusso principale degli eventi:] \ 
	\begin{enumerate}
		\item L'utente inserisce i dati per la ricerca	;
		
	\end{enumerate}
	\item[Postcondizione:] Il sistema mostra la lista delle domande che soddisfano la ricerca.
\end{description}
\hypertarget{UC26}{}
\subsection{Caso d'uso UC26: Ricerca domanda per keywords}\begin{description}
	\item[Attori:] Docente;
	\item[Scopo e descrizione:] Il docente  ricerca una domanda per keywords
	
	\item[Precondizione:] Il docente è autenticato presso il sistema
	;
	
	\item[Flusso principale degli eventi:] \ 
	\begin{enumerate}
		\item L'utente inserisce le keywords presenti nel testo delle domande che vuole cercare;
		
	\end{enumerate}
	\item[Postcondizione:] Il sistema mostra la lista delle domande che contengono nel titolo o nel corpo le parole chiave specificate.
\end{description}
\hypertarget{UC27}{}
\subsection{Caso d'uso UC27: Ricerca domanda per argomento}\begin{description}
	\item[Attori:] Docente;
	\item[Scopo e descrizione:] Il docente ricerca una domanda per argomento
	\item[Precondizione:] Il docente è autenticato presso il sistema
	;
	
	\item[Flusso principale degli eventi:] \ 
	\begin{enumerate}
		\item L'utente seleziona gli argomenti di cui vuole visualizzare le domande;
		
	\end{enumerate}
	\item[Postcondizione:] Il sistema mostra la lista delle domande che contengono gli argomenti selezionati.
\end{description}
\hypertarget{UC28}{}
\subsection{Caso d'uso UC28: Ricerca domanda per difficoltà}\begin{description}
	\item[Attori:] Docente;
	\item[Scopo e descrizione:] Il docente ricerca una domanda per difficoltà
	
	\item[Precondizione:] Il docente è autenticato presso il sistema
	;
	
	\item[Flusso principale degli eventi:] \ 
	\begin{enumerate}
		\item L'utente inserisce il limite minimo e massimo di difficoltà di cui vuole visualizzare le domande;
		
	\end{enumerate}
	\item[Postcondizione:] Il sistema mostra la lista delle domande con la difficoltà selezionata.
\end{description}
\hypertarget{UC29}{}
\subsection{Caso d'uso UC29: Ricerca domanda per docente}\begin{description}
	\item[Attori:] Docente;
	\item[Scopo e descrizione:] Il docente ricerca una domanda per docente
	\item[Precondizione:] Il docente è autenticato presso il sistema
	;
	
	\item[Flusso principale degli eventi:] \ 
	\begin{enumerate}
		\item L'utente inserisce il nome del docente di cui vuole visualizzare le domande;
		
	\end{enumerate}
	\item[Postcondizione:] Il sistema mostra la lista delle domande create dal docente selezionato.
\end{description}
\hypertarget{UC30}{}
\subsection{Caso d'uso UC30: Azioni Amministratore}
\begin{figure}[H]
	\centering
	\begin{resizedtikzpicture}{\textwidth}
		\umlactor[x=0, y=0]{Amministratore}
		\begin{umlsystem}[x=0, fill=lightgray!20]{Quizzipedia}
			\umlusecase[x=5, y=-1.75, fill=white, width=4cm, name=35]{\textbf{UC30.2:} Rimozione utente}
			\umlassoc{Amministratore}{35}
			\umlusecase[x=5, y=2, fill=white, width=4cm, name=34]{\textbf{UC30.1:} Cambia ruolo}
			\umlassoc{Amministratore}{34}
		\end{umlsystem}
	\end{resizedtikzpicture}
	\caption{\textbf{UC30}: Azioni Amministratore}
	\label{UC30}
\end{figure}
\begin{description}
	\item[Attori:] Amministratore;
	\item[Scopo e descrizione:] Descrive le azione che può compiere un aministratore
	\item[Precondizione:] L'amministratore è autenticato nel sistema;
	
	\item[Flusso principale degli eventi:] \ 
	\begin{enumerate}
		\item L'amministratore può impostare il ruolo di un utente con ruolo inferiore al proprio (\hyperlink{UC30.1}{UC30.1});
		\item L'amministratore può rimuovere un utente con ruolo inferiore al proprio (\hyperlink{UC30.2}{UC30.2});
		
	\end{enumerate}
	\item[Postcondizione:] Il sistema ha ottenuto le informazioni sulle operazioni che l’amministratore desidera eseguire.
\end{description}
\hypertarget{UC30.1}{}
\subsection{Caso d'uso UC30.1: Cambia ruolo}\begin{description}
	\item[Attori:] Amministratore;
	\item[Scopo e descrizione:] L'amministratore imposta il ruolo dello studente a docente o viceversa
	\item[Precondizione:] L'utente selezionato è uno studente o docente;
	
	\item[Flusso principale degli eventi:] \ 
	\begin{enumerate}
		\item L'amministratore seleziona un utente con ruolo inferiore al proprio;
		\item L'amministratore imposta il ruolo dell'utente selezionato ad uno inferiore al proprio;
		\item L'amministratore conferma l'operazione da eseguire;
		
	\end{enumerate}
	\item[Postcondizione:] Il ruolo dell'utente è diventato quello selezionato.
\end{description}
\hypertarget{UC30.2}{}
\subsection{Caso d'uso UC30.2: Rimozione utente}\begin{description}
	\item[Attori:] Amministratore;
	\item[Scopo e descrizione:] L'amministratore rimuove un utente con ruolo inferiore al proprio
	\item[Precondizione:] L'utente selezionato esiste nel sistema ed ha un ruolo inferiore al proprio;
	
	\item[Flusso principale degli eventi:] \ 
	\begin{enumerate}
		\item Amministratore ricerca l'utente con ruolo inferiore al proprio;
		\item Amministratore seleziona l'utente con ruolo inferiore al proprio;
		\item Amministratore conferma la rimozione;
		
	\end{enumerate}
	\item[Postcondizione:] L'utente selezionato è stato rimosso dal sistema.
\end{description}
\hypertarget{UC31}{}
\subsection{Caso d'uso UC31: Ricerca utente}
\begin{figure}[H]
	\centering
	\begin{resizedtikzpicture}{\textwidth}
		\umlactor[x=0, y=0]{Amministratore}
		\begin{umlsystem}[x=0, fill=lightgray!20]{Quizzipedia}
			\umlusecase[x=5, y=0, fill=white, width=4cm, name=200]{\textbf{UC31:} Ricerca utente}
			\umlassoc{Amministratore}{200}
			\umlusecase[x=15, y=-2, fill=white, width=4cm, name=203]{\textbf{UC33:} Ricerca utente per username}
			\umlinherit{203}{200}
			\umlusecase[x=15, y=2, fill=white, width=4cm, name=201]{\textbf{UC32:} Ricerca utente per nome completo}
			\umlinherit{201}{200}
		\end{umlsystem}
	\end{resizedtikzpicture}
	\caption{\textbf{UC31}: Ricerca utente}
	\label{UC31}
\end{figure}
\begin{description}
	\item[Attori:] Amministratore;
	\item[Scopo e descrizione:] L'amministratore ricerca un utente
	\item[Precondizione:] L'amministratore è autenticato presso il sistema;
	
	\item[Flusso principale degli eventi:] \ 
	\begin{enumerate}
		\item L'amministratore inserisce i dati per la ricerca;
		
	\end{enumerate}
	\item[Postcondizione:] Il sistema mostra la lista degli utenti che soddisfano i criteri di ricerca.
\end{description}
\hypertarget{UC32}{}
\subsection{Caso d'uso UC32: Ricerca utente per nome completo}
\begin{figure}[H]
	\centering
	\begin{resizedtikzpicture}{\textwidth}
		\umlactor[x=0, y=0]{Amministratore}
		\begin{umlsystem}[x=0, fill=lightgray!20]{Quizzipedia}
			\umlusecase[x=5, y=0, fill=white, width=4cm, name=201]{\textbf{UC32:} Ricerca utente per nome completo}
			\umlassoc{Amministratore}{201}
			\umlusecase[x=15, y=0, fill=white, width=4cm, name=204]{\textbf{UC34:} Ricerca utente per ruolo}
			\umlinherit{204}{201}
		\end{umlsystem}
	\end{resizedtikzpicture}
	\caption{\textbf{UC32}: Ricerca utente per nome completo}
	\label{UC32}
\end{figure}
\begin{description}
	\item[Attori:] Amministratore;
	\item[Scopo e descrizione:] L'amministratore ricerca un utente per nome completo
	\item[Precondizione:] L'amministratore è autenticato presso il sistema;
	
	\item[Flusso principale degli eventi:] \ 
	\begin{enumerate}
		\item L'amministratore ricerca un utente per nome completo;
		
	\end{enumerate}
	\item[Postcondizione:] Il sistema mostra la lista degli utenti il cui nome completo corrisponde al nome completo ricercato.
\end{description}
\hypertarget{UC33}{}
\subsection{Caso d'uso UC33: Ricerca utente per username}\begin{description}
	\item[Attori:] Amministratore;
	\item[Scopo e descrizione:] L'amministratore ricerca un utente per username
	\item[Precondizione:] L'amministratore è autenticato presso il sistema;
	
	\item[Flusso principale degli eventi:] \ 
	\begin{enumerate}
		\item L'amministratore ricerca un utente per username;
		
	\end{enumerate}
	\item[Postcondizione:] Il sistema mostra la lista degli utenti il cui username corrisponde all'username ricercato.
\end{description}
\hypertarget{UC34}{}
\subsection{Caso d'uso UC34: Ricerca utente per ruolo}\begin{description}
	\item[Attori:] Amministratore;
	\item[Scopo e descrizione:] L'amministratore ricerca un utente in base al ruolo 
	
	\item[Precondizione:] L'amministratore è autenticato presso il sistema
	;
	
	\item[Flusso principale degli eventi:] \ 
	\begin{enumerate}
		\item L'amministratore ricerca un utente in base al ruolo 
		;
		
	\end{enumerate}
	\item[Postcondizione:] Il sistema mostra la lista degli utenti il cui ruolo completo corrisponde al nome completo ricercato
	.
\end{description}


\newpage
\section{Requisiti}
\subsection{Struttura di un requisito}
È compito degli \ANpl\ stilare una lista dei requisiti emersi dal capitolato e da eventuali riunioni con il proponente. Questi requisiti dovranno essere classificati per tipo e importanza, utilizzando la seguente struttura:
\begin{center}
	R[importanza][tipo][codice]
\end{center}
\begin{itemize}
	\item \textbf{Importanza} può assumere i seguenti valori:
	\begin{description}
		\item[1:] requisito desiderabile
		\item[2:] requisito opzionale
		\item[3:] requisito obbligatorio
	\end{description}
	\item \textbf{Tipo} può assumere i seguenti valori:
	\begin{description}
		\item[F:] Funzionale, descrive i servizi o le funzioni offerte dal sistema
		\item[Q:] Di Qualità, descrive i requisiti sulla qualità offerte dal sistema
		\item[P:] Prestazionale, descrive i requisiti sulle prestazioni offerte dal sistema
		\item[V:] Vincolo, descrive i vincoli sui servizi offerti dal sistema
	\end{description}
	\item \textbf{Codice} è un numero progressivo univoco per ogni requisito, indipendente da importanza e tipo. Nel caso si abbia un sotto-requisito codice può anche essere espresso in modo gerarchico tramite la notazione:
	\begin{center}
		\textit{CodiceRequistoPadre.CodiceSottorequisito}
	\end{center}
	\item \textbf{Validazione:} verrà inserito un link al metodo deciso per validare il requisito
\end{itemize}
Ogni requisito deve essere correlato da una sintetica ma precisa descrizione. Per ogni requisito bisogna indicarne le fonti, che posso essere il capitolato o uno o più casi d'uso.

\subsection{Tabella dei requisiti}

% Output di tracy sotto a:
% tracy-tpiga.rhcloud.com/tracy/latex/RequirementSource?type=latex&raw=1
\begin{longtable}{|r l|p{2cm}|p{6cm}|p{2cm}|}
\hline
\multicolumn{2}{|c|}{Requisito} & Tipologia & Descrizione & Fonti\tabularnewline
\hline
 & \hypertarget{R-3V1}{R-3V1} & Vincolo

Obbligatorio & La parte destinata ai creatori di domande e di questionari dovrà essere fruibile attraverso un PC & Capitolato\tabularnewline
\hline
 & \hypertarget{R-3V2}{R-3V2} & Vincolo

Obbligatorio & La parte server deve essere realizzata in Javascript con server node.js
 & Capitolato\tabularnewline
\hline
 & \hypertarget{R-3V3}{R-3V3} & Vincolo

Obbligatorio & La parte client dovrà essere eseguibile in un browser HTML5 & Capitolato\tabularnewline
\hline
\begin{tikzpicture}
\draw [->, thick] (0.2,0.2) -- (0.2,0.1) -- (1,0.1);
\end{tikzpicture} & \hypertarget{R-3V3.1}{R-3V3.1} & Vincolo

Obbligatorio & L'applicazione funziona correttamente su Chrome versione 47.0.2526 e successive & Capitolato\tabularnewline
\hline
\begin{tikzpicture}
\draw [->, thick] (0.2,0.2) -- (0.2,0.1) -- (1,0.1);
\end{tikzpicture} & \hypertarget{R-3V3.2}{R-3V3.2} & Vincolo

Obbligatorio & L'applicazione funziona correttamente su Firefox versione 38.5.2esr e successive & Capitolato\tabularnewline
\hline
\begin{tikzpicture}
\draw [->, thick] (0.2,0.2) -- (0.2,0.1) -- (1,0.1);
\end{tikzpicture} & \hypertarget{R-3V3.3}{R-3V3.3} & Vincolo

Obbligatorio & L'applicazione funziona correttamente su Microsoft Edge versione 25.10586 e successive & Capitolato\tabularnewline
\hline
\begin{tikzpicture}
\draw [->, thick] (0.2,0.2) -- (0.2,0.1) -- (1,0.1);
\end{tikzpicture} & \hypertarget{R-3V3.4}{R-3V3.4} & Vincolo

Obbligatorio & L'applicazione funziona correttamente su Safari versione 9.0.2 e successive & Capitolato\tabularnewline
\hline
 & \hypertarget{R-3V4}{R-3V4} & Vincolo

Obbligatorio & L'applicazione deve utilizzare fogli stile CSS3 per l’ aspetto estetico & Capitolato\tabularnewline
\hline
 & \hypertarget{R-3V5}{R-3V5} & Vincolo

Obbligatorio & L'applicazione deve utilizzare Javascript per la parte attiva del client & Capitolato\tabularnewline
\hline
 & \hypertarget{R-3V6}{R-3V6} & Vincolo

Obbligatorio & La parte destinata agli esaminandi dovrà funzionare con qualunque dispositivo: dal Personal Computer, ai Tablet ed agli SmartPhone attraverso browser & Capitolato\tabularnewline
\hline
\begin{tikzpicture}
\draw [->, thick] (0.2,0.2) -- (0.2,0.1) -- (1,0.1);
\end{tikzpicture} & \hypertarget{R-3V6.1}{R-3V6.1} & Vincolo

Obbligatorio & La parte destinata agli esaminandi dovrà funzionare si smartphone e tablet iOS dalla versione iOS6 & Capitolato\tabularnewline
\hline
\begin{tikzpicture}
\draw [->, thick] (0.2,0.2) -- (0.2,0.1) -- (1,0.1);
\end{tikzpicture} & \hypertarget{R-3V6.2}{R-3V6.2} & Vincolo

Obbligatorio & La parte destinata agli esaminandi dovrà funzionare su Android dalla versione 4.0 & Capitolato\tabularnewline
\hline
 & \hypertarget{R-3F7}{R-3F7} & Funzionale

Obbligatorio & Il sistema deve gestire questionari & Capitolato

\hyperlink{UC8}{UC8}\tabularnewline
\hline
\begin{tikzpicture}
\draw [->, thick] (0.2,0.2) -- (0.2,0.1) -- (1,0.1);
\end{tikzpicture} & \hypertarget{R-3F7.1}{R-3F7.1} & Funzionale

Obbligatorio & L'applicazione deve archiviare questionari in un server, suddivisi per argomento
 & Capitolato\tabularnewline
\hline
\begin{tikzpicture}
\draw [->, thick] (0.2,0.2) -- (0.2,0.1) -- (1,0.1);
\end{tikzpicture} & \hypertarget{R-3F7.2}{R-3F7.2} & Funzionale

Obbligatorio & L'applicazione deve tradurre da QML a HTML le domande archiviate & Capitolato\tabularnewline
\hline
\begin{tikzpicture}
\draw [->, thick] (0.2,0.2) -- (0.2,0.1) -- (1,0.1);
\end{tikzpicture} & \hypertarget{R-3F7.3}{R-3F7.3} & Funzionale

Obbligatorio & L'applicazione deve proporre questionari preconfezionati & Capitolato\tabularnewline
\hline
\begin{tikzpicture}
\draw [->, thick] (0.2,0.2) -- (0.2,0.1) -- (1,0.1);
\end{tikzpicture} & \hypertarget{R-3F7.4}{R-3F7.4} & Funzionale

Obbligatorio & L'applicazione deve valutare le risposte fornite dall'utente & Capitolato

\hyperlink{UC17}{UC17}\tabularnewline
\hline
\begin{tikzpicture}
\draw [->, thick] (0.2,0.2) -- (0.2,0.1) -- (1,0.1);
\end{tikzpicture} & \hypertarget{R-3F7.5}{R-3F7.5} & Funzionale

Obbligatorio & Le domande dovranno essere raccolte attraverso uno specifico linguaggio, che chiameremo Quiz Markup Language (QML), studiato per descrivere quiz
 & Capitolato

\hyperlink{UC7.1.2}{UC7.1.2}\tabularnewline
\hline
\begin{tikzpicture}
\draw [->, thick] (0.4,0.2) -- (0.4,0.1) -- (1,0.1);
\end{tikzpicture} & \hypertarget{R-3F7.5.1}{R-3F7.5.1} & Funzionale

Obbligatorio & Il QML deve essere un linguaggio semplice con un parsing ad hoc definito dal gruppo. La sintassi per scrivere il testo della domanda deve essere unica, la sintassi per le risposte è dipendente dal tipo di domanda & Capitolato

\hyperlink{UC7.1.2}{UC7.1.2}\tabularnewline
\hline
\begin{tikzpicture}
\draw [->, thick] (0.6,0.2) -- (0.6,0.1) -- (1,0.1);
\end{tikzpicture} & \hypertarget{R-2F7.5.1.1}{R-2F7.5.1.1} & Funzionale

Opzionale & Il QML deve permettere l'inserimento di testo in grassetto & \hyperlink{UC7.1.2}{UC7.1.2}\tabularnewline
\hline
\begin{tikzpicture}
\draw [->, thick] (0.6,0.2) -- (0.6,0.1) -- (1,0.1);
\end{tikzpicture} & \hypertarget{R-3F7.5.1.2}{R-3F7.5.1.2} & Funzionale

Obbligatorio & Il QML deve permettere l'inserimento di immagini & Capitolato

\hyperlink{UC7.1.2}{UC7.1.2}\tabularnewline
\hline
\begin{tikzpicture}
\draw [->, thick] (0.6,0.2) -- (0.6,0.1) -- (1,0.1);
\end{tikzpicture} & \hypertarget{R-2F7.5.1.3}{R-2F7.5.1.3} & Funzionale

Opzionale & Il QML deve permettere l'inserimento di testo in corsivo & \hyperlink{UC7.1.2}{UC7.1.2}\tabularnewline
\hline
\begin{tikzpicture}
\draw [->, thick] (0.6,0.2) -- (0.6,0.1) -- (1,0.1);
\end{tikzpicture} & \hypertarget{R-2F7.5.1.4}{R-2F7.5.1.4} & Funzionale

Opzionale & Il QML deve permettere l'inserimento di tabelle & \hyperlink{UC7.1.2}{UC7.1.2}\tabularnewline
\hline
\begin{tikzpicture}
\draw [->, thick] (0.6,0.2) -- (0.6,0.1) -- (1,0.1);
\end{tikzpicture} & \hypertarget{R-3F7.5.1.5}{R-3F7.5.1.5} & Funzionale

Obbligatorio & Il QML deve permettere di definire un punteggio per ogni risposta & \hyperlink{UC7.1.2}{UC7.1.2}\tabularnewline
\hline
\begin{tikzpicture}
\draw [->, thick] (0.6,0.2) -- (0.6,0.1) -- (1,0.1);
\end{tikzpicture} & \hypertarget{R-3F7.5.1.6}{R-3F7.5.1.6} & Funzionale

Obbligatorio & Il QML deve permettere di definire il tipo di domanda & \hyperlink{UC7.1.2}{UC7.1.2}\tabularnewline
\hline
\begin{tikzpicture}
\draw [->, thick] (0.6,0.2) -- (0.6,0.1) -- (1,0.1);
\end{tikzpicture} & \hypertarget{R-2F7.5.1.7}{R-2F7.5.1.7} & Funzionale

Opzionale & Il QML deve permettere l'inserimento di video & \hyperlink{UC7.1.2}{UC7.1.2}

\hyperlink{UC7.1.3}{UC7.1.3}\tabularnewline
\hline
\begin{tikzpicture}
\draw [->, thick] (0.6,0.2) -- (0.6,0.1) -- (1,0.1);
\end{tikzpicture} & \hypertarget{R-2F7.5.1.8}{R-2F7.5.1.8} & Funzionale

Opzionale & Il QML deve permettere l'inserimento di audio & \hyperlink{UC7.1.2}{UC7.1.2}

\hyperlink{UC7.1.3}{UC7.1.3}\tabularnewline
\hline
\begin{tikzpicture}
\draw [->, thick] (0.4,0.2) -- (0.4,0.1) -- (1,0.1);
\end{tikzpicture} & \hypertarget{R-2F7.5.2}{R-2F7.5.2} & Funzionale

Opzionale & Il QML può gestire risposte a scelta multipla, a riempimento di parole omesse, ad associazione, a risposta aperta, con testi e immagini & Capitolato

\hyperlink{UC7.10}{UC7.10}

\hyperlink{UC7.11}{UC7.11}

\hyperlink{UC7.1}{UC7.1}

\hyperlink{UC7.8}{UC7.8}

\hyperlink{UC7.9}{UC7.9}\tabularnewline
\hline
\begin{tikzpicture}
\draw [->, thick] (0.4,0.2) -- (0.4,0.1) -- (1,0.1);
\end{tikzpicture} & \hypertarget{R-3F7.5.3}{R-3F7.5.3} & Funzionale

Obbligatorio & Le domande scritte in QML devono essere editabili dall'utente & Proponente

\hyperlink{UC7.1}{UC7.1}\tabularnewline
\hline
\begin{tikzpicture}
\draw [->, thick] (0.4,0.2) -- (0.4,0.1) -- (1,0.1);
\end{tikzpicture} & \hypertarget{R-2F7.5.4}{R-2F7.5.4} & Funzionale

Opzionale & Le domande possono essere inserite e editabili attraverso interfaccia grafica & Proponente

\hyperlink{UC7.1.3}{UC7.1.3}\tabularnewline
\hline
\begin{tikzpicture}
\draw [->, thick] (0.4,0.2) -- (0.4,0.1) -- (1,0.1);
\end{tikzpicture} & \hypertarget{R-3F7.5.5}{R-3F7.5.5} & Funzionale

Obbligatorio & Il QML deve gestire risposte vero/falso e risposte a scelta multipla con testi e immagini & Capitolato

\hyperlink{UC7.1}{UC7.1}

\hyperlink{UC7.6}{UC7.6}

\hyperlink{UC7.8}{UC7.8}\tabularnewline
\hline
\begin{tikzpicture}
\draw [->, thick] (0.2,0.2) -- (0.2,0.1) -- (1,0.1);
\end{tikzpicture} & \hypertarget{R-2F7.6}{R-2F7.6} & Funzionale

Opzionale & Il sistema deve gestire sia questionari pubblici che privati attraverso la creazione di classi la cui iscrizione è protetta da password & Proponente

\hyperlink{UC16}{UC16}

\hyperlink{UC9.1}{UC9.1}\tabularnewline
\hline
\begin{tikzpicture}
\draw [->, thick] (0.2,0.2) -- (0.2,0.1) -- (1,0.1);
\end{tikzpicture} & \hypertarget{R-3F7.7}{R-3F7.7} & Funzionale

Obbligatorio & I docenti costruiranno questionari scegliendo tra le domande archiviate nel server & Capitolato

\hyperlink{UC8.1}{UC8.1}

\hyperlink{UC8}{UC8}\tabularnewline
\hline
\begin{tikzpicture}
\draw [->, thick] (0.4,0.2) -- (0.4,0.1) -- (1,0.1);
\end{tikzpicture} & \hypertarget{R-3F7.7.1}{R-3F7.7.1} & Funzionale

Obbligatorio & Il sistema deve segnalare un errore nel caso un questionario creato non contenga domande & \hyperlink{UC8.1}{UC8.1}

\hyperlink{UC8.4}{UC8.4}

\hyperlink{UC8}{UC8}\tabularnewline
\hline
\begin{tikzpicture}
\draw [->, thick] (0.2,0.2) -- (0.2,0.1) -- (1,0.1);
\end{tikzpicture} & \hypertarget{R-2F7.8}{R-2F7.8} & Funzionale

Opzionale & Il sistema deve offrire la possibilità di creare dinamicamente i questionari & Capitolato\tabularnewline
\hline
\begin{tikzpicture}
\draw [->, thick] (0.2,0.2) -- (0.2,0.1) -- (1,0.1);
\end{tikzpicture} & \hypertarget{R-2F7.9}{R-2F7.9} & Funzionale

Opzionale & L'applicazione deve essere predisposta ad un sistema di feedback per i questionari & Capitolato

\hyperlink{UC15.10}{UC15.10}\tabularnewline
\hline
\begin{tikzpicture}
\draw [->, thick] (0.2,0.2) -- (0.2,0.1) -- (1,0.1);
\end{tikzpicture} & \hypertarget{R-3F7.10}{R-3F7.10} & Funzionale

Obbligatorio & Il sistema deve permettere la raccolta dei dati relativi a questionari svolti & Capitolato

\hyperlink{UC17}{UC17}\tabularnewline
\hline
\begin{tikzpicture}
\draw [->, thick] (0.4,0.2) -- (0.4,0.1) -- (1,0.1);
\end{tikzpicture} & \hypertarget{R-3F7.10.1}{R-3F7.10.1} & Funzionale

Obbligatorio & Il sistema deve permettere la rivisitazione del test svolto e dei risultati ottenuti & Capitolato

\hyperlink{UC17}{UC17}\tabularnewline
\hline
\begin{tikzpicture}
\draw [->, thick] (0.4,0.2) -- (0.4,0.1) -- (1,0.1);
\end{tikzpicture} & \hypertarget{R-3F7.10.2}{R-3F7.10.2} & Funzionale

Obbligatorio & Il sistema deve offrire la possibilità di analizzare i dati raccolti dai test effettuati & Capitolato

\hyperlink{UC17.2}{UC17.2}

\hyperlink{UC17}{UC17}\tabularnewline
\hline
\begin{tikzpicture}
\draw [->, thick] (0.4,0.2) -- (0.4,0.1) -- (1,0.1);
\end{tikzpicture} & \hypertarget{R-2F7.10.3}{R-2F7.10.3} & Funzionale

Opzionale & Il sistema deve offrire la possibilità di visualizzazione sotto forma di grafico delle statistiche dei dati raccolti & Capitolato

\hyperlink{UC17.3}{UC17.3}

\hyperlink{UC17}{UC17}\tabularnewline
\hline
\begin{tikzpicture}
\draw [->, thick] (0.2,0.2) -- (0.2,0.1) -- (1,0.1);
\end{tikzpicture} & \hypertarget{R-3F7.11}{R-3F7.11} & Funzionale

Obbligatorio & Il sistema deve prevedere la gestione di domande & \hyperlink{UC7}{UC7}

\hyperlink{UC8.1.1}{UC8.1.1}

\hyperlink{UC8.1.2}{UC8.1.2}

\hyperlink{UC8}{UC8}\tabularnewline
\hline
\begin{tikzpicture}
\draw [->, thick] (0.4,0.2) -- (0.4,0.1) -- (1,0.1);
\end{tikzpicture} & \hypertarget{R-3F7.11.1}{R-3F7.11.1} & Funzionale

Obbligatorio & Il sistema deve consentire ad un docente di inserire una nuova domanda & \hyperlink{UC7.1}{UC7.1}

\hyperlink{UC7}{UC7}\tabularnewline
\hline
\begin{tikzpicture}
\draw [->, thick] (0.6,0.2) -- (0.6,0.1) -- (1,0.1);
\end{tikzpicture} & \hypertarget{R-3F7.11.1.1}{R-3F7.11.1.1} & Funzionale

Obbligatorio & Un docente deve poter selezionare gli argomenti di una domanda & \hyperlink{UC7.1.1}{UC7.1.1}

\hyperlink{UC7.1}{UC7.1}

\hyperlink{UC7}{UC7}\tabularnewline
\hline
\begin{tikzpicture}
\draw [->, thick] (0.8,0.2) -- (0.8,0.1) -- (1,0.1);
\end{tikzpicture} & \hypertarget{R-3F7.11.1.1.1}{R-3F7.11.1.1.1} & Funzionale

Obbligatorio & Il sistema deve segnalare un errore nel caso non venga selezionato alcun argomento per una nuova domanda & \hyperlink{UC7.1.1}{UC7.1.1}

\hyperlink{UC7.5}{UC7.5}

\hyperlink{UC7}{UC7}

\hyperlink{UC8.1.1}{UC8.1.1}\tabularnewline
\hline
\begin{tikzpicture}
\draw [->, thick] (0.6,0.2) -- (0.6,0.1) -- (1,0.1);
\end{tikzpicture} & \hypertarget{R-3F7.11.1.2}{R-3F7.11.1.2} & Funzionale

Obbligatorio & Un docente deve comporre la domanda da inserire in QML & \hyperlink{UC7.1.2}{UC7.1.2}

\hyperlink{UC7}{UC7}

\hyperlink{UC8.1.1}{UC8.1.1}\tabularnewline
\hline
\begin{tikzpicture}
\draw [->, thick] (0.8,0.2) -- (0.8,0.1) -- (1,0.1);
\end{tikzpicture} & \hypertarget{R-3F7.11.1.2.1}{R-3F7.11.1.2.1} & Funzionale

Obbligatorio & Il sistema deve segnalare un errore nel caso di inserimento QML non valido & \hyperlink{UC7.1.2}{UC7.1.2}

\hyperlink{UC7.4}{UC7.4}

\hyperlink{UC7}{UC7}

\hyperlink{UC8.1.1}{UC8.1.1}\tabularnewline
\hline
\begin{tikzpicture}
\draw [->, thick] (0.8,0.2) -- (0.8,0.1) -- (1,0.1);
\end{tikzpicture} & \hypertarget{R-3F7.11.1.2.2}{R-3F7.11.1.2.2} & Funzionale

Obbligatorio & Il docente deve poter inserire domande di tipo vero/falso & \hyperlink{UC7.1}{UC7.1}

\hyperlink{UC7.6}{UC7.6}\tabularnewline
\hline
\begin{tikzpicture}
\draw [->, thick] (0.8,0.2) -- (0.8,0.1) -- (1,0.1);
\end{tikzpicture} & \hypertarget{R-3F7.11.1.2.3}{R-3F7.11.1.2.3} & Funzionale

Obbligatorio & Il docente deve poter inserire domande a scelta multipla & \hyperlink{UC7.1}{UC7.1}

\hyperlink{UC7.7}{UC7.7}\tabularnewline
\hline
\begin{tikzpicture}
\draw [->, thick] (0.8,0.2) -- (0.8,0.1) -- (1,0.1);
\end{tikzpicture} & \hypertarget{R-3F7.11.1.2.4}{R-3F7.11.1.2.4} & Funzionale

Obbligatorio & Il docente deve poter inserire domande a risposta multipla & \hyperlink{UC7.1}{UC7.1}

\hyperlink{UC7.8}{UC7.8}\tabularnewline
\hline
\begin{tikzpicture}
\draw [->, thick] (0.8,0.2) -- (0.8,0.1) -- (1,0.1);
\end{tikzpicture} & \hypertarget{R-2F7.11.1.2.5}{R-2F7.11.1.2.5} & Funzionale

Opzionale & Il docente deve poter inserire domande di tipo testo con parole omesse & \hyperlink{UC7.1}{UC7.1}

\hyperlink{UC7.9}{UC7.9}\tabularnewline
\hline
\begin{tikzpicture}
\draw [->, thick] (0.8,0.2) -- (0.8,0.1) -- (1,0.1);
\end{tikzpicture} & \hypertarget{R-2F7.11.1.2.6}{R-2F7.11.1.2.6} & Funzionale

Opzionale & Il docente deve poter inserire domande con associazione di parole & \hyperlink{UC7.10}{UC7.10}

\hyperlink{UC7.1}{UC7.1}\tabularnewline
\hline
\begin{tikzpicture}
\draw [->, thick] (0.8,0.2) -- (0.8,0.1) -- (1,0.1);
\end{tikzpicture} & \hypertarget{R-2F7.11.1.2.7}{R-2F7.11.1.2.7} & Funzionale

Opzionale & Il docente deve poter inserire domande a risposta aperta & \hyperlink{UC7.11}{UC7.11}

\hyperlink{UC7.1}{UC7.1}\tabularnewline
\hline
\begin{tikzpicture}
\draw [->, thick] (0.4,0.2) -- (0.4,0.1) -- (1,0.1);
\end{tikzpicture} & \hypertarget{R-3F7.11.2}{R-3F7.11.2} & Funzionale

Obbligatorio & Il sistema deve consentire ad un docente di modificare una domanda & \hyperlink{UC7.2}{UC7.2}

\hyperlink{UC7}{UC7}\tabularnewline
\hline
\begin{tikzpicture}
\draw [->, thick] (0.4,0.2) -- (0.4,0.1) -- (1,0.1);
\end{tikzpicture} & \hypertarget{R-3F7.11.3}{R-3F7.11.3} & Funzionale

Obbligatorio & Il sistema deve consentire ad un docente di eliminare una domanda presente nel sistema & \hyperlink{UC7.3}{UC7.3}

\hyperlink{UC7}{UC7}\tabularnewline
\hline
\begin{tikzpicture}
\draw [->, thick] (0.2,0.2) -- (0.2,0.1) -- (1,0.1);
\end{tikzpicture} & \hypertarget{R-3F7.12}{R-3F7.12} & Funzionale

Obbligatorio & il sistema deve dare la possibilità al docente di modificare il proprio questionario & \hyperlink{UC8.2}{UC8.2}

\hyperlink{UC8}{UC8}\tabularnewline
\hline
\begin{tikzpicture}
\draw [->, thick] (0.4,0.2) -- (0.4,0.1) -- (1,0.1);
\end{tikzpicture} & \hypertarget{R-3F7.12.1}{R-3F7.12.1} & Funzionale

Obbligatorio & Il sistema deve dare la possibilità di ricercare un questionario & \hyperlink{UC18}{UC18}\tabularnewline
\hline
\begin{tikzpicture}
\draw [->, thick] (0.4,0.2) -- (0.4,0.1) -- (1,0.1);
\end{tikzpicture} & \hypertarget{R-3F7.12.2}{R-3F7.12.2} & Funzionale

Obbligatorio & Il sistema deve dare la possibile al docente di aggiungere al questionario già fatto una domanda. & \hyperlink{UC8.1.1}{UC8.1.1}

\hyperlink{UC8}{UC8}\tabularnewline
\hline
\begin{tikzpicture}
\draw [->, thick] (0.4,0.2) -- (0.4,0.1) -- (1,0.1);
\end{tikzpicture} & \hypertarget{R-3F7.12.3}{R-3F7.12.3} & Funzionale

Obbligatorio & Il sistema deve dare la possibilità al docente di togliere domande precedentemente aggiunte al questionario & \hyperlink{UC8.1.2}{UC8.1.2}

\hyperlink{UC8}{UC8}\tabularnewline
\hline
\begin{tikzpicture}
\draw [->, thick] (0.4,0.2) -- (0.4,0.1) -- (1,0.1);
\end{tikzpicture} & \hypertarget{R-3F7.12.4}{R-3F7.12.4} & Funzionale

Obbligatorio & Il sistema deve consentire ad un docente di modificare gli argomenti di un proprio questionario & \hyperlink{UC8.1.3}{UC8.1.3}

\hyperlink{UC8.2}{UC8.2}

\hyperlink{UC8}{UC8}\tabularnewline
\hline
\begin{tikzpicture}
\draw [->, thick] (0.2,0.2) -- (0.2,0.1) -- (1,0.1);
\end{tikzpicture} & \hypertarget{R-3F7.13}{R-3F7.13} & Funzionale

Obbligatorio & Il sistema deve dare la possibilità al docente di eliminare il proprio questionario & \hyperlink{UC8.3}{UC8.3}

\hyperlink{UC8}{UC8}\tabularnewline
\hline
 & \hypertarget{R-3F8}{R-3F8} & Funzionale

Obbligatorio & Il sistema deve permettere l’autenticazione
di un utente & \hyperlink{UC1}{UC1}\tabularnewline
\hline
\begin{tikzpicture}
\draw [->, thick] (0.2,0.2) -- (0.2,0.1) -- (1,0.1);
\end{tikzpicture} & \hypertarget{R-3F8.1}{R-3F8.1} & Funzionale

Obbligatorio & L’autenticazione richiede l’inserimento dell’e-mail e della password & 



\hyperlink{UC1}{UC1}\tabularnewline
\hline
\begin{tikzpicture}
\draw [->, thick] (0.2,0.2) -- (0.2,0.1) -- (1,0.1);
\end{tikzpicture} & \hypertarget{R-3F8.2}{R-3F8.2} & Funzionale

Obbligatorio & Il sistema deve permettere la visualizzazione di un messaggio di errore in caso di email o password non valida & 



\hyperlink{UC1}{UC1}\tabularnewline
\hline
 & \hypertarget{R-3F9}{R-3F9} & Funzionale

Obbligatorio & Il candidato attraverso una specifica interfaccia potrà rispondere alle domande & Capitolato

\hyperlink{UC15.1}{UC15.1}\tabularnewline
\hline
 & \hypertarget{R-3F10}{R-3F10} & Funzionale

Obbligatorio & Il sistema deve avere almeno 3 tipi di utenti: amministratore, docente e studente & Capitolato\tabularnewline
\hline
 & \hypertarget{R-3F11}{R-3F11} & Funzionale

Obbligatorio & Il sistema deve gestire la registrazione di un utente & \hyperlink{UC1.1}{UC1.1}\tabularnewline
\hline
\begin{tikzpicture}
\draw [->, thick] (0.2,0.2) -- (0.2,0.1) -- (1,0.1);
\end{tikzpicture} & \hypertarget{R-3F11.1}{R-3F11.1} & Funzionale

Obbligatorio & Il sistema deve contenere già un account proprietario & Interno\tabularnewline
\hline
\begin{tikzpicture}
\draw [->, thick] (0.2,0.2) -- (0.2,0.1) -- (1,0.1);
\end{tikzpicture} & \hypertarget{R-3F11.2}{R-3F11.2} & Funzionale

Obbligatorio & Il sistema deve permettere la registrazione
di uno studente & \hyperlink{UC1.1}{UC1.1}\tabularnewline
\hline
\begin{tikzpicture}
\draw [->, thick] (0.4,0.2) -- (0.4,0.1) -- (1,0.1);
\end{tikzpicture} & \hypertarget{R-3F11.2.1}{R-3F11.2.1} & Funzionale

Obbligatorio & La registrazione di uno studente richiede l’inserimento
della mail, password, nome e cognome & 







\hyperlink{UC1.1}{UC1.1}\tabularnewline
\hline
\begin{tikzpicture}
\draw [->, thick] (0.6,0.2) -- (0.6,0.1) -- (1,0.1);
\end{tikzpicture} & \hypertarget{R-3F11.2.1.1}{R-3F11.2.1.1} & Funzionale

Obbligatorio & Il sistema deve segnalare un errore nel caso non vengano inseriti tutti i campi necessari per la registrazione & 

\hyperlink{UC1.1}{UC1.1}\tabularnewline
\hline
\begin{tikzpicture}
\draw [->, thick] (0.6,0.2) -- (0.6,0.1) -- (1,0.1);
\end{tikzpicture} & \hypertarget{R-3F11.2.1.2}{R-3F11.2.1.2} & Funzionale

Obbligatorio & La registrazione di uno studente deve richiedere l'inserimento delle informazioni personali & 

\hyperlink{UC1.1}{UC1.1}\tabularnewline
\hline
\begin{tikzpicture}
\draw [->, thick] (0.8,0.2) -- (0.8,0.1) -- (1,0.1);
\end{tikzpicture} & \hypertarget{R-3F11.2.1.2.1}{R-3F11.2.1.2.1} & Funzionale

Obbligatorio & Il sistema deve segnalare un errore se non sono state inserite tutte le informazioni personali richieste & 



\hyperlink{UC1.1}{UC1.1}\tabularnewline
\hline
\begin{tikzpicture}
\draw [->, thick] (0.6,0.2) -- (0.6,0.1) -- (1,0.1);
\end{tikzpicture} & \hypertarget{R-3F11.2.1.3}{R-3F11.2.1.3} & Funzionale

Obbligatorio & Il sistema deve richiedere di ripetere l'inserimento della password durante la registrazione di uno studente & 

\hyperlink{UC1.1}{UC1.1}\tabularnewline
\hline
\begin{tikzpicture}
\draw [->, thick] (0.4,0.2) -- (0.4,0.1) -- (1,0.1);
\end{tikzpicture} & \hypertarget{R-3F11.2.2}{R-3F11.2.2} & Funzionale

Obbligatorio & Il sistema deve permettere la visualizzazione di un messaggio di errore se la mail dello studente è già presente nel sistema & 

\hyperlink{UC1.1}{UC1.1}\tabularnewline
\hline
\begin{tikzpicture}
\draw [->, thick] (0.4,0.2) -- (0.4,0.1) -- (1,0.1);
\end{tikzpicture} & \hypertarget{R-3F11.2.3}{R-3F11.2.3} & Funzionale

Obbligatorio & Il sistema deve permettere la visualizzazione di un messaggio di errore se la password dello studente è troppo corta & 

\hyperlink{UC1.1}{UC1.1}\tabularnewline
\hline
\begin{tikzpicture}
\draw [->, thick] (0.4,0.2) -- (0.4,0.1) -- (1,0.1);
\end{tikzpicture} & \hypertarget{R-3F11.2.4}{R-3F11.2.4} & Funzionale

Obbligatorio & Il sistema deve segnalare un errore se non viene inserito un indirizzo email valido & 



\hyperlink{UC1.1}{UC1.1}\tabularnewline
\hline
\begin{tikzpicture}
\draw [->, thick] (0.4,0.2) -- (0.4,0.1) -- (1,0.1);
\end{tikzpicture} & \hypertarget{R-3F11.2.5}{R-3F11.2.5} & Funzionale

Obbligatorio & Il sistema deve segnalare un errore se il secondo inserimento password non coincide con il primo & 



\hyperlink{UC1.1}{UC1.1}\tabularnewline
\hline
\begin{tikzpicture}
\draw [->, thick] (0.2,0.2) -- (0.2,0.1) -- (1,0.1);
\end{tikzpicture} & \hypertarget{R-3F11.3}{R-3F11.3} & Funzionale

Obbligatorio & L'amministratore deve poter inviare un invito di registrazione ad un docente di cui conosce l'indirizzo email & \hyperlink{UC32.1}{UC32.1}

\hyperlink{UC32}{UC32}\tabularnewline
\hline
\begin{tikzpicture}
\draw [->, thick] (0.2,0.2) -- (0.2,0.1) -- (1,0.1);
\end{tikzpicture} & \hypertarget{R-3F11.4}{R-3F11.4} & Funzionale

Obbligatorio & Il sistema deve poter registrate un docente che ha ricevuto un invito di registrazione & \hyperlink{UC1.1}{UC1.1}\tabularnewline
\hline
\begin{tikzpicture}
\draw [->, thick] (0.4,0.2) -- (0.4,0.1) -- (1,0.1);
\end{tikzpicture} & \hypertarget{R-3F11.4.1}{R-3F11.4.1} & Funzionale

Obbligatorio & La registrazione di un docente richiede l’inserimento della mail, password, nome e cognome & 







\hyperlink{UC1.1}{UC1.1}\tabularnewline
\hline
\begin{tikzpicture}
\draw [->, thick] (0.4,0.2) -- (0.4,0.1) -- (1,0.1);
\end{tikzpicture} & \hypertarget{R-3F11.4.2}{R-3F11.4.2} & Funzionale

Obbligatorio & Il sistema deve permettere la visualizzazione di un messaggio di errore se la mail del docente è già presente nel sistema & 

\hyperlink{UC1.1}{UC1.1}\tabularnewline
\hline
\begin{tikzpicture}
\draw [->, thick] (0.4,0.2) -- (0.4,0.1) -- (1,0.1);
\end{tikzpicture} & \hypertarget{R-3F11.4.3}{R-3F11.4.3} & Funzionale

Obbligatorio & Il sistema deve permettere la visualizzazione di un messaggio di errore se la password del docente è troppo corta & 

\hyperlink{UC1.1}{UC1.1}\tabularnewline
\hline
\begin{tikzpicture}
\draw [->, thick] (0.2,0.2) -- (0.2,0.1) -- (1,0.1);
\end{tikzpicture} & \hypertarget{R-3F11.5}{R-3F11.5} & Funzionale

Obbligatorio & Il sistema deve permettere la creazione di un nuovo amministratore da parte del proprietario & \hyperlink{UC33.1}{UC33.1}

\hyperlink{UC33}{UC33}\tabularnewline
\hline
 & \hypertarget{R-3V12}{R-3V12} & Vincolo

Obbligatorio & Il database server side deve essere MongoDB & Capitolato\tabularnewline
\hline
 & \hypertarget{R-3F13}{R-3F13} & Funzionale

Obbligatorio & Il sistema deve gestire le azioni di amministratore & \hyperlink{UC32}{UC32}\tabularnewline
\hline
\begin{tikzpicture}
\draw [->, thick] (0.2,0.2) -- (0.2,0.1) -- (1,0.1);
\end{tikzpicture} & \hypertarget{R-3F13.1}{R-3F13.1} & Funzionale

Obbligatorio & Il sistema deve dare la possibilità ad un amministratore di invitare un docente alla registrazione nel sistema & \hyperlink{UC32.1}{UC32.1}

\hyperlink{UC32}{UC32}\tabularnewline
\hline
\begin{tikzpicture}
\draw [->, thick] (0.2,0.2) -- (0.2,0.1) -- (1,0.1);
\end{tikzpicture} & \hypertarget{R-3F13.2}{R-3F13.2} & Funzionale

Obbligatorio & Il sistema deve dare la possibilità ad un amministratore la possibilità di rimuovere un docente dal sistema & \hyperlink{UC32.2}{UC32.2}

\hyperlink{UC32}{UC32}\tabularnewline
\hline
\begin{tikzpicture}
\draw [->, thick] (0.2,0.2) -- (0.2,0.1) -- (1,0.1);
\end{tikzpicture} & \hypertarget{R-3F13.3}{R-3F13.3} & Funzionale

Obbligatorio & Il sistema deve dare la possibilità ad un amministratore di rimuovere uno studente dal sistema  & 

\hyperlink{UC32}{UC32}\tabularnewline
\hline
\begin{tikzpicture}
\draw [->, thick] (0.2,0.2) -- (0.2,0.1) -- (1,0.1);
\end{tikzpicture} & \hypertarget{R-3F13.4}{R-3F13.4} & Funzionale

Obbligatorio & Il sistema deve segnalare un errore nel caso un amministratore tenti di inviare un invito ad un indirizzo email già presente nel sistema & 

\hyperlink{UC32.1}{UC32.1}

\hyperlink{UC32}{UC32}\tabularnewline
\hline
 & \hypertarget{R-2F14}{R-2F14} & Funzionale

Opzionale & Il sistema deve prevedere la gestione di classi da parte di un docente & \hyperlink{UC9}{UC9}\tabularnewline
\hline
\begin{tikzpicture}
\draw [->, thick] (0.2,0.2) -- (0.2,0.1) -- (1,0.1);
\end{tikzpicture} & \hypertarget{R-2F14.1}{R-2F14.1} & Funzionale

Opzionale & Un docente deve avere la possibilità di creare una nuova classe & \hyperlink{UC9.1}{UC9.1}

\hyperlink{UC9}{UC9}\tabularnewline
\hline
\begin{tikzpicture}
\draw [->, thick] (0.4,0.2) -- (0.4,0.1) -- (1,0.1);
\end{tikzpicture} & \hypertarget{R-2F14.1.1}{R-2F14.1.1} & Funzionale

Opzionale & Un docente deve poter specificare il nome della nuova classe & \hyperlink{UC9.1.1}{UC9.1.1}

\hyperlink{UC9.1}{UC9.1}

\hyperlink{UC9}{UC9}\tabularnewline
\hline
\begin{tikzpicture}
\draw [->, thick] (0.6,0.2) -- (0.6,0.1) -- (1,0.1);
\end{tikzpicture} & \hypertarget{R-2F14.1.1.1}{R-2F14.1.1.1} & Funzionale

Opzionale & Nel caso il nome della classe sia già presente il sistema deve segnalare un errore & \hyperlink{UC9.1.1}{UC9.1.1}

\hyperlink{UC9.1}{UC9.1}

\hyperlink{UC9.4}{UC9.4}

\hyperlink{UC9}{UC9}\tabularnewline
\hline
\begin{tikzpicture}
\draw [->, thick] (0.4,0.2) -- (0.4,0.1) -- (1,0.1);
\end{tikzpicture} & \hypertarget{R-2F14.1.2}{R-2F14.1.2} & Funzionale

Opzionale & Un docente deve poter specificare gli argomenti di una nuova classe & \hyperlink{UC9.1.2}{UC9.1.2}

\hyperlink{UC9.1}{UC9.1}

\hyperlink{UC9}{UC9}\tabularnewline
\hline
\begin{tikzpicture}
\draw [->, thick] (0.4,0.2) -- (0.4,0.1) -- (1,0.1);
\end{tikzpicture} & \hypertarget{R-2F14.1.3}{R-2F14.1.3} & Funzionale

Opzionale & Un docente deve poter specificare la password relativa ad una nuova classe & \hyperlink{UC9.1.3}{UC9.1.3}

\hyperlink{UC9.1}{UC9.1}

\hyperlink{UC9}{UC9}\tabularnewline
\hline
\begin{tikzpicture}
\draw [->, thick] (0.2,0.2) -- (0.2,0.1) -- (1,0.1);
\end{tikzpicture} & \hypertarget{R-2F14.2}{R-2F14.2} & Funzionale

Opzionale & Un docente deve avere la possibilità di eliminare una classe esistente & \hyperlink{UC9.3}{UC9.3}

\hyperlink{UC9}{UC9}\tabularnewline
\hline
\begin{tikzpicture}
\draw [->, thick] (0.2,0.2) -- (0.2,0.1) -- (1,0.1);
\end{tikzpicture} & \hypertarget{R-2F14.3}{R-2F14.3} & Funzionale

Opzionale & Un docente deve poter modificare una classe & \hyperlink{UC9.2}{UC9.2}

\hyperlink{UC9}{UC9}\tabularnewline
\hline
\begin{tikzpicture}
\draw [->, thick] (0.4,0.2) -- (0.4,0.1) -- (1,0.1);
\end{tikzpicture} & \hypertarget{R-2F14.3.1}{R-2F14.3.1} & Funzionale

Opzionale & Un docente deve poter modificare il nome di una propria classe esistente & \hyperlink{UC9.1.1}{UC9.1.1}

\hyperlink{UC9.2}{UC9.2}

\hyperlink{UC9}{UC9}\tabularnewline
\hline
\begin{tikzpicture}
\draw [->, thick] (0.4,0.2) -- (0.4,0.1) -- (1,0.1);
\end{tikzpicture} & \hypertarget{R-2F14.3.2}{R-2F14.3.2} & Funzionale

Opzionale & Un docente deve poter modificare gli argomenti di una sua classe esistente & \hyperlink{UC9.1.2}{UC9.1.2}

\hyperlink{UC9.2}{UC9.2}

\hyperlink{UC9}{UC9}\tabularnewline
\hline
\begin{tikzpicture}
\draw [->, thick] (0.4,0.2) -- (0.4,0.1) -- (1,0.1);
\end{tikzpicture} & \hypertarget{R-2F14.3.3}{R-2F14.3.3} & Funzionale

Opzionale & Un docente deve poter modificare la password di una delle proprie classi & \hyperlink{UC9.1.3}{UC9.1.3}

\hyperlink{UC9.2}{UC9.2}

\hyperlink{UC9}{UC9}\tabularnewline
\hline
 & \hypertarget{R-3F15}{R-3F15} & Funzionale

Obbligatorio & Il sistema deve consentire ad un utente di gestire i propri dati memorizzati nel sistema & \hyperlink{UC5}{UC5}\tabularnewline
\hline
\begin{tikzpicture}
\draw [->, thick] (0.2,0.2) -- (0.2,0.1) -- (1,0.1);
\end{tikzpicture} & \hypertarget{R-3F15.1}{R-3F15.1} & Funzionale

Obbligatorio & Il sistema deve consentire ad un utente di modificare la propria mail & 

\hyperlink{UC5}{UC5}\tabularnewline
\hline
\begin{tikzpicture}
\draw [->, thick] (0.4,0.2) -- (0.4,0.1) -- (1,0.1);
\end{tikzpicture} & \hypertarget{R-3F15.1.1}{R-3F15.1.1} & Funzionale

Obbligatorio & Il sistema deve gestire eventuali errori di inserimento mail  & 



\hyperlink{UC5}{UC5}\tabularnewline
\hline
\begin{tikzpicture}
\draw [->, thick] (0.4,0.2) -- (0.4,0.1) -- (1,0.1);
\end{tikzpicture} & \hypertarget{R-3F15.1.2}{R-3F15.1.2} & Funzionale

Obbligatorio & Il sistema deve segnalare un errore in caso di inserimento di email già presente nel sistema & 



\hyperlink{UC5}{UC5}\tabularnewline
\hline
\begin{tikzpicture}
\draw [->, thick] (0.2,0.2) -- (0.2,0.1) -- (1,0.1);
\end{tikzpicture} & \hypertarget{R-3F15.2}{R-3F15.2} & Funzionale

Obbligatorio & Il sistema deve consentire ad un utente di modificare la propria password di accesso al sistema & 

\hyperlink{UC5}{UC5}\tabularnewline
\hline
\begin{tikzpicture}
\draw [->, thick] (0.4,0.2) -- (0.4,0.1) -- (1,0.1);
\end{tikzpicture} & \hypertarget{R-3F15.2.1}{R-3F15.2.1} & Funzionale

Obbligatorio & Il sistema deve richiedere l'inserimento della vecchia password all'utente per consentire la modifica della password & 



\hyperlink{UC5}{UC5}\tabularnewline
\hline
\begin{tikzpicture}
\draw [->, thick] (0.6,0.2) -- (0.6,0.1) -- (1,0.1);
\end{tikzpicture} & \hypertarget{R-3F15.2.1.1}{R-3F15.2.1.1} & Funzionale

Obbligatorio & Il sistema deve segnalare l'inserimento di una password scorretta & 



\hyperlink{UC5}{UC5}\tabularnewline
\hline
\begin{tikzpicture}
\draw [->, thick] (0.4,0.2) -- (0.4,0.1) -- (1,0.1);
\end{tikzpicture} & \hypertarget{R-3F15.2.2}{R-3F15.2.2} & Funzionale

Obbligatorio & Il sistema deve richiedere di inserire la nuova password due volte & 



\hyperlink{UC5}{UC5}\tabularnewline
\hline
\begin{tikzpicture}
\draw [->, thick] (0.6,0.2) -- (0.6,0.1) -- (1,0.1);
\end{tikzpicture} & \hypertarget{R-3F15.2.2.1}{R-3F15.2.2.1} & Funzionale

Obbligatorio & Il sistema deve verificare che la nuova password sia di almeno sei caratteri e segnalarlo in caso contrario & 





\hyperlink{UC5}{UC5}\tabularnewline
\hline
\begin{tikzpicture}
\draw [->, thick] (0.6,0.2) -- (0.6,0.1) -- (1,0.1);
\end{tikzpicture} & \hypertarget{R-3F15.2.2.2}{R-3F15.2.2.2} & Funzionale

Obbligatorio & Il sistema deve controllare che i due campi per l'inserimento della nuova password coincidano e segnalarlo in caso contrario. & 





\hyperlink{UC5}{UC5}\tabularnewline
\hline
\begin{tikzpicture}
\draw [->, thick] (0.2,0.2) -- (0.2,0.1) -- (1,0.1);
\end{tikzpicture} & \hypertarget{R-3F15.3}{R-3F15.3} & Funzionale

Obbligatorio & Il sistema deve consentire all'utente di modificare i propri dati personali & \hyperlink{UC5.1}{UC5.1}

\hyperlink{UC5}{UC5}\tabularnewline
\hline
\begin{tikzpicture}
\draw [->, thick] (0.4,0.2) -- (0.4,0.1) -- (1,0.1);
\end{tikzpicture} & \hypertarget{R-3F15.3.1}{R-3F15.3.1} & Funzionale

Obbligatorio & Il sistema deve consentire all'utente la modifica del proprio nome  & \hyperlink{UC5.1.1}{UC5.1.1}

\hyperlink{UC5.1}{UC5.1}

\hyperlink{UC5}{UC5}\tabularnewline
\hline
\begin{tikzpicture}
\draw [->, thick] (0.4,0.2) -- (0.4,0.1) -- (1,0.1);
\end{tikzpicture} & \hypertarget{R-3F15.3.2}{R-3F15.3.2} & Funzionale

Obbligatorio & Il sistema deve consentire all'utente di modificare il proprio cognome & 

\hyperlink{UC5.1}{UC5.1}

\hyperlink{UC5}{UC5}\tabularnewline
\hline
\begin{tikzpicture}
\draw [->, thick] (0.4,0.2) -- (0.4,0.1) -- (1,0.1);
\end{tikzpicture} & \hypertarget{R-3F15.3.3}{R-3F15.3.3} & Funzionale

Obbligatorio & Il sistema deve controllare che i campi obbligatori nome e cognome siano stati inseriti e segnalarlo in caso contrario & 

\hyperlink{UC5.1.1}{UC5.1.1}

\hyperlink{UC5.1}{UC5.1}

\hyperlink{UC5.2}{UC5.2}

\hyperlink{UC5}{UC5}\tabularnewline
\hline
 & \hypertarget{R-3F16}{R-3F16} & Funzionale

Obbligatorio & Il sistema deve gestire la ricerca di un questionario & \hyperlink{UC18}{UC18}\tabularnewline
\hline
\begin{tikzpicture}
\draw [->, thick] (0.2,0.2) -- (0.2,0.1) -- (1,0.1);
\end{tikzpicture} & \hypertarget{R-3F16.1}{R-3F16.1} & Funzionale

Obbligatorio & Il sistema deve gestire la ricerca di un questionario per titolo & \hyperlink{UC18}{UC18}

\hyperlink{UC19}{UC19}\tabularnewline
\hline
\begin{tikzpicture}
\draw [->, thick] (0.2,0.2) -- (0.2,0.1) -- (1,0.1);
\end{tikzpicture} & \hypertarget{R-2F16.2}{R-2F16.2} & Funzionale

Opzionale & Il sistema deve gestire la ricerca di un questionario per classe & \hyperlink{UC18}{UC18}

\hyperlink{UC20}{UC20}\tabularnewline
\hline
\begin{tikzpicture}
\draw [->, thick] (0.2,0.2) -- (0.2,0.1) -- (1,0.1);
\end{tikzpicture} & \hypertarget{R-3F16.3}{R-3F16.3} & Funzionale

Obbligatorio & Il sistema deve gestire la ricerca di un questionario per docente & \hyperlink{UC18}{UC18}

\hyperlink{UC22}{UC22}\tabularnewline
\hline
\begin{tikzpicture}
\draw [->, thick] (0.2,0.2) -- (0.2,0.1) -- (1,0.1);
\end{tikzpicture} & \hypertarget{R-3F16.4}{R-3F16.4} & Funzionale

Obbligatorio & Il sistema deve gestire la ricerca di un questionario per argomento & \hyperlink{UC18}{UC18}

\hyperlink{UC21}{UC21}\tabularnewline
\hline
\begin{tikzpicture}
\draw [->, thick] (0.2,0.2) -- (0.2,0.1) -- (1,0.1);
\end{tikzpicture} & \hypertarget{R-3F16.5}{R-3F16.5} & Funzionale

Obbligatorio & Il sistema deve gestire la ricerca di un questionario per difficoltà & \hyperlink{UC18}{UC18}

\hyperlink{UC23}{UC23}\tabularnewline
\hline
 & \hypertarget{R-2F17}{R-2F17} & Funzionale

Opzionale & Il sistema deve consentire l'iscrizione ad una classe da parte di uno studente & \hyperlink{UC16}{UC16}\tabularnewline
\hline
\begin{tikzpicture}
\draw [->, thick] (0.2,0.2) -- (0.2,0.1) -- (1,0.1);
\end{tikzpicture} & \hypertarget{R-2F17.1}{R-2F17.1} & Funzionale

Opzionale & Il sistema deve chiedere l'inserimento della password prima di confermare l'iscrizione ad una classe & \hyperlink{UC16.1}{UC16.1}

\hyperlink{UC16}{UC16}\tabularnewline
\hline
\begin{tikzpicture}
\draw [->, thick] (0.4,0.2) -- (0.4,0.1) -- (1,0.1);
\end{tikzpicture} & \hypertarget{R-2F17.1.1}{R-2F17.1.1} & Funzionale

Opzionale & Il sistema deve verificare la password prima di confermare l'iscrizione dello studente alla classe. & \hyperlink{UC16.2}{UC16.2}

\hyperlink{UC16}{UC16}\tabularnewline
\hline
\begin{tikzpicture}
\draw [->, thick] (0.2,0.2) -- (0.2,0.1) -- (1,0.1);
\end{tikzpicture} & \hypertarget{R-2F17.2}{R-2F17.2} & Funzionale

Opzionale & Il sistema deve consentire di ricercare la classe a cui iscriversi & \hyperlink{UC29}{UC29}

\hyperlink{UC30}{UC30}

\hyperlink{UC31}{UC31}\tabularnewline
\hline
 & \hypertarget{R-3F18}{R-3F18} & Funzionale

Obbligatorio & Il sistema deve dare la possibilità di eseguire un questionario ad uno studente & \hyperlink{UC15}{UC15}\tabularnewline
\hline
\begin{tikzpicture}
\draw [->, thick] (0.2,0.2) -- (0.2,0.1) -- (1,0.1);
\end{tikzpicture} & \hypertarget{R-3F18.1}{R-3F18.1} & Funzionale

Obbligatorio & Il sistema deve dare la possibilità ad uno studente di rispondere ad una domanda & \hyperlink{UC15.1}{UC15.1}

\hyperlink{UC15}{UC15}\tabularnewline
\hline
\begin{tikzpicture}
\draw [->, thick] (0.4,0.2) -- (0.4,0.1) -- (1,0.1);
\end{tikzpicture} & \hypertarget{R-3F18.1.1}{R-3F18.1.1} & Funzionale

Obbligatorio & Il sistema deve dare la possibilità ad uno studente di rispondere ad una domanda di tipo vero/falso & \hyperlink{UC15.1}{UC15.1}

\hyperlink{UC15.4}{UC15.4}\tabularnewline
\hline
\begin{tikzpicture}
\draw [->, thick] (0.4,0.2) -- (0.4,0.1) -- (1,0.1);
\end{tikzpicture} & \hypertarget{R-3F18.1.2}{R-3F18.1.2} & Funzionale

Obbligatorio & Il sistema deve dare la possibilità ad uno studente di rispondere ad una domanda a scelta multipla & \hyperlink{UC15.1}{UC15.1}

\hyperlink{UC15.5}{UC15.5}\tabularnewline
\hline
\begin{tikzpicture}
\draw [->, thick] (0.4,0.2) -- (0.4,0.1) -- (1,0.1);
\end{tikzpicture} & \hypertarget{R-2F18.1.3}{R-2F18.1.3} & Funzionale

Opzionale & Il sistema deve dare la possibilità ad uno studente di rispondere ad una domanda a risposta multipla ovvero una domanda a scelta multipla  con più risposte giuste & \hyperlink{UC15.1}{UC15.1}

\hyperlink{UC15.6}{UC15.6}\tabularnewline
\hline
\begin{tikzpicture}
\draw [->, thick] (0.4,0.2) -- (0.4,0.1) -- (1,0.1);
\end{tikzpicture} & \hypertarget{R-2F18.1.4}{R-2F18.1.4} & Funzionale

Opzionale & Il sistema deve dare la possibilità ad uno studente di rispondere ad una domanda di tipo testo con parole omesse & \hyperlink{UC15.1}{UC15.1}

\hyperlink{UC15.7}{UC15.7}\tabularnewline
\hline
\begin{tikzpicture}
\draw [->, thick] (0.4,0.2) -- (0.4,0.1) -- (1,0.1);
\end{tikzpicture} & \hypertarget{R-2F18.1.5}{R-2F18.1.5} & Funzionale

Opzionale & Il sistema deve dare la possibilità ad uno studente di rispondere ad una domanda con associazione di parole & \hyperlink{UC15.1}{UC15.1}

\hyperlink{UC15.8}{UC15.8}\tabularnewline
\hline
\begin{tikzpicture}
\draw [->, thick] (0.4,0.2) -- (0.4,0.1) -- (1,0.1);
\end{tikzpicture} & \hypertarget{R-2F18.1.6}{R-2F18.1.6} & Funzionale

Opzionale & Il sistema deve dare la possibilità ad uno studente di rispondere ad una domanda a risposta aperta & \hyperlink{UC15.1}{UC15.1}

\hyperlink{UC15.9}{UC15.9}\tabularnewline
\hline
\begin{tikzpicture}
\draw [->, thick] (0.2,0.2) -- (0.2,0.1) -- (1,0.1);
\end{tikzpicture} & \hypertarget{R-3F18.2}{R-3F18.2} & Funzionale

Obbligatorio & Il sistema deve richiedere allo studente la conferma alla fine della compilazione di un questionario, perché questo sia considerato eseguito & \hyperlink{UC15.3}{UC15.3}

\hyperlink{UC15}{UC15}\tabularnewline
\hline
\begin{tikzpicture}
\draw [->, thick] (0.4,0.2) -- (0.4,0.1) -- (1,0.1);
\end{tikzpicture} & \hypertarget{R-3F18.2.1}{R-3F18.2.1} & Funzionale

Obbligatorio & Il sistema deve segnalare un errore nel caso una domanda del questionario non sia stata compilata & \hyperlink{UC15.2}{UC15.2}

\hyperlink{UC15.3}{UC15.3}

\hyperlink{UC15}{UC15}\tabularnewline
\hline
\begin{tikzpicture}
\draw [->, thick] (0.2,0.2) -- (0.2,0.1) -- (1,0.1);
\end{tikzpicture} & \hypertarget{R-3F18.3}{R-3F18.3} & Funzionale

Obbligatorio & Il sistema deve dare la possibilità ad uno studente di andare alla domanda successiva durante l'esecuzione del questionario & \hyperlink{UC15}{UC15}\tabularnewline
\hline
\begin{tikzpicture}
\draw [->, thick] (0.2,0.2) -- (0.2,0.1) -- (1,0.1);
\end{tikzpicture} & \hypertarget{R-3F18.4}{R-3F18.4} & Funzionale

Obbligatorio & Il sistema deve dare la possibilità ad uno studente di andare alla domanda precedente durante l'esecuzione del questionario & \hyperlink{UC15}{UC15}\tabularnewline
\hline
 & \hypertarget{R-3F19}{R-3F19} & Funzionale

Obbligatorio & Il sistema deve consentire all'utente autenticato di effettuare il logout & \hyperlink{UC6}{UC6}\tabularnewline
\hline
 & \hypertarget{R-3F20}{R-3F20} & Funzionale

Obbligatorio & Il sistema deve gestire le operazioni che può fare il proprietario del sistema & \hyperlink{UC33}{UC33}\tabularnewline
\hline
\begin{tikzpicture}
\draw [->, thick] (0.2,0.2) -- (0.2,0.1) -- (1,0.1);
\end{tikzpicture} & \hypertarget{R-3F20.1}{R-3F20.1} & Funzionale

Obbligatorio & Il sistema deve assicurare l'esistenza di un unico utente proprietario & \hyperlink{UC33}{UC33}\tabularnewline
\hline
\begin{tikzpicture}
\draw [->, thick] (0.2,0.2) -- (0.2,0.1) -- (1,0.1);
\end{tikzpicture} & \hypertarget{R-3F20.2}{R-3F20.2} & Funzionale

Obbligatorio & Il sistema deve consentire solo all'utente proprietario di aggiungere un amministratore & \hyperlink{UC33.1}{UC33.1}

\hyperlink{UC33}{UC33}\tabularnewline
\hline
\begin{tikzpicture}
\draw [->, thick] (0.4,0.2) -- (0.4,0.1) -- (1,0.1);
\end{tikzpicture} & \hypertarget{R-3F20.2.1}{R-3F20.2.1} & Funzionale

Obbligatorio & Il sistema deve verificare che l'e-mail non corrisponda già ad un amministratore del sistema & 

\hyperlink{UC33.1}{UC33.1}

\hyperlink{UC33}{UC33}\tabularnewline
\hline
\begin{tikzpicture}
\draw [->, thick] (0.2,0.2) -- (0.2,0.1) -- (1,0.1);
\end{tikzpicture} & \hypertarget{R-3F20.3}{R-3F20.3} & Funzionale

Obbligatorio & Il sistema deve consentire solo all'utente proprietario di eliminare un amministratore & \hyperlink{UC33.2}{UC33.2}

\hyperlink{UC33}{UC33}\tabularnewline
\hline
 & \hypertarget{R-3F21}{R-3F21} & Funzionale

Obbligatorio & Il sistema deve consentire ad un docente o ad uno studente di ricercare una domanda presente nel sistema & \hyperlink{UC24}{UC24}\tabularnewline
\hline
\begin{tikzpicture}
\draw [->, thick] (0.2,0.2) -- (0.2,0.1) -- (1,0.1);
\end{tikzpicture} & \hypertarget{R-3F21.1}{R-3F21.1} & Funzionale

Obbligatorio & Il sistema deve consentire a uno studente o ad un docente di ricercare una domanda per argomento & \hyperlink{UC24}{UC24}

\hyperlink{UC26}{UC26}\tabularnewline
\hline
\begin{tikzpicture}
\draw [->, thick] (0.2,0.2) -- (0.2,0.1) -- (1,0.1);
\end{tikzpicture} & \hypertarget{R-3F21.2}{R-3F21.2} & Funzionale

Obbligatorio & Il sistema deve consentire a uno studente o ad un docente di ricercare una domanda per docente & \hyperlink{UC24}{UC24}

\hyperlink{UC28}{UC28}\tabularnewline
\hline
\begin{tikzpicture}
\draw [->, thick] (0.2,0.2) -- (0.2,0.1) -- (1,0.1);
\end{tikzpicture} & \hypertarget{R-3F21.3}{R-3F21.3} & Funzionale

Obbligatorio & Il sistema deve consentire a uno studente o ad un docente di ricercare una domanda per difficoltà & \hyperlink{UC24}{UC24}

\hyperlink{UC27}{UC27}\tabularnewline
\hline
\begin{tikzpicture}
\draw [->, thick] (0.2,0.2) -- (0.2,0.1) -- (1,0.1);
\end{tikzpicture} & \hypertarget{R-3F21.4}{R-3F21.4} & Funzionale

Obbligatorio & Il sistema deve consentire a uno studente o ad un docente di ricercare una domanda per keywords & \hyperlink{UC24}{UC24}

\hyperlink{UC25}{UC25}\tabularnewline
\hline
 & \hypertarget{R-2F22}{R-2F22} & Funzionale

Opzionale & Il sistema deve fornire la possibilità ad uno studente  di ricerca di una classe & \hyperlink{UC29}{UC29}\tabularnewline
\hline
\begin{tikzpicture}
\draw [->, thick] (0.2,0.2) -- (0.2,0.1) -- (1,0.1);
\end{tikzpicture} & \hypertarget{R-2F22.1}{R-2F22.1} & Funzionale

Opzionale & Il sistema deve fornire la possibilità ad uno studente  di ricerca di una classe per argomento & \hyperlink{UC29}{UC29}

\hyperlink{UC31}{UC31}\tabularnewline
\hline
\begin{tikzpicture}
\draw [->, thick] (0.2,0.2) -- (0.2,0.1) -- (1,0.1);
\end{tikzpicture} & \hypertarget{R-2F22.2}{R-2F22.2} & Funzionale

Opzionale & Il sistema deve fornire la possibilità ad uno studente  di ricerca di una classe per docente & \hyperlink{UC29}{UC29}

\hyperlink{UC30}{UC30}\tabularnewline
\hline
 & \hypertarget{R-3F23}{R-3F23} & Funzionale

Obbligatorio & Il sistema deve dare la possibilità di recupero password ad un utente non autenticato inviando una mail con la nuova password
 & \tabularnewline
\hline
\begin{tikzpicture}
\draw [->, thick] (0.2,0.2) -- (0.2,0.1) -- (1,0.1);
\end{tikzpicture} & \hypertarget{R-3F23.1}{R-3F23.1} & Funzionale

Obbligatorio & Il sistema deve dare la possibilità di inserire  mail per recupero password ad un utente non autenticato & 

\tabularnewline
\hline
\begin{tikzpicture}
\draw [->, thick] (0.2,0.2) -- (0.2,0.1) -- (1,0.1);
\end{tikzpicture} & \hypertarget{R-3F23.2}{R-3F23.2} & Funzionale

Obbligatorio & Il sistema deve segnalare ad un utente non autenticato che l'email inserito per recupero password non è presente nel sistema & 

\tabularnewline
\hline
 & \hypertarget{R-3F24}{R-3F24} & Funzionale

Obbligatorio & Il sistema deve permettere ad un docente di gestire gli argomenti & \hyperlink{UC10}{UC10}\tabularnewline
\hline
\begin{tikzpicture}
\draw [->, thick] (0.2,0.2) -- (0.2,0.1) -- (1,0.1);
\end{tikzpicture} & \hypertarget{R-3F24.1}{R-3F24.1} & Funzionale

Obbligatorio & Il sistema deve consentire ad un docente di creare un nuovo argomento & \hyperlink{UC10.2}{UC10.2}

\hyperlink{UC10}{UC10}\tabularnewline
\hline
\begin{tikzpicture}
\draw [->, thick] (0.4,0.2) -- (0.4,0.1) -- (1,0.1);
\end{tikzpicture} & \hypertarget{R-3F24.1.1}{R-3F24.1.1} & Funzionale

Obbligatorio & Il sistema deve verificare se l'argomento sia già presente nel sistema e notificarlo in caso contrario & \hyperlink{UC10.2}{UC10.2}

\hyperlink{UC10.3}{UC10.3}

\hyperlink{UC10}{UC10}\tabularnewline
\hline
\begin{tikzpicture}
\draw [->, thick] (0.2,0.2) -- (0.2,0.1) -- (1,0.1);
\end{tikzpicture} & \hypertarget{R-3F24.2}{R-3F24.2} & Funzionale

Obbligatorio & Il sistema deve permettere ad un docente di eliminare un argomento & \hyperlink{UC10.5}{UC10.5}

\hyperlink{UC10}{UC10}\tabularnewline
\hline
\begin{tikzpicture}
\draw [->, thick] (0.4,0.2) -- (0.4,0.1) -- (1,0.1);
\end{tikzpicture} & \hypertarget{R-3F24.2.1}{R-3F24.2.1} & Funzionale

Obbligatorio & Il sistema deve permettere al docente di selezionare un argomento presente da eliminare & \hyperlink{UC10.1}{UC10.1}

\hyperlink{UC10.5}{UC10.5}

\hyperlink{UC10}{UC10}\tabularnewline
\hline
\begin{tikzpicture}
\draw [->, thick] (0.4,0.2) -- (0.4,0.1) -- (1,0.1);
\end{tikzpicture} & \hypertarget{R-3F24.2.2}{R-3F24.2.2} & Funzionale

Obbligatorio & Il sistema deve impedire al docente di eliminare un argomento corrisponde a domande o questionari presenti nel sistema & \hyperlink{UC10.5}{UC10.5}

\hyperlink{UC10.6}{UC10.6}

\hyperlink{UC10}{UC10}\tabularnewline
\hline
\begin{tikzpicture}
\draw [->, thick] (0.2,0.2) -- (0.2,0.1) -- (1,0.1);
\end{tikzpicture} & \hypertarget{R-3F24.3}{R-3F24.3} & Funzionale

Obbligatorio & Il sistema deve permettere ad un docente di modificare un argomento & \hyperlink{UC10.4}{UC10.4}

\hyperlink{UC10}{UC10}\tabularnewline
\hline
\begin{tikzpicture}
\draw [->, thick] (0.2,0.2) -- (0.2,0.1) -- (1,0.1);
\end{tikzpicture} & \hypertarget{R-3F24.4}{R-3F24.4} & Funzionale

Obbligatorio & Il sistema deve permettere ad un docente di esplorare gli argomenti presenti & \hyperlink{UC10.1}{UC10.1}

\hyperlink{UC10}{UC10}\tabularnewline
\hline
 & \hypertarget{R-3F25}{R-3F25} & Funzionale

Obbligatorio & Il sistema deve consentire ad un docente di visualizzare statistiche & \hyperlink{UC11}{UC11}\tabularnewline
\hline
\begin{tikzpicture}
\draw [->, thick] (0.2,0.2) -- (0.2,0.1) -- (1,0.1);
\end{tikzpicture} & \hypertarget{R-3F25.1}{R-3F25.1} & Funzionale

Obbligatorio & Il sistema deve permettere ad un docente di visualizzare le statistiche relative ad una domanda & \hyperlink{UC11}{UC11}

\hyperlink{UC12}{UC12}\tabularnewline
\hline
\begin{tikzpicture}
\draw [->, thick] (0.4,0.2) -- (0.4,0.1) -- (1,0.1);
\end{tikzpicture} & \hypertarget{R-3F25.1.1}{R-3F25.1.1} & Funzionale

Obbligatorio & Il sistema deve permettere ad un docente di ricercare la domanda di cui visualizzare le statistiche  & \hyperlink{UC11}{UC11}

\hyperlink{UC12}{UC12}

\hyperlink{UC24}{UC24}\tabularnewline
\hline
\begin{tikzpicture}
\draw [->, thick] (0.2,0.2) -- (0.2,0.1) -- (1,0.1);
\end{tikzpicture} & \hypertarget{R-3F25.2}{R-3F25.2} & Funzionale

Obbligatorio & Il sistema deve permettere ad un docente di visualizzare le statistiche relative ad un questionario & \hyperlink{UC11}{UC11}

\hyperlink{UC13}{UC13}\tabularnewline
\hline
\begin{tikzpicture}
\draw [->, thick] (0.4,0.2) -- (0.4,0.1) -- (1,0.1);
\end{tikzpicture} & \hypertarget{R-3F25.2.1}{R-3F25.2.1} & Funzionale

Obbligatorio & Il sistema deve consentire ad un docente di ricercare il questionario di cui visualizzare le statistiche & \hyperlink{UC11}{UC11}

\hyperlink{UC13}{UC13}

\hyperlink{UC18}{UC18}\tabularnewline
\hline
\begin{tikzpicture}
\draw [->, thick] (0.2,0.2) -- (0.2,0.1) -- (1,0.1);
\end{tikzpicture} & \hypertarget{R-2F25.3}{R-2F25.3} & Funzionale

Opzionale & Il sistema deve permettere ad un docente di visualizzare le statistiche relative ad una delle proprie classi & \hyperlink{UC11}{UC11}

\hyperlink{UC14}{UC14}\tabularnewline
\hline
\begin{tikzpicture}
\draw [->, thick] (0.4,0.2) -- (0.4,0.1) -- (1,0.1);
\end{tikzpicture} & \hypertarget{R-2F25.3.1}{R-2F25.3.1} & Funzionale

Opzionale & Un docente deve poter visualizzare i risultati dei suoi questionari in base alla classe & \hyperlink{UC11}{UC11}

\hyperlink{UC14.2}{UC14.2}

\hyperlink{UC14}{UC14}\tabularnewline
\hline
\begin{tikzpicture}
\draw [->, thick] (0.4,0.2) -- (0.4,0.1) -- (1,0.1);
\end{tikzpicture} & \hypertarget{R-2F25.3.2}{R-2F25.3.2} & Funzionale

Opzionale & Un docente deve poter selezionare tra le proprie classi quella di cui visualizzare le statistiche & \hyperlink{UC11}{UC11}

\hyperlink{UC14}{UC14}\tabularnewline
\hline
\begin{tikzpicture}
\draw [->, thick] (0.4,0.2) -- (0.4,0.1) -- (1,0.1);
\end{tikzpicture} & \hypertarget{R-2F25.3.3}{R-2F25.3.3} & Funzionale

Opzionale & Un docente deve poter visualizzare un sommario delle statistiche della classe & \hyperlink{UC11}{UC11}

\hyperlink{UC14.3}{UC14.3}

\hyperlink{UC14}{UC14}\tabularnewline
\hline
\begin{tikzpicture}
\draw [->, thick] (0.4,0.2) -- (0.4,0.1) -- (1,0.1);
\end{tikzpicture} & \hypertarget{R-2F25.3.4}{R-2F25.3.4} & Funzionale

Opzionale & Un docente deve poter visualizzare le statistiche relative ad un particolare studente della classe & \hyperlink{UC11}{UC11}

\hyperlink{UC14.4}{UC14.4}

\hyperlink{UC14}{UC14}\tabularnewline
\hline
\begin{tikzpicture}
\draw [->, thick] (0.6,0.2) -- (0.6,0.1) -- (1,0.1);
\end{tikzpicture} & \hypertarget{R-3F25.3.4.1}{R-3F25.3.4.1} & Funzionale

Obbligatorio & Un docente deve poter cercare e selezionare un particolare questionario tra quelli eseguiti dallo studente di cui visualizzare le statistiche & \hyperlink{UC11}{UC11}

\hyperlink{UC14.4.1}{UC14.4.1}

\hyperlink{UC14.4}{UC14.4}

\hyperlink{UC14}{UC14}

\hyperlink{UC18}{UC18}\tabularnewline
\hline
\begin{tikzpicture}
\draw [->, thick] (0.4,0.2) -- (0.4,0.1) -- (1,0.1);
\end{tikzpicture} & \hypertarget{R-2F25.3.5}{R-2F25.3.5} & Funzionale

Opzionale & Un docente deve poter visualizzare i risultati delle sue domande in base alla classe & \hyperlink{UC11}{UC11}

\hyperlink{UC14.1}{UC14.1}

\hyperlink{UC14}{UC14}\tabularnewline
\hline
\begin{tikzpicture}
\draw [->, thick] (0.6,0.2) -- (0.6,0.1) -- (1,0.1);
\end{tikzpicture} & \hypertarget{R-2F25.3.5.1}{R-2F25.3.5.1} & Funzionale

Opzionale & Il sistema deve consentire ad un docente di ricercare una domanda tra quelle della classe & \hyperlink{UC11}{UC11}

\hyperlink{UC14.1}{UC14.1}

\hyperlink{UC14}{UC14}

\hyperlink{UC24}{UC24}\tabularnewline
\hline
 & \hypertarget{R-3F26}{R-3F26} & Funzionale

Obbligatorio & Il sistema deve consentire ad uno studente di visualizzare le proprie statistiche & \hyperlink{UC17}{UC17}\tabularnewline
\hline
\begin{tikzpicture}
\draw [->, thick] (0.2,0.2) -- (0.2,0.1) -- (1,0.1);
\end{tikzpicture} & \hypertarget{R-3F26.1}{R-3F26.1} & Funzionale

Obbligatorio & Il sistema deve permettere ad uno studente di visualizzare un sommario delle proprie statistiche & \hyperlink{UC17.3}{UC17.3}

\hyperlink{UC17}{UC17}\tabularnewline
\hline
\begin{tikzpicture}
\draw [->, thick] (0.4,0.2) -- (0.4,0.1) -- (1,0.1);
\end{tikzpicture} & \hypertarget{R-3F26.1.1}{R-3F26.1.1} & Funzionale

Obbligatorio & Il sistema deve mostrare allo studente il totale delle domande eseguite e il numero di risposte corrette & \hyperlink{UC17.3}{UC17.3}

\hyperlink{UC17}{UC17}\tabularnewline
\hline
\begin{tikzpicture}
\draw [->, thick] (0.4,0.2) -- (0.4,0.1) -- (1,0.1);
\end{tikzpicture} & \hypertarget{R-3F26.1.2}{R-3F26.1.2} & Funzionale

Obbligatorio & Il sistema deve mostrare allo studente il totale dei questionari eseguiti e la media dei risultati ottenuti & \hyperlink{UC17.3}{UC17.3}

\hyperlink{UC17}{UC17}\tabularnewline
\hline
\begin{tikzpicture}
\draw [->, thick] (0.4,0.2) -- (0.4,0.1) -- (1,0.1);
\end{tikzpicture} & \hypertarget{R-3F26.1.3}{R-3F26.1.3} & Funzionale

Obbligatorio & Il sistema deve mostrare allo studente la media delle difficoltà delle domande a cui ha risposto correttamente & \hyperlink{UC17.3}{UC17.3}

\hyperlink{UC17}{UC17}\tabularnewline
\hline
\begin{tikzpicture}
\draw [->, thick] (0.4,0.2) -- (0.4,0.1) -- (1,0.1);
\end{tikzpicture} & \hypertarget{R-3F26.1.4}{R-3F26.1.4} & Funzionale

Obbligatorio & Il sistema deve mostrare allo studente la media delle difficoltà delle domande a cui non ha risposto correttamente & \hyperlink{UC17.3}{UC17.3}

\hyperlink{UC17}{UC17}\tabularnewline
\hline
\begin{tikzpicture}
\draw [->, thick] (0.4,0.2) -- (0.4,0.1) -- (1,0.1);
\end{tikzpicture} & \hypertarget{R-1F26.1.5}{R-1F26.1.5} & Funzionale

Desiderabile & Per ogni argomento il sistema deve mostrare allo studente il totale delle domande eseguite, il numero di risposte corrette, la media delle difficoltà delle risposte corrette e la media delle difficoltà delle risposte non corrette. & \hyperlink{UC17.3}{UC17.3}

\hyperlink{UC17}{UC17}\tabularnewline
\hline
\begin{tikzpicture}
\draw [->, thick] (0.2,0.2) -- (0.2,0.1) -- (1,0.1);
\end{tikzpicture} & \hypertarget{R-3F26.2}{R-3F26.2} & Funzionale

Obbligatorio & Il sistema deve permettere ad uno studente di visualizzare le statistiche relative ad una domanda eseguita & \hyperlink{UC17.1}{UC17.1}

\hyperlink{UC17}{UC17}\tabularnewline
\hline
\begin{tikzpicture}
\draw [->, thick] (0.4,0.2) -- (0.4,0.1) -- (1,0.1);
\end{tikzpicture} & \hypertarget{R-3F26.2.1}{R-3F26.2.1} & Funzionale

Obbligatorio & Il sistema deve mostrare allo studente le risposte già date in passato alla domanda selezionata & \hyperlink{UC17.1}{UC17.1}

\hyperlink{UC17}{UC17}\tabularnewline
\hline
\begin{tikzpicture}
\draw [->, thick] (0.4,0.2) -- (0.4,0.1) -- (1,0.1);
\end{tikzpicture} & \hypertarget{R-3F26.2.2}{R-3F26.2.2} & Funzionale

Obbligatorio & Il sistema deve mostrare allo studente la difficoltà della domanda calcolata come funzione delle risposte date da tutti gli utenti & \hyperlink{UC17.1}{UC17.1}

\hyperlink{UC17}{UC17}\tabularnewline
\hline
\begin{tikzpicture}
\draw [->, thick] (0.2,0.2) -- (0.2,0.1) -- (1,0.1);
\end{tikzpicture} & \hypertarget{R-3F26.3}{R-3F26.3} & Funzionale

Obbligatorio & Il sistema deve permettere ad uno studente di visualizzare le statistiche e le risposte relative ad un questionario eseguito & \hyperlink{UC17.2}{UC17.2}

\hyperlink{UC17}{UC17}\tabularnewline
\hline
\begin{tikzpicture}
\draw [->, thick] (0.4,0.2) -- (0.4,0.1) -- (1,0.1);
\end{tikzpicture} & \hypertarget{R-3F26.3.1}{R-3F26.3.1} & Funzionale

Obbligatorio & Il sistema deve consentire ad uno studente di ricercare e selezionare uno dei questionari eseguiti & \hyperlink{UC17.2}{UC17.2}

\hyperlink{UC17}{UC17}

\hyperlink{UC18}{UC18}\tabularnewline
\hline
\begin{tikzpicture}
\draw [->, thick] (0.4,0.2) -- (0.4,0.1) -- (1,0.1);
\end{tikzpicture} & \hypertarget{R-3F26.3.2}{R-3F26.3.2} & Funzionale

Obbligatorio & Il sistema deve mostrare allo studente le statistiche di ogni singola domanda del questionario selezionato & \hyperlink{UC17.2}{UC17.2}

\hyperlink{UC17}{UC17}\tabularnewline
\hline
\begin{tikzpicture}
\draw [->, thick] (0.4,0.2) -- (0.4,0.1) -- (1,0.1);
\end{tikzpicture} & \hypertarget{R-3F26.3.3}{R-3F26.3.3} & Funzionale

Obbligatorio & Il sistema deve mostrare allo studente le proprie risposte al questionario selezionato e il punteggio totale & \hyperlink{UC17.2}{UC17.2}

\hyperlink{UC17}{UC17}\tabularnewline
\hline
\begin{tikzpicture}
\draw [->, thick] (0.4,0.2) -- (0.4,0.1) -- (1,0.1);
\end{tikzpicture} & \hypertarget{R-3F26.3.4}{R-3F26.3.4} & Funzionale

Obbligatorio & Il sistema deve mostrare allo studente la difficoltà del questionario calcolata come funzione delle risposte date da tutti gli utenti & \hyperlink{UC17.2}{UC17.2}

\hyperlink{UC17}{UC17}\tabularnewline
\hline
 & \hypertarget{R-3Q27}{R-3Q27} & Qualità

Obbligatorio & La progettazione rispetta le norme e le
metriche riportate nei documenti \NdP\ \textit{v2.0.0} e \PdQ\ \textit{v2.0.0} & Interno\tabularnewline
\hline
 & \hypertarget{R-3Q28}{R-3Q28} & Qualità

Obbligatorio & Devono essere rispettate tutte le norme e le metriche sulla stesura di codice riportate nei documenti \NdP\ \textit{v2.0.0} e \PdQ\ \textit{v2.0.0} & Interno\tabularnewline
\hline
 & \hypertarget{R-3Q29}{R-3Q29} & Qualità

Obbligatorio & Deve essere realizzato un manualeutente che descriva l’utilizzo del software & Interno\tabularnewline
\hline
 & \hypertarget{R-2Q30}{R-2Q30} & Qualità

Opzionale & Deve essere prodotta documentazione in Javadoc-like del codice sorgente del software scritto in Javascript & Interno\tabularnewline
\hline
\caption{Tabella requisiti / fonti} \tabularnewline
\end{longtable}


\subsection{Tabella fonti/requisiti}

% Output di tracy sotto a:
% tracy-tpiga.rhcloud.com/tracy/latex/SourceRequirement?type=latex&raw=1
\begin{longtable}[H]{p{0.1\textwidth} p{0.1\textwidth} p{0.6\textwidth} p{0.2\textwidth}}\\
	\toprule
	&	&	\textbf{Fonte}	&	\textbf{Requisito}\\*
	\midrule
	\midrule
	&  & Capitolato & \hyperlink{R-3V1}{R-3V1}
	
	\hyperlink{R-3V2}{R-3V2}
	
	\hyperlink{R-3V3}{R-3V3}
	
	\hyperlink{R-3V3.1}{R-3V3.1}
	
	\hyperlink{R-3V3.2}{R-3V3.2}
	
	\hyperlink{R-3V4}{R-3V4}
	
	\hyperlink{R-3V3.3}{R-3V3.3}
	
	\hyperlink{R-3V5}{R-3V5}
	
	\hyperlink{R-3V6}{R-3V6}
	
	\hyperlink{R-3V6.1}{R-3V6.1}
	
	\hyperlink{R-3V6.2}{R-3V6.2}
	
	\hyperlink{R-3F7.1}{R-3F7.1}
	
	\hyperlink{R-3F7.5.1}{R-3F7.5.1}
	
	\hyperlink{R-3F7.2}{R-3F7.2}
	
	\hyperlink{R-3F7.3}{R-3F7.3}
	
	\hyperlink{R-2F7.5.2}{R-2F7.5.2}
	
	\hyperlink{R-3F7.4}{R-3F7.4}
	
	\hyperlink{R-3F7.5}{R-3F7.5}
	
	\hyperlink{R-3F7}{R-3F7}
	
	\hyperlink{R-3F7.7}{R-3F7.7}
	
	\hyperlink{R-3F8}{R-3F8}
	
	\hyperlink{R-3V10}{R-3V10}
	
	\hyperlink{R-3V3.4}{R-3V3.4}
	
	\hyperlink{R-2F7.8}{R-2F7.8}
	
	\hyperlink{R-2F7.9}{R-2F7.9}
	
	\hyperlink{R-2F7.10}{R-2F7.10}
	
	\hyperlink{R-2F7.10.1}{R-2F7.10.1}
	
	\hyperlink{R-2F7.10.2}{R-2F7.10.2}
	
	\hyperlink{R-2F7.10.3}{R-2F7.10.3}
	
	\hyperlink{R-3F7.5.5}{R-3F7.5.5}
	
	\hyperlink{R-3F7.5.1.2}{R-3F7.5.1.2}\\*
	\midrule
	&  & Proponente & \hyperlink{R-3F7.5.3}{R-3F7.5.3}
	
	\hyperlink{R-2F7.5.4}{R-2F7.5.4}
	
	\hyperlink{R-2F7.6}{R-2F7.6}\\*
	\midrule
	&  & Committente & \\*
	\midrule
	&  & Interno & \hyperlink{R-3V10}{R-3V10}
	
	\hyperlink{R-3F9.1}{R-3F9.1}
	
	\hyperlink{R-3F18.1}{R-3F18.1}
	
	\hyperlink{R-3Q25}{R-3Q25}
	
	\hyperlink{R-3Q26}{R-3Q26}
	
	\hyperlink{R-3Q27}{R-3Q27}
	
	\hyperlink{R-2Q28}{R-2Q28}\\*
	\midrule
	& \hyperlink{UC1}{UC1} & \hyperlink{UC1}{Autenticazione} & \hyperlink{R-3F31}{R-3F31}\\*
	\midrule
	\begin{tikzpicture}
	\draw [->, thick] (0.2,0.2) -- (0.2,0.1) -- (1,0.1);
	\end{tikzpicture} & \hyperlink{UC1.1}{UC1.1} & \hyperlink{UC1.1}{Inserimento username} & \hyperlink{R-3F31.1}{R-3F31.1}\\*
	\midrule
	\begin{tikzpicture}
	\draw [->, thick] (0.2,0.2) -- (0.2,0.1) -- (1,0.1);
	\end{tikzpicture} & \hyperlink{UC1.2}{UC1.2} & \hyperlink{UC1.2}{Inserimento password} & \hyperlink{R-3F31.1}{R-3F31.1}\\*
	\midrule
	\begin{tikzpicture}
	\draw [->, thick] (0.2,0.2) -- (0.2,0.1) -- (1,0.1);
	\end{tikzpicture} & \hyperlink{UC1.3}{UC1.3} & \hyperlink{UC1.3}{Errore username non presente} & \hyperlink{R-3F31.1.1}{R-3F31.1.1}\\*
	\midrule
	\begin{tikzpicture}
	\draw [->, thick] (0.2,0.2) -- (0.2,0.1) -- (1,0.1);
	\end{tikzpicture} & \hyperlink{UC1.4}{UC1.4} & \hyperlink{UC1.4}{Errore password errata} & \hyperlink{R-3F31.1.2}{R-3F31.1.2}\\*
	\midrule
	& \hyperlink{UC2}{UC2} & \hyperlink{UC2}{Registrazione} & \hyperlink{R-3F9}{R-3F9}\\*
	\midrule
	\begin{tikzpicture}
	\draw [->, thick] (0.2,0.2) -- (0.2,0.1) -- (1,0.1);
	\end{tikzpicture} & \hyperlink{UC2.1}{UC2.1} & \hyperlink{UC2.1}{Inserimento nome completo} & \hyperlink{R-3F9.2}{R-3F9.2}
	
	\hyperlink{R-3F9.2.2}{R-3F9.2.2}\\*
	\midrule
	\begin{tikzpicture}
	\draw [->, thick] (0.2,0.2) -- (0.2,0.1) -- (1,0.1);
	\end{tikzpicture} & \hyperlink{UC2.2}{UC2.2} & \hyperlink{UC2.2}{Inserimento username} & \hyperlink{R-3F9.2}{R-3F9.2}\\*
	\midrule
	\begin{tikzpicture}
	\draw [->, thick] (0.2,0.2) -- (0.2,0.1) -- (1,0.1);
	\end{tikzpicture} & \hyperlink{UC2.3}{UC2.3} & \hyperlink{UC2.3}{Inserimento password} & \hyperlink{R-3F9.2}{R-3F9.2}\\*
	\midrule
	\begin{tikzpicture}
	\draw [->, thick] (0.2,0.2) -- (0.2,0.1) -- (1,0.1);
	\end{tikzpicture} & \hyperlink{UC2.4}{UC2.4} & \hyperlink{UC2.4}{Errore nome completo troppo corto} & \hyperlink{R-3F9.2.1}{R-3F9.2.1}\\*
	\midrule
	\begin{tikzpicture}
	\draw [->, thick] (0.2,0.2) -- (0.2,0.1) -- (1,0.1);
	\end{tikzpicture} & \hyperlink{UC2.5}{UC2.5} & \hyperlink{UC2.5}{Errore username troppo corto} & \hyperlink{R-3F9.2.1}{R-3F9.2.1}\\*
	\midrule
	\begin{tikzpicture}
	\draw [->, thick] (0.2,0.2) -- (0.2,0.1) -- (1,0.1);
	\end{tikzpicture} & \hyperlink{UC2.6}{UC2.6} & \hyperlink{UC2.6}{Errore password troppo corta} & \hyperlink{R-3F9.4}{R-3F9.4}
	
	\hyperlink{R-3F9.2.1}{R-3F9.2.1}\\*
	\midrule
	\begin{tikzpicture}
	\draw [->, thick] (0.2,0.2) -- (0.2,0.1) -- (1,0.1);
	\end{tikzpicture} & \hyperlink{UC2.7}{UC2.7} & \hyperlink{UC2.7}{Errore username già utilizzato} & \hyperlink{R-3F9.3}{R-3F9.3}
	
	\hyperlink{R-3F9.5}{R-3F9.5}
	
	\hyperlink{R-3F9.2.1}{R-3F9.2.1}\\*
	\midrule
	& \hyperlink{UC4}{UC4} & \hyperlink{UC4}{Gestione profilo} & \hyperlink{R-3F13}{R-3F13}
	
	\hyperlink{R-3F13.1}{R-3F13.1}
	
	\hyperlink{R-3F13.2}{R-3F13.2}
	
	\hyperlink{R-3F13.1.1}{R-3F13.1.1}
	
	\hyperlink{R-3F13.2.1}{R-3F13.2.1}
	
	\hyperlink{R-3F13.3}{R-3F13.3}
	
	\hyperlink{R-3F13.3.1}{R-3F13.3.1}
	
	\hyperlink{R-3F13.3.2}{R-3F13.3.2}
	
	\hyperlink{R-3F13.2.1.1}{R-3F13.2.1.1}
	
	\hyperlink{R-3F13.2.2}{R-3F13.2.2}
	
	\hyperlink{R-3F13.2.2.1}{R-3F13.2.2.1}
	
	\hyperlink{R-3F13.3.3}{R-3F13.3.3}
	
	\hyperlink{R-3F13.1.2}{R-3F13.1.2}
	
	\hyperlink{R-3F13.3.4}{R-3F13.3.4}\\*
	\midrule
	\begin{tikzpicture}
	\draw [->, thick] (0.2,0.2) -- (0.2,0.1) -- (1,0.1);
	\end{tikzpicture} & \hyperlink{UC4.1}{UC4.1} & \hyperlink{UC4.1}{Modifica nome completo} & \hyperlink{R-3F13.3}{R-3F13.3}
	
	\hyperlink{R-3F13.3.1}{R-3F13.3.1}
	
	\hyperlink{R-3F13.3.3}{R-3F13.3.3}\\*
	\midrule
	\begin{tikzpicture}
	\draw [->, thick] (0.2,0.2) -- (0.2,0.1) -- (1,0.1);
	\end{tikzpicture} & \hyperlink{UC4.2}{UC4.2} & \hyperlink{UC4.2}{Modifica username} & \hyperlink{R-3F13.1}{R-3F13.1}
	
	\hyperlink{R-3F13.1.1}{R-3F13.1.1}
	
	\hyperlink{R-3F13.3}{R-3F13.3}
	
	\hyperlink{R-3F13.3.2}{R-3F13.3.2}
	
	\hyperlink{R-3F13.3.4}{R-3F13.3.4}\\*
	\midrule
	\begin{tikzpicture}
	\draw [->, thick] (0.2,0.2) -- (0.2,0.1) -- (1,0.1);
	\end{tikzpicture} & \hyperlink{UC4.3}{UC4.3} & \hyperlink{UC4.3}{Modifica password} & \hyperlink{R-3F13.2}{R-3F13.2}
	
	\hyperlink{R-3F13.2.1}{R-3F13.2.1}
	
	\hyperlink{R-3F13.2.2}{R-3F13.2.2}
	
	\hyperlink{R-3F13.2.2.1}{R-3F13.2.2.1}\\*
	\midrule
	\begin{tikzpicture}
	\draw [->, thick] (0.4,0.2) -- (0.4,0.1) -- (1,0.1);
	\end{tikzpicture} & \hyperlink{UC4.3.1}{UC4.3.1} & \hyperlink{UC4.3.1}{Inserimento vecchia password} & \hyperlink{R-3F13.2.1}{R-3F13.2.1}\\*
	\midrule
	\begin{tikzpicture}
	\draw [->, thick] (0.4,0.2) -- (0.4,0.1) -- (1,0.1);
	\end{tikzpicture} & \hyperlink{UC4.3.2}{UC4.3.2} & \hyperlink{UC4.3.2}{Errore password non corrispondente} & \hyperlink{R-3F13.2.1.1}{R-3F13.2.1.1}\\*
	\midrule
	\begin{tikzpicture}
	\draw [->, thick] (0.4,0.2) -- (0.4,0.1) -- (1,0.1);
	\end{tikzpicture} & \hyperlink{UC4.3.3}{UC4.3.3} & \hyperlink{UC4.3.3}{Inserimento nuova password} & \hyperlink{R-3F13.2}{R-3F13.2}\\*
	\midrule
	\begin{tikzpicture}
	\draw [->, thick] (0.2,0.2) -- (0.2,0.1) -- (1,0.1);
	\end{tikzpicture} & \hyperlink{UC4.4}{UC4.4} & \hyperlink{UC4.4}{Errore username troppo corto} & \hyperlink{R-3F13.1.1}{R-3F13.1.1}
	
	\hyperlink{R-3F13.3.4}{R-3F13.3.4}\\*
	\midrule
	\begin{tikzpicture}
	\draw [->, thick] (0.2,0.2) -- (0.2,0.1) -- (1,0.1);
	\end{tikzpicture} & \hyperlink{UC4.5}{UC4.5} & \hyperlink{UC4.5}{Errore nome completo troppo corto} & \hyperlink{R-3F13.3.3}{R-3F13.3.3}\\*
	\midrule
	\begin{tikzpicture}
	\draw [->, thick] (0.2,0.2) -- (0.2,0.1) -- (1,0.1);
	\end{tikzpicture} & \hyperlink{UC4.6}{UC4.6} & \hyperlink{UC4.6}{Errore username già utilizzato} & \hyperlink{R-3F13.1.1}{R-3F13.1.1}
	
	\hyperlink{R-3F13.1.2}{R-3F13.1.2}\\*
	\midrule
	\begin{tikzpicture}
	\draw [->, thick] (0.2,0.2) -- (0.2,0.1) -- (1,0.1);
	\end{tikzpicture} & \hyperlink{UC4.7}{UC4.7} & \hyperlink{UC4.7}{Errore password troppo corta} & \hyperlink{R-3F13.2.1.1}{R-3F13.2.1.1}
	
	\hyperlink{R-3F13.2.2.1}{R-3F13.2.2.1}\\*
	\midrule
	& \hyperlink{UC5}{UC5} & \hyperlink{UC5}{Logout} & \hyperlink{R-3F17}{R-3F17}\\*
	\midrule
	& \hyperlink{UC18}{UC18} & \hyperlink{UC18}{Gestione domande} & \hyperlink{R-3F7.2}{R-3F7.2}
	
	\hyperlink{R-3F7.11}{R-3F7.11}
	
	\hyperlink{R-3F7.11.1}{R-3F7.11.1}
	
	\hyperlink{R-3F7.11.2}{R-3F7.11.2}
	
	\hyperlink{R-3F7.11.3}{R-3F7.11.3}
	
	\hyperlink{R-3F7.11.1.1}{R-3F7.11.1.1}
	
	\hyperlink{R-3F7.11.1.2}{R-3F7.11.1.2}
	
	\hyperlink{R-3F7.11.1.1.1}{R-3F7.11.1.1.1}
	
	\hyperlink{R-3F7.11.1.2.1}{R-3F7.11.1.2.1}\\*
	\midrule
	\begin{tikzpicture}
	\draw [->, thick] (0.2,0.2) -- (0.2,0.1) -- (1,0.1);
	\end{tikzpicture} & \hyperlink{UC18.1}{UC18.1} & \hyperlink{UC18.1}{Inserimento domanda} & \hyperlink{R-2F7.5.2}{R-2F7.5.2}
	
	\hyperlink{R-3F7.5.3}{R-3F7.5.3}
	
	\hyperlink{R-3F7.11.1}{R-3F7.11.1}
	
	\hyperlink{R-3F7.11.1.1}{R-3F7.11.1.1}
	
	\hyperlink{R-3F7.11.1.2.2}{R-3F7.11.1.2.2}
	
	\hyperlink{R-3F7.11.1.2.3}{R-3F7.11.1.2.3}
	
	\hyperlink{R-1F7.11.1.2.4}{R-1F7.11.1.2.4}
	
	\hyperlink{R-1F7.11.1.2.5}{R-1F7.11.1.2.5}
	
	\hyperlink{R-1F7.11.1.2.6}{R-1F7.11.1.2.6}
	
	\hyperlink{R-2F7.11.1.2.7}{R-2F7.11.1.2.7}
	
	\hyperlink{R-3F7.5.5}{R-3F7.5.5}\\*
	\midrule
	\begin{tikzpicture}
	\draw [->, thick] (0.4,0.2) -- (0.4,0.1) -- (1,0.1);
	\end{tikzpicture} & \hyperlink{UC18.1.1}{UC18.1.1} & \hyperlink{UC18.1.1}{Selezione argomenti nuova domanda} & \hyperlink{R-3F7.11.1.1}{R-3F7.11.1.1}
	
	\hyperlink{R-3F7.11.1.1.1}{R-3F7.11.1.1.1}\\*
	\midrule
	\begin{tikzpicture}
	\draw [->, thick] (0.4,0.2) -- (0.4,0.1) -- (1,0.1);
	\end{tikzpicture} & \hyperlink{UC18.1.2}{UC18.1.2} & \hyperlink{UC18.1.2}{Scrittura domanda in QML della nuova domanda} & \hyperlink{R-3F7.5.1}{R-3F7.5.1}
	
	\hyperlink{R-3F7.5}{R-3F7.5}
	
	\hyperlink{R-3F7.11.1.2}{R-3F7.11.1.2}
	
	\hyperlink{R-3F7.11.1.2.1}{R-3F7.11.1.2.1}
	
	\hyperlink{R-2F7.5.1.1}{R-2F7.5.1.1}
	
	\hyperlink{R-3F7.5.1.2}{R-3F7.5.1.2}
	
	\hyperlink{R-2F7.5.1.3}{R-2F7.5.1.3}
	
	\hyperlink{R-2F7.5.1.4}{R-2F7.5.1.4}
	
	\hyperlink{R-2F7.5.1.5}{R-2F7.5.1.5}
	
	\hyperlink{R-3F7.5.1.6}{R-3F7.5.1.6}
	
	\hyperlink{R-2F7.5.1.7}{R-2F7.5.1.7}
	
	\hyperlink{R-2F7.5.1.8}{R-2F7.5.1.8}\\*
	\midrule
	\begin{tikzpicture}
	\draw [->, thick] (0.4,0.2) -- (0.4,0.1) -- (1,0.1);
	\end{tikzpicture} & \hyperlink{UC18.1.3}{UC18.1.3} & \hyperlink{UC18.1.3}{Errore QML non valido} & \hyperlink{R-3F7.11.1.2.1}{R-3F7.11.1.2.1}\\*
	\midrule
	\begin{tikzpicture}
	\draw [->, thick] (0.4,0.2) -- (0.4,0.1) -- (1,0.1);
	\end{tikzpicture} & \hyperlink{UC18.1.4}{UC18.1.4} & \hyperlink{UC18.1.4}{Errore argomento mancante} & \hyperlink{R-3F7.11.1.1.1}{R-3F7.11.1.1.1}\\*
	\midrule
	\begin{tikzpicture}
	\draw [->, thick] (0.4,0.2) -- (0.4,0.1) -- (1,0.1);
	\end{tikzpicture} & \hyperlink{UC18.1.5}{UC18.1.5} & \hyperlink{UC18.1.5}{Scrittura nuova domanda da interfaccia grafica} & \hyperlink{R-2F7.5.4}{R-2F7.5.4}
	
	\hyperlink{R-2F7.5.1.7}{R-2F7.5.1.7}
	
	\hyperlink{R-2F7.5.1.8}{R-2F7.5.1.8}\\*
	\midrule
	\begin{tikzpicture}
	\draw [->, thick] (0.2,0.2) -- (0.2,0.1) -- (1,0.1);
	\end{tikzpicture} & \hyperlink{UC18.2}{UC18.2} & \hyperlink{UC18.2}{Modifica domanda} & \hyperlink{R-3F7.11.2}{R-3F7.11.2}\\*
	\midrule
	\begin{tikzpicture}
	\draw [->, thick] (0.4,0.2) -- (0.4,0.1) -- (1,0.1);
	\end{tikzpicture} & \hyperlink{UC18.2.1}{UC18.2.1} & \hyperlink{UC18.2.1}{Selezione argomenti modifica domanda} & \hyperlink{R-3F7.11.2}{R-3F7.11.2}\\*
	\midrule
	\begin{tikzpicture}
	\draw [->, thick] (0.4,0.2) -- (0.4,0.1) -- (1,0.1);
	\end{tikzpicture} & \hyperlink{UC18.2.2}{UC18.2.2} & \hyperlink{UC18.2.2}{Scrittura domanda in QML della domanda da modificare} & \hyperlink{R-3F7.11.2}{R-3F7.11.2}\\*
	\midrule
	\begin{tikzpicture}
	\draw [->, thick] (0.4,0.2) -- (0.4,0.1) -- (1,0.1);
	\end{tikzpicture} & \hyperlink{UC18.2.3}{UC18.2.3} & \hyperlink{UC18.2.3}{Modifica domanda da interfaccia grafica} & \hyperlink{R-3F7.11.2}{R-3F7.11.2}\\*
	\midrule
	\begin{tikzpicture}
	\draw [->, thick] (0.4,0.2) -- (0.4,0.1) -- (1,0.1);
	\end{tikzpicture} & \hyperlink{UC18.2.4}{UC18.2.4} & \hyperlink{UC18.2.4}{Errore modifica QML non valido} & \hyperlink{R-3F7.11.1.2.1}{R-3F7.11.1.2.1}\\*
	\midrule
	\begin{tikzpicture}
	\draw [->, thick] (0.4,0.2) -- (0.4,0.1) -- (1,0.1);
	\end{tikzpicture} & \hyperlink{UC18.2.5}{UC18.2.5} & \hyperlink{UC18.2.5}{Errore argomento mancante} & \hyperlink{R-3F7.11.1.1.2}{R-3F7.11.1.1.2}\\*
	\midrule
	\begin{tikzpicture}
	\draw [->, thick] (0.2,0.2) -- (0.2,0.1) -- (1,0.1);
	\end{tikzpicture} & \hyperlink{UC18.3}{UC18.3} & \hyperlink{UC18.3}{Elimina domanda} & \hyperlink{R-3F7.11.3}{R-3F7.11.3}\\*
	\midrule
	\begin{tikzpicture}
	\draw [->, thick] (0.2,0.2) -- (0.2,0.1) -- (1,0.1);
	\end{tikzpicture} & \hyperlink{UC18.4}{UC18.4} & \hyperlink{UC18.4}{Inserimento domanda di tipo vero/falso} & \hyperlink{R-3F7.11.1.2.2}{R-3F7.11.1.2.2}
	
	\hyperlink{R-3F7.5.5}{R-3F7.5.5}\\*
	\midrule
	\begin{tikzpicture}
	\draw [->, thick] (0.2,0.2) -- (0.2,0.1) -- (1,0.1);
	\end{tikzpicture} & \hyperlink{UC18.5}{UC18.5} & \hyperlink{UC18.5}{Inserimento domanda a scelta multipla} & \hyperlink{R-2F7.5.2}{R-2F7.5.2}
	
	\hyperlink{R-3F7.11.1.2.3}{R-3F7.11.1.2.3}\\*
	\midrule
	\begin{tikzpicture}
	\draw [->, thick] (0.2,0.2) -- (0.2,0.1) -- (1,0.1);
	\end{tikzpicture} & \hyperlink{UC18.6}{UC18.6} & \hyperlink{UC18.6}{Inserimento domanda a risposta multipla} & \hyperlink{R-1F7.11.1.2.4}{R-1F7.11.1.2.4}
	
	\hyperlink{R-3F7.5.5}{R-3F7.5.5}\\*
	\midrule
	\begin{tikzpicture}
	\draw [->, thick] (0.2,0.2) -- (0.2,0.1) -- (1,0.1);
	\end{tikzpicture} & \hyperlink{UC18.7}{UC18.7} & \hyperlink{UC18.7}{Inserimento domanda di tipo testo con parole omesse} & \hyperlink{R-2F7.5.2}{R-2F7.5.2}
	
	\hyperlink{R-1F7.11.1.2.5}{R-1F7.11.1.2.5}\\*
	\midrule
	\begin{tikzpicture}
	\draw [->, thick] (0.2,0.2) -- (0.2,0.1) -- (1,0.1);
	\end{tikzpicture} & \hyperlink{UC18.8}{UC18.8} & \hyperlink{UC18.8}{Inserimento domanda con l'associazione di parole} & \hyperlink{R-2F7.5.2}{R-2F7.5.2}
	
	\hyperlink{R-1F7.11.1.2.6}{R-1F7.11.1.2.6}\\*
	\midrule
	\begin{tikzpicture}
	\draw [->, thick] (0.2,0.2) -- (0.2,0.1) -- (1,0.1);
	\end{tikzpicture} & \hyperlink{UC18.9}{UC18.9} & \hyperlink{UC18.9}{Inserimento domanda a risposta aperta} & \hyperlink{R-2F7.5.2}{R-2F7.5.2}
	
	\hyperlink{R-2F7.11.1.2.7}{R-2F7.11.1.2.7}\\*
	\midrule
	\begin{tikzpicture}
	\draw [->, thick] (0.2,0.2) -- (0.2,0.1) -- (1,0.1);
	\end{tikzpicture} & \hyperlink{UC18.10}{UC18.10} & \hyperlink{UC18.10}{Visualizza domanda} & \hyperlink{R-3F7.2}{R-3F7.2}
	
	\hyperlink{R-3F29}{R-3F29}\\*
	\midrule
	\begin{tikzpicture}
	\draw [->, thick] (0.2,0.2) -- (0.2,0.1) -- (1,0.1);
	\end{tikzpicture} & \hyperlink{UC18.11}{UC18.11} & \hyperlink{UC18.11}{Errore eliminazione domanda utilizzata} & \hyperlink{R-3F7.11.3.1}{R-3F7.11.3.1}\\*
	\midrule
	& \hyperlink{UC19}{UC19} & \hyperlink{UC19}{Gestione questionari} & \hyperlink{R-3F7.2}{R-3F7.2}
	
	\hyperlink{R-3F7.3}{R-3F7.3}
	
	\hyperlink{R-3F7}{R-3F7}
	
	\hyperlink{R-3F7.7}{R-3F7.7}
	
	\hyperlink{R-3F7.11}{R-3F7.11}
	
	\hyperlink{R-3F7.12}{R-3F7.12}
	
	\hyperlink{R-3F7.12.2}{R-3F7.12.2}
	
	\hyperlink{R-3F7.12.3}{R-3F7.12.3}
	
	\hyperlink{R-3F7.13}{R-3F7.13}
	
	\hyperlink{R-3F7.12.4}{R-3F7.12.4}
	
	\hyperlink{R-3F7.7.1}{R-3F7.7.1}\\*
	\midrule
	\begin{tikzpicture}
	\draw [->, thick] (0.2,0.2) -- (0.2,0.1) -- (1,0.1);
	\end{tikzpicture} & \hyperlink{UC19.1}{UC19.1} & \hyperlink{UC19.1}{Inserisci questionario} & \hyperlink{R-3F7.7}{R-3F7.7}
	
	\hyperlink{R-3F7.7.1}{R-3F7.7.1}
	
	\hyperlink{R-3F7.14}{R-3F7.14}\\*
	\midrule
	\begin{tikzpicture}
	\draw [->, thick] (0.4,0.2) -- (0.4,0.1) -- (1,0.1);
	\end{tikzpicture} & \hyperlink{UC19.1.1}{UC19.1.1} & \hyperlink{UC19.1.1}{Aggiungi domanda in un nuovo questionario } & \hyperlink{R-3F7.11}{R-3F7.11}
	
	\hyperlink{R-3F7.11.1.2}{R-3F7.11.1.2}
	
	\hyperlink{R-3F7.11.1.1.1}{R-3F7.11.1.1.1}
	
	\hyperlink{R-3F7.11.1.2.1}{R-3F7.11.1.2.1}\\*
	\midrule
	\begin{tikzpicture}
	\draw [->, thick] (0.4,0.2) -- (0.4,0.1) -- (1,0.1);
	\end{tikzpicture} & \hyperlink{UC19.1.2}{UC19.1.2} & \hyperlink{UC19.1.2}{Elimina domanda da un nuovo questionario} & \hyperlink{R-3F7.11}{R-3F7.11}\\*
	\midrule
	\begin{tikzpicture}
	\draw [->, thick] (0.4,0.2) -- (0.4,0.1) -- (1,0.1);
	\end{tikzpicture} & \hyperlink{UC19.1.3}{UC19.1.3} & \hyperlink{UC19.1.3}{Seleziona argomenti del nuovo questionario} & \hyperlink{R-3F7.14.1}{R-3F7.14.1}\\*
	\midrule
	\begin{tikzpicture}
	\draw [->, thick] (0.2,0.2) -- (0.2,0.1) -- (1,0.1);
	\end{tikzpicture} & \hyperlink{UC19.2}{UC19.2} & \hyperlink{UC19.2}{Modifica questionario} & \hyperlink{R-3F7.12}{R-3F7.12}
	
	\hyperlink{R-3F7.12.4}{R-3F7.12.4}\\*
	\midrule
	\begin{tikzpicture}
	\draw [->, thick] (0.4,0.2) -- (0.4,0.1) -- (1,0.1);
	\end{tikzpicture} & \hyperlink{UC19.2.1}{UC19.2.1} & \hyperlink{UC19.2.1}{Aggiungi domanda in un questionario} & \hyperlink{R-3F7.12.2}{R-3F7.12.2}\\*
	\midrule
	\begin{tikzpicture}
	\draw [->, thick] (0.4,0.2) -- (0.4,0.1) -- (1,0.1);
	\end{tikzpicture} & \hyperlink{UC19.2.2}{UC19.2.2} & \hyperlink{UC19.2.2}{Elimina domanda da un questionario da  modificare} & \hyperlink{R-3F7.12.3}{R-3F7.12.3}\\*
	\midrule
	\begin{tikzpicture}
	\draw [->, thick] (0.4,0.2) -- (0.4,0.1) -- (1,0.1);
	\end{tikzpicture} & \hyperlink{UC19.2.3}{UC19.2.3} & \hyperlink{UC19.2.3}{Selezione argomenti modifica questionario} & \hyperlink{R-3F7.12.4}{R-3F7.12.4}\\*
	\midrule
	\begin{tikzpicture}
	\draw [->, thick] (0.2,0.2) -- (0.2,0.1) -- (1,0.1);
	\end{tikzpicture} & \hyperlink{UC19.3}{UC19.3} & \hyperlink{UC19.3}{Elimina questionario} & \hyperlink{R-3F7.13}{R-3F7.13}\\*
	\midrule
	\begin{tikzpicture}
	\draw [->, thick] (0.2,0.2) -- (0.2,0.1) -- (1,0.1);
	\end{tikzpicture} & \hyperlink{UC19.4}{UC19.4} & \hyperlink{UC19.4}{Errore questionario vuoto} & \hyperlink{R-3F7.7.1}{R-3F7.7.1}\\*
	\midrule
	\begin{tikzpicture}
	\draw [->, thick] (0.2,0.2) -- (0.2,0.1) -- (1,0.1);
	\end{tikzpicture} & \hyperlink{UC19.5}{UC19.5} & \hyperlink{UC19.5}{Visualizza questionario} & \hyperlink{R-3F30}{R-3F30}\\*
	\midrule
	& \hyperlink{UC20}{UC20} & \hyperlink{UC20}{Gestione classi} & \hyperlink{R-2F12.1.1}{R-2F12.1.1}
	
	\hyperlink{R-2F12}{R-2F12}
	
	\hyperlink{R-2F12.1}{R-2F12.1}
	
	\hyperlink{R-2F12.2}{R-2F12.2}
	
	\hyperlink{R-2F12.3}{R-2F12.3}
	
	\hyperlink{R-2F12.1.2}{R-2F12.1.2}
	
	\hyperlink{R-2F12.1.3}{R-2F12.1.3}
	
	\hyperlink{R-2F12.3.1}{R-2F12.3.1}
	
	\hyperlink{R-2F12.3.2}{R-2F12.3.2}
	
	\hyperlink{R-2F12.3.3}{R-2F12.3.3}
	
	\hyperlink{R-2F12.1.1.1}{R-2F12.1.1.1}\\*
	\midrule
	\begin{tikzpicture}
	\draw [->, thick] (0.2,0.2) -- (0.2,0.1) -- (1,0.1);
	\end{tikzpicture} & \hyperlink{UC20.1}{UC20.1} & \hyperlink{UC20.1}{Inserisci classe} & \hyperlink{R-2F7.6}{R-2F7.6}
	
	\hyperlink{R-2F12.1.1}{R-2F12.1.1}
	
	\hyperlink{R-2F12.1}{R-2F12.1}
	
	\hyperlink{R-2F12.1.2}{R-2F12.1.2}
	
	\hyperlink{R-2F12.1.3}{R-2F12.1.3}
	
	\hyperlink{R-2F12.1.1.1}{R-2F12.1.1.1}\\*
	\midrule
	\begin{tikzpicture}
	\draw [->, thick] (0.4,0.2) -- (0.4,0.1) -- (1,0.1);
	\end{tikzpicture} & \hyperlink{UC20.1.1}{UC20.1.1} & \hyperlink{UC20.1.1}{Inserisci nome classe} & \hyperlink{R-2F12.1.1}{R-2F12.1.1}
	
	\hyperlink{R-2F12.3.1}{R-2F12.3.1}
	
	\hyperlink{R-2F12.1.1.1}{R-2F12.1.1.1}\\*
	\midrule
	\begin{tikzpicture}
	\draw [->, thick] (0.4,0.2) -- (0.4,0.1) -- (1,0.1);
	\end{tikzpicture} & \hyperlink{UC20.1.2}{UC20.1.2} & \hyperlink{UC20.1.2}{Inserisci argomenti classe} & \hyperlink{R-2F12.1.2}{R-2F12.1.2}
	
	\hyperlink{R-2F12.3.2}{R-2F12.3.2}\\*
	\midrule
	\begin{tikzpicture}
	\draw [->, thick] (0.4,0.2) -- (0.4,0.1) -- (1,0.1);
	\end{tikzpicture} & \hyperlink{UC20.1.3}{UC20.1.3} & \hyperlink{UC20.1.3}{Inserisci password classe} & \hyperlink{R-2F12.1.3}{R-2F12.1.3}
	
	\hyperlink{R-2F12.3.3}{R-2F12.3.3}\\*
	\midrule
	\begin{tikzpicture}
	\draw [->, thick] (0.2,0.2) -- (0.2,0.1) -- (1,0.1);
	\end{tikzpicture} & \hyperlink{UC20.2}{UC20.2} & \hyperlink{UC20.2}{Modifica classe} & \hyperlink{R-2F12.3}{R-2F12.3}
	
	\hyperlink{R-2F12.3.1}{R-2F12.3.1}
	
	\hyperlink{R-2F12.3.2}{R-2F12.3.2}
	
	\hyperlink{R-2F12.3.3}{R-2F12.3.3}\\*
	\midrule
	\begin{tikzpicture}
	\draw [->, thick] (0.2,0.2) -- (0.2,0.1) -- (1,0.1);
	\end{tikzpicture} & \hyperlink{UC20.3}{UC20.3} & \hyperlink{UC20.3}{Elimina classe} & \hyperlink{R-2F12.2}{R-2F12.2}\\*
	\midrule
	\begin{tikzpicture}
	\draw [->, thick] (0.2,0.2) -- (0.2,0.1) -- (1,0.1);
	\end{tikzpicture} & \hyperlink{UC20.4}{UC20.4} & \hyperlink{UC20.4}{Errore nome classe già presente} & \hyperlink{R-2F12.1.1.1}{R-2F12.1.1.1}\\*
	\midrule
	& \hyperlink{UC21}{UC21} & \hyperlink{UC21}{Gestione argomenti} & \hyperlink{R-3F22}{R-3F22}
	
	\hyperlink{R-3F22.1}{R-3F22.1}
	
	\hyperlink{R-3F22.2}{R-3F22.2}
	
	\hyperlink{R-3F22.3}{R-3F22.3}
	
	\hyperlink{R-3F22.4}{R-3F22.4}
	
	\hyperlink{R-3F22.1.1}{R-3F22.1.1}
	
	\hyperlink{R-3F22.2.1}{R-3F22.2.1}
	
	\hyperlink{R-3F22.2.2}{R-3F22.2.2}\\*
	\midrule
	\begin{tikzpicture}
	\draw [->, thick] (0.2,0.2) -- (0.2,0.1) -- (1,0.1);
	\end{tikzpicture} & \hyperlink{UC21.1}{UC21.1} & \hyperlink{UC21.1}{Crea argomento} & \hyperlink{R-3F22.1}{R-3F22.1}
	
	\hyperlink{R-3F22.1.1}{R-3F22.1.1}\\*
	\midrule
	\begin{tikzpicture}
	\draw [->, thick] (0.2,0.2) -- (0.2,0.1) -- (1,0.1);
	\end{tikzpicture} & \hyperlink{UC21.2}{UC21.2} & \hyperlink{UC21.2}{Errore argomento già presente nel sistema} & \hyperlink{R-3F22.1.1}{R-3F22.1.1}\\*
	\midrule
	\begin{tikzpicture}
	\draw [->, thick] (0.2,0.2) -- (0.2,0.1) -- (1,0.1);
	\end{tikzpicture} & \hyperlink{UC21.3}{UC21.3} & \hyperlink{UC21.3}{Modifica argomento} & \hyperlink{R-3F22.3}{R-3F22.3}\\*
	\midrule
	\begin{tikzpicture}
	\draw [->, thick] (0.2,0.2) -- (0.2,0.1) -- (1,0.1);
	\end{tikzpicture} & \hyperlink{UC21.4}{UC21.4} & \hyperlink{UC21.4}{Eliminazione argomento} & \hyperlink{R-3F22.2}{R-3F22.2}
	
	\hyperlink{R-3F22.2.1}{R-3F22.2.1}
	
	\hyperlink{R-3F22.2.2}{R-3F22.2.2}\\*
	\midrule
	\begin{tikzpicture}
	\draw [->, thick] (0.2,0.2) -- (0.2,0.1) -- (1,0.1);
	\end{tikzpicture} & \hyperlink{UC21.5}{UC21.5} & \hyperlink{UC21.5}{Errore l'argomento ha domande o questionari} & \hyperlink{R-3F22.2.2}{R-3F22.2.2}\\*
	\midrule
	& \hyperlink{UC22}{UC22} & \hyperlink{UC22}{Visualizza statistiche} & \hyperlink{R-2F23.3.1}{R-2F23.3.1}
	
	\hyperlink{R-2F23}{R-2F23}
	
	\hyperlink{R-2F23.1}{R-2F23.1}
	
	\hyperlink{R-2F23.2}{R-2F23.2}
	
	\hyperlink{R-2F23.3}{R-2F23.3}
	
	\hyperlink{R-2F23.1.1}{R-2F23.1.1}
	
	\hyperlink{R-2F23.2.1}{R-2F23.2.1}
	
	\hyperlink{R-2F23.3.2}{R-2F23.3.2}
	
	\hyperlink{R-2F23.3.3}{R-2F23.3.3}
	
	\hyperlink{R-2F23.3.4}{R-2F23.3.4}
	
	\hyperlink{R-2F23.3.4.1}{R-2F23.3.4.1}
	
	\hyperlink{R-2F23.3.5}{R-2F23.3.5}
	
	\hyperlink{R-2F23.3.5.1}{R-2F23.3.5.1}\\*
	\midrule
	& \hyperlink{UC23}{UC23} & \hyperlink{UC23}{Visualizza statistiche domanda} & \hyperlink{R-2F23.1}{R-2F23.1}
	
	\hyperlink{R-2F23.1.1}{R-2F23.1.1}\\*
	\midrule
	& \hyperlink{UC24}{UC24} & \hyperlink{UC24}{Visualizza statistiche questionario} & \hyperlink{R-2F23.2}{R-2F23.2}
	
	\hyperlink{R-2F23.2.1}{R-2F23.2.1}\\*
	\midrule
	& \hyperlink{UC25}{UC25} & \hyperlink{UC25}{Visualizza statistiche classe} & \hyperlink{R-2F23.3.1}{R-2F23.3.1}
	
	\hyperlink{R-2F23.3}{R-2F23.3}
	
	\hyperlink{R-2F23.3.2}{R-2F23.3.2}
	
	\hyperlink{R-2F23.3.3}{R-2F23.3.3}
	
	\hyperlink{R-2F23.3.4}{R-2F23.3.4}
	
	\hyperlink{R-2F23.3.4.1}{R-2F23.3.4.1}
	
	\hyperlink{R-2F23.3.5}{R-2F23.3.5}
	
	\hyperlink{R-2F23.3.5.1}{R-2F23.3.5.1}\\*
	\midrule
	\begin{tikzpicture}
	\draw [->, thick] (0.2,0.2) -- (0.2,0.1) -- (1,0.1);
	\end{tikzpicture} & \hyperlink{UC25.1}{UC25.1} & \hyperlink{UC25.1}{Visualizza risultati domande della classe} & \hyperlink{R-2F23.3.5}{R-2F23.3.5}
	
	\hyperlink{R-2F23.3.5.1}{R-2F23.3.5.1}\\*
	\midrule
	\begin{tikzpicture}
	\draw [->, thick] (0.2,0.2) -- (0.2,0.1) -- (1,0.1);
	\end{tikzpicture} & \hyperlink{UC25.2}{UC25.2} & \hyperlink{UC25.2}{Visualizza risultati questionari della classe} & \hyperlink{R-2F23.3.1}{R-2F23.3.1}\\*
	\midrule
	\begin{tikzpicture}
	\draw [->, thick] (0.2,0.2) -- (0.2,0.1) -- (1,0.1);
	\end{tikzpicture} & \hyperlink{UC25.3}{UC25.3} & \hyperlink{UC25.3}{Visualizza sommario statistiche classe} & \hyperlink{R-2F23.3.3}{R-2F23.3.3}\\*
	\midrule
	\begin{tikzpicture}
	\draw [->, thick] (0.2,0.2) -- (0.2,0.1) -- (1,0.1);
	\end{tikzpicture} & \hyperlink{UC25.4}{UC25.4} & \hyperlink{UC25.4}{Visualizza statistiche studente della classe} & \hyperlink{R-2F23.3.4}{R-2F23.3.4}
	
	\hyperlink{R-2F23.3.4.1}{R-2F23.3.4.1}\\*
	\midrule
	\begin{tikzpicture}
	\draw [->, thick] (0.4,0.2) -- (0.4,0.1) -- (1,0.1);
	\end{tikzpicture} & \hyperlink{UC25.4.1}{UC25.4.1} & \hyperlink{UC25.4.1}{Visualizza risultati questionari dello studente} & \hyperlink{R-2F23.3.4.1}{R-2F23.3.4.1}\\*
	\midrule
	& \hyperlink{UC6}{UC6} & \hyperlink{UC6}{Esegui questionario} & \hyperlink{R-3F7.4}{R-3F7.4}
	
	\hyperlink{R-3F16}{R-3F16}
	
	\hyperlink{R-3F16.1}{R-3F16.1}
	
	\hyperlink{R-3F16.2}{R-3F16.2}
	
	\hyperlink{R-3F16.3}{R-3F16.3}
	
	\hyperlink{R-3F16.4}{R-3F16.4}
	
	\hyperlink{R-3F16.2.1}{R-3F16.2.1}\\*
	\midrule
	\begin{tikzpicture}
	\draw [->, thick] (0.2,0.2) -- (0.2,0.1) -- (1,0.1);
	\end{tikzpicture} & \hyperlink{UC6.1}{UC6.1} & \hyperlink{UC6.1}{Rispondi domanda} & \hyperlink{R-3F7.2}{R-3F7.2}
	
	\hyperlink{R-3F16.1}{R-3F16.1}
	
	\hyperlink{R-3F16.1.1}{R-3F16.1.1}
	
	\hyperlink{R-3F16.1.2}{R-3F16.1.2}
	
	\hyperlink{R-1F16.1.3}{R-1F16.1.3}
	
	\hyperlink{R-1F16.1.4}{R-1F16.1.4}
	
	\hyperlink{R-1F16.1.5}{R-1F16.1.5}
	
	\hyperlink{R-2F16.1.6}{R-2F16.1.6}\\*
	\midrule
	\begin{tikzpicture}
	\draw [->, thick] (0.2,0.2) -- (0.2,0.1) -- (1,0.1);
	\end{tikzpicture} & \hyperlink{UC6.2}{UC6.2} & \hyperlink{UC6.2}{Errore domanda non risposta} & \hyperlink{R-3F16.2.1}{R-3F16.2.1}\\*
	\midrule
	\begin{tikzpicture}
	\draw [->, thick] (0.2,0.2) -- (0.2,0.1) -- (1,0.1);
	\end{tikzpicture} & \hyperlink{UC6.3}{UC6.3} & \hyperlink{UC6.3}{Conferma questionario} & \hyperlink{R-3F16.2}{R-3F16.2}
	
	\hyperlink{R-3F16.2.1}{R-3F16.2.1}\\*
	\midrule
	\begin{tikzpicture}
	\draw [->, thick] (0.2,0.2) -- (0.2,0.1) -- (1,0.1);
	\end{tikzpicture} & \hyperlink{UC6.4}{UC6.4} & \hyperlink{UC6.4}{Rispondi domanda vero/falso} & \hyperlink{R-3F16.1.1}{R-3F16.1.1}\\*
	\midrule
	\begin{tikzpicture}
	\draw [->, thick] (0.2,0.2) -- (0.2,0.1) -- (1,0.1);
	\end{tikzpicture} & \hyperlink{UC6.5}{UC6.5} & \hyperlink{UC6.5}{Rispondi domanda a scelta multipla} & \hyperlink{R-3F16.1.2}{R-3F16.1.2}\\*
	\midrule
	\begin{tikzpicture}
	\draw [->, thick] (0.2,0.2) -- (0.2,0.1) -- (1,0.1);
	\end{tikzpicture} & \hyperlink{UC6.6}{UC6.6} & \hyperlink{UC6.6}{Rispondi domanda a risposta multipla} & \hyperlink{R-1F16.1.3}{R-1F16.1.3}\\*
	\midrule
	\begin{tikzpicture}
	\draw [->, thick] (0.2,0.2) -- (0.2,0.1) -- (1,0.1);
	\end{tikzpicture} & \hyperlink{UC6.7}{UC6.7} & \hyperlink{UC6.7}{Rispondi domanda di tipo testo con parole omesse} & \hyperlink{R-1F16.1.4}{R-1F16.1.4}\\*
	\midrule
	\begin{tikzpicture}
	\draw [->, thick] (0.2,0.2) -- (0.2,0.1) -- (1,0.1);
	\end{tikzpicture} & \hyperlink{UC6.8}{UC6.8} & \hyperlink{UC6.8}{Rispondi domanda con associazione di parole} & \hyperlink{R-1F16.1.5}{R-1F16.1.5}\\*
	\midrule
	\begin{tikzpicture}
	\draw [->, thick] (0.2,0.2) -- (0.2,0.1) -- (1,0.1);
	\end{tikzpicture} & \hyperlink{UC6.9}{UC6.9} & \hyperlink{UC6.9}{Rispondi domanda a risposta aperta} & \hyperlink{R-2F16.1.6}{R-2F16.1.6}\\*
	\midrule
	\begin{tikzpicture}
	\draw [->, thick] (0.2,0.2) -- (0.2,0.1) -- (1,0.1);
	\end{tikzpicture} & \hyperlink{UC6.10}{UC6.10} & \hyperlink{UC6.10}{Feedback questionario} & \hyperlink{R-2F7.9}{R-2F7.9}\\*
	\midrule
	\begin{tikzpicture}
	\draw [->, thick] (0.2,0.2) -- (0.2,0.1) -- (1,0.1);
	\end{tikzpicture} & \hyperlink{UC6.11}{UC6.11} & \hyperlink{UC6.11}{Visualizza valutazione questionario} & \hyperlink{R-3F7.4}{R-3F7.4}\\*
	\midrule
	& \hyperlink{UC7}{UC7} & \hyperlink{UC7}{Iscrizione ad una classe} & \hyperlink{R-2F7.6}{R-2F7.6}
	
	\hyperlink{R-2F15}{R-2F15}
	
	\hyperlink{R-2F15.1}{R-2F15.1}
	
	\hyperlink{R-2F15.1.1}{R-2F15.1.1}\\*
	\midrule
	\begin{tikzpicture}
	\draw [->, thick] (0.2,0.2) -- (0.2,0.1) -- (1,0.1);
	\end{tikzpicture} & \hyperlink{UC7.1}{UC7.1} & \hyperlink{UC7.1}{Inserisci password classe} & \hyperlink{R-2F15.1}{R-2F15.1}\\*
	\midrule
	\begin{tikzpicture}
	\draw [->, thick] (0.2,0.2) -- (0.2,0.1) -- (1,0.1);
	\end{tikzpicture} & \hyperlink{UC7.2}{UC7.2} & \hyperlink{UC7.2}{Errore password classe} & \hyperlink{R-2F15.1.1}{R-2F15.1.1}\\*
	\midrule
	& \hyperlink{UC8}{UC8} & \hyperlink{UC8}{Visualizza storico studente} & \hyperlink{R-3F7.4}{R-3F7.4}
	
	\hyperlink{R-2F7.10}{R-2F7.10}
	
	\hyperlink{R-2F7.10.1}{R-2F7.10.1}
	
	\hyperlink{R-2F7.10.2}{R-2F7.10.2}
	
	\hyperlink{R-2F7.10.3}{R-2F7.10.3}
	
	\hyperlink{R-2F24}{R-2F24}
	
	\hyperlink{R-2F24.1}{R-2F24.1}
	
	\hyperlink{R-2F24.2}{R-2F24.2}
	
	\hyperlink{R-2F24.3}{R-2F24.3}
	
	\hyperlink{R-2F24.3.1}{R-2F24.3.1}
	
	\hyperlink{R-2F24.1.1}{R-2F24.1.1}
	
	\hyperlink{R-2F24.1.2}{R-2F24.1.2}
	
	\hyperlink{R-2F24.1.3}{R-2F24.1.3}
	
	\hyperlink{R-2F24.1.4}{R-2F24.1.4}
	
	\hyperlink{R-2F24.1.5}{R-2F24.1.5}
	
	\hyperlink{R-2F24.2.1}{R-2F24.2.1}
	
	\hyperlink{R-2F24.2.2}{R-2F24.2.2}
	
	\hyperlink{R-2F24.3.2}{R-2F24.3.2}
	
	\hyperlink{R-2F24.3.3}{R-2F24.3.3}
	
	\hyperlink{R-2F24.3.4}{R-2F24.3.4}\\*
	\midrule
	\begin{tikzpicture}
	\draw [->, thick] (0.2,0.2) -- (0.2,0.1) -- (1,0.1);
	\end{tikzpicture} & \hyperlink{UC8.1}{UC8.1} & \hyperlink{UC8.1}{Visualizza statistiche domande studente} & \hyperlink{R-2F24.2}{R-2F24.2}
	
	\hyperlink{R-2F24.2.1}{R-2F24.2.1}
	
	\hyperlink{R-2F24.2.2}{R-2F24.2.2}\\*
	\midrule
	\begin{tikzpicture}
	\draw [->, thick] (0.2,0.2) -- (0.2,0.1) -- (1,0.1);
	\end{tikzpicture} & \hyperlink{UC8.2}{UC8.2} & \hyperlink{UC8.2}{Visualizza statistiche questionari studente} & \hyperlink{R-2F7.10.2}{R-2F7.10.2}
	
	\hyperlink{R-2F24.3}{R-2F24.3}
	
	\hyperlink{R-2F24.3.1}{R-2F24.3.1}
	
	\hyperlink{R-2F24.3.2}{R-2F24.3.2}
	
	\hyperlink{R-2F24.3.3}{R-2F24.3.3}
	
	\hyperlink{R-2F24.3.4}{R-2F24.3.4}\\*
	\midrule
	\begin{tikzpicture}
	\draw [->, thick] (0.2,0.2) -- (0.2,0.1) -- (1,0.1);
	\end{tikzpicture} & \hyperlink{UC8.3}{UC8.3} & \hyperlink{UC8.3}{Visualizza sommario statistiche studente} & \hyperlink{R-2F7.10.3}{R-2F7.10.3}
	
	\hyperlink{R-2F24.1}{R-2F24.1}
	
	\hyperlink{R-2F24.1.1}{R-2F24.1.1}
	
	\hyperlink{R-2F24.1.2}{R-2F24.1.2}
	
	\hyperlink{R-2F24.1.3}{R-2F24.1.3}
	
	\hyperlink{R-2F24.1.4}{R-2F24.1.4}
	
	\hyperlink{R-2F24.1.5}{R-2F24.1.5}\\*
	\midrule
	& \hyperlink{UC9}{UC9} & \hyperlink{UC9}{Ricerca questionario} & \hyperlink{R-3F7.12.1}{R-3F7.12.1}
	
	\hyperlink{R-3F14}{R-3F14}
	
	\hyperlink{R-3F14.1}{R-3F14.1}
	
	\hyperlink{R-2F14.2}{R-2F14.2}
	
	\hyperlink{R-3F14.3}{R-3F14.3}
	
	\hyperlink{R-3F14.4}{R-3F14.4}
	
	\hyperlink{R-2F14.5}{R-2F14.5}
	
	\hyperlink{R-2F23.2.1}{R-2F23.2.1}
	
	\hyperlink{R-2F23.3.4.1}{R-2F23.3.4.1}
	
	\hyperlink{R-2F24.3.1}{R-2F24.3.1}\\*
	\midrule
	& \hyperlink{UC10}{UC10} & \hyperlink{UC10}{Ricerca questionario per titolo} & \hyperlink{R-3F14.1}{R-3F14.1}\\*
	\midrule
	& \hyperlink{UC11}{UC11} & \hyperlink{UC11}{Ricerca questionario per classe} & \hyperlink{R-2F14.2}{R-2F14.2}\\*
	\midrule
	& \hyperlink{UC12}{UC12} & \hyperlink{UC12}{Ricerca questionario per argomento} & \hyperlink{R-3F14.4}{R-3F14.4}\\*
	\midrule
	& \hyperlink{UC13}{UC13} & \hyperlink{UC13}{Ricerca questionario per docente} & \hyperlink{R-3F14.3}{R-3F14.3}\\*
	\midrule
	& \hyperlink{UC14}{UC14} & \hyperlink{UC14}{Ricerca questionario per difficoltà} & \hyperlink{R-2F14.5}{R-2F14.5}\\*
	\midrule
	& \hyperlink{UC26}{UC26} & \hyperlink{UC26}{Ricerca domanda} & \hyperlink{R-3F19}{R-3F19}
	
	\hyperlink{R-3F19.1}{R-3F19.1}
	
	\hyperlink{R-3F19.2}{R-3F19.2}
	
	\hyperlink{R-2F19.3}{R-2F19.3}
	
	\hyperlink{R-3F19.4}{R-3F19.4}
	
	\hyperlink{R-2F23.1.1}{R-2F23.1.1}
	
	\hyperlink{R-2F23.3.5.1}{R-2F23.3.5.1}\\*
	\midrule
	& \hyperlink{UC27}{UC27} & \hyperlink{UC27}{Ricerca domanda per keywords} & \hyperlink{R-3F19.4}{R-3F19.4}\\*
	\midrule
	& \hyperlink{UC28}{UC28} & \hyperlink{UC28}{Ricerca domanda per argomento} & \hyperlink{R-3F19.1}{R-3F19.1}\\*
	\midrule
	& \hyperlink{UC29}{UC29} & \hyperlink{UC29}{Ricerca domanda per difficoltà} & \hyperlink{R-2F19.3}{R-2F19.3}\\*
	\midrule
	& \hyperlink{UC30}{UC30} & \hyperlink{UC30}{Ricerca domanda per docente} & \hyperlink{R-3F19.2}{R-3F19.2}\\*
	\midrule
	& \hyperlink{UC15}{UC15} & \hyperlink{UC15}{Ricerca classe} & \hyperlink{R-2F15.2}{R-2F15.2}
	
	\hyperlink{R-2F20}{R-2F20}
	
	\hyperlink{R-2F20.1}{R-2F20.1}
	
	\hyperlink{R-2F20.2}{R-2F20.2}\\*
	\midrule
	& \hyperlink{UC16}{UC16} & \hyperlink{UC16}{Ricerca classe per docente} & \hyperlink{R-2F15.2}{R-2F15.2}
	
	\hyperlink{R-2F20.2}{R-2F20.2}\\*
	\midrule
	& \hyperlink{UC17}{UC17} & \hyperlink{UC17}{Ricerca classe per argomento} & \hyperlink{R-2F15.2}{R-2F15.2}
	
	\hyperlink{R-2F20.1}{R-2F20.1}\\*
	\midrule
	& \hyperlink{UC31}{UC31} & \hyperlink{UC31}{Azioni Amministratore} & \hyperlink{R-3F11}{R-3F11}
	
	\hyperlink{R-3F11.1}{R-3F11.1}
	
	\hyperlink{R-3F11.2}{R-3F11.2}
	
	\hyperlink{R-3F18}{R-3F18}
	
	\hyperlink{R-3F18.2}{R-3F18.2}
	
	\hyperlink{R-3F18.3}{R-3F18.3}\\*
	\midrule
	\begin{tikzpicture}
	\draw [->, thick] (0.2,0.2) -- (0.2,0.1) -- (1,0.1);
	\end{tikzpicture} & \hyperlink{UC31.1}{UC31.1} & \hyperlink{UC31.1}{Cambia ruolo} & \hyperlink{R-3F11.1}{R-3F11.1}
	
	\hyperlink{R-3F18.2}{R-3F18.2}\\*
	\midrule
	\begin{tikzpicture}
	\draw [->, thick] (0.2,0.2) -- (0.2,0.1) -- (1,0.1);
	\end{tikzpicture} & \hyperlink{UC31.2}{UC31.2} & \hyperlink{UC31.2}{Rimozione utente} & \hyperlink{R-3F11.2}{R-3F11.2}
	
	\hyperlink{R-3F18.3}{R-3F18.3}\\*
	\midrule
	& \hyperlink{UC32}{UC32} & \hyperlink{UC32}{Ricerca utente} & \hyperlink{R-3F32}{R-3F32}
	
	\hyperlink{R-3F32.1}{R-3F32.1}
	
	\hyperlink{R-3F32.2}{R-3F32.2}
	
	\hyperlink{R-3F32.3}{R-3F32.3}\\*
	\midrule
	& \hyperlink{UC33}{UC33} & \hyperlink{UC33}{Ricerca utente per nome completo} & \hyperlink{R-3F32.2}{R-3F32.2}\\*
	\midrule
	& \hyperlink{UC34}{UC34} & \hyperlink{UC34}{Ricerca utente per username} & \hyperlink{R-3F32.3}{R-3F32.3}\\*
	\midrule
	& \hyperlink{UC35}{UC35} & \hyperlink{UC35}{Ricerca utente per ruolo} & \hyperlink{R-3F32.1}{R-3F32.1}\\*
	\midrule
	& \hyperlink{UC3}{UC3} & \hyperlink{UC3}{Recupero password} & \hyperlink{R-2F21}{R-2F21}
	
	\hyperlink{R-2F21.1}{R-2F21.1}
	
	\hyperlink{R-2F21.2}{R-2F21.2}\\*
	\midrule
	\begin{tikzpicture}
	\draw [->, thick] (0.2,0.2) -- (0.2,0.1) -- (1,0.1);
	\end{tikzpicture} & \hyperlink{UC3.1}{UC3.1} & \hyperlink{UC3.1}{Errore mail per recupero password non presente} & \hyperlink{R-2F21.2}{R-2F21.2}\\*
	\midrule
	\caption{Tabella fonti/requisiti} \\*
\end{longtable}

\end{document}
